% concatenated from Guo_JGEA_Revised_Manuscript.tex
\documentclass[11pt]{amsart}
\usepackage[utf8x]{inputenc}
\usepackage[english]{babel}
\usepackage[T1]{fontenc}
\usepackage{microtype}
\usepackage[none]{hyphenat}

\renewcommand{\baselinestretch}{1.2} 
\usepackage{graphicx}
\usepackage{caption}
\usepackage{subcaption}


\usepackage{amssymb,amsmath}
\usepackage{mathtools}


\usepackage[unicode]{hyperref}
\usepackage[capitalise]{cleveref}

\usepackage{indentfirst}
\usepackage{enumerate,amsmath,amssymb, mathrsfs,mathtools}
\usepackage{appendix}
\usepackage{latexsym}
\usepackage{url}
\usepackage{color}
\usepackage{accents}
\usepackage{setspace}
\usepackage{pdfpages}
\usepackage{stmaryrd}

\usepackage{amsrefs}

\usepackage[margin=2.5cm]{geometry}
\allowdisplaybreaks
\newcommand{\Acirc}{\accentset{\circ}{A}}
\newcommand{\hcirc}{\accentset{\circ}{h}}
\newcommand{\hbarcirc}{\accentset{\circ}{\bar{h}}}
\newcommand{\sigmacirc}{\accentset{\circ}{\sigma}}
\newcommand{\ubar}[1]{\underaccent{\bar}{#1}}
\BibSpec{article}{%
	+{}{\PrintAuthors} {author}
	+{,}{ } {title}
	+{, }{\textit } {journal}
	+{}{ \parenthesize} {date}
		+{,  }{no. } {volume}
	+{,}{ } {pages}
	+{,}{ } {note}
%	+{, }{\PrintDOI} {url}
}
\usepackage{bbm}


\ExplSyntaxOn


\NewExpandableDocumentCommand{\gobblefirst}{m}
{
	\tl_tail:n { #1 }
}

\ExplSyntaxOff
\BibSpec{arXiv}{%
  +{}{\PrintAuthors}{author}
  +{,}{ \textit}{title}
  +{}{ \parenthesize}{date}
  +{,}{ arXiv }{eprint}
%  +{,}{ primary class }{primaryclass}
}

\BibSpec{book}{%
	+{}{\PrintAuthors}  {author}
	+{. }{}{title}
	+{,}{ }{series}
	+{,}{ vol.~}{volume}
	+{. }{\textit}{publisher}
    +{}{ \parenthesize} {date}
%	+{,}{ \DashPages} {pages}
%	+{,}{ ISBN}{isbn} 
	+{,}{ ISBN \gobblefirst}{isbn}
}
\BibSpec{collection.article}{%
	+{}{\PrintAuthors}{author}
	+{, }{}{title}
	+{, }{\textit}{booktitle}
	+{, }{ \DashPages}{pages}
	+{,}{ }{series}
	+{, }{}{volume}
	+{, }{\textit}{publisher}
	+{,}{ }{date}
}


%\BibSpec{misc}{%
%	+{}{\PrintAuthors}  {author}
%	+{. }{}{title}
%        +{,}{ }{date}
%        +{.}{ }{url}
%}

\parskip 0.0cm

\mathcode`l="8000
\begingroup
\makeatletter
\lccode`\~=`\l
\DeclareMathSymbol{\lsb@l}{\mathalpha}{letters}{`l}
\lowercase{\gdef~{\ifnum\the\mathgroup=\m@ne \ell \else \lsb@l \fi}}%
\endgroup
\def\Xint#1{\mathchoice
	{\XXint\displaystyle\textstyle{#1}}%
	{\XXint\textstyle\scriptstyle{#1}}%
	{\XXint\scriptstyle\scriptscriptstyle{#1}}%
	{\XXint\scriptscriptstyle\scriptscriptstyle{#1}}%
	\!\int}
\def\XXint#1#2#3{{\setbox0=\hbox{$#1{#2#3}{\int}$ }
		\vcenter{\hbox{$#2#3$ }}\kern-.6\wd0}}
\def\ddashint{\Xint=}
\def\dashint{\Xint-}

\newtheorem{proposition}{Proposition}
\newtheorem{notation}[proposition]{Notation}
\newtheorem{example}[proposition]{Example}
\newtheorem{definition}[proposition]{Definition}
\newtheorem{theorem}[proposition]{Theorem}
\newtheorem{lemma}[proposition]{Lemma}
\newtheorem{corollary}[proposition]{Corollary}
\newtheorem{remark}[proposition]{Remark}
\newtheorem{convention}[proposition]{Convention}




\title[Two Fefferman-type constructions]{Two Fefferman-type constructions involving almost Grassmann structures and path geometries}
\author{Zhangwen Guo}
\thanks{The author is grateful to her doctoral advisor A. \v Cap for introducing her to the problems solved in this paper and his steady support. This research was funded in part by the Austrian Science Fund (FWF): 10.55776/P33559 and 10.55776/Y963. %For open access purposes, the author has applied a CC BY public copyright license to any author-accepted manuscript version arising from this submission.
}
\address{\textnormal{Zhangwen Guo  \newline \indent
		University of Vienna \newline \indent
		Faculty of Mathematics  \newline \indent
		Oskar-Morgenstern-Platz 1 \newline \indent 1090 Vienna,	Austria %\newline\indent 		 \href{https://orcid.org/0009-0007-4027-4902}{https://orcid.org/0009-0007-4027-4902}
        \newline \indent
		  \href{mailto:zhangwen.guo@univie.ac.atm}{zhangwen.guo@univie.ac.at}}
}


\begin{document}
%\date{\today}\onehalfspacing
\begin{abstract}
We introduce a Fefferman-type construction that associates an almost Grassmann structure of type $(2,n+1)$ to every $(n+1)$-dimensional path geometry. We prove that the construction is normal and provide two equivalent characterizing conditions for all almost Grassmann structures which locally arise from this construction: one in terms of certain parallel tractors and the other in terms of a Weyl connection of an almost Grassmann structure. We prove that the latter condition is independent of the choice of Weyl connection.

We then introduce a related Fefferman-type construction associating an almost Grassmann structure of type $(2,n+1)$ to every almost Grassmann structure of type $(2,n)$. We prove that this construction is non-normal and characterize all almost Grassmann structures which locally arise in this way in Cartan geometric terms.
\end{abstract}
	\maketitle 

\section{Introduction}
Assume that $N^{n+1}$ is an $(n+1)$-dimensional manifold, $n\geq 1$. A path geometry on $N$ describes a smooth collection of paths, i.e., immersed $1$-dimensional submanifolds, on $N$ such that there is a unique such path for every point and associated direction. On a formal level, the paths are encoded on the projectivized tangent bundle $M^{2\,n+1}=\mathcal PTN$ as the unique line distribution $E\subseteq TM$ such that (i) the rank $(n+1)$-tautological distribution of $M$ can be written as the direct sum of $E$ and the vertical bundle $V\subseteq TM$ of $M\to N$ and (ii) the integral manifolds of $E$ descend to the paths on $N$. In particular, there is a construction $(N,\{\text{paths}\})\mapsto(M,E,V)$
leading to a $(2\,n+1)$-dimensional generalized path geometry; see \cref{2.4}. Conversely,
it has been proven that if the defining distribution $V$ of a $(2\,n+1)$-dimensional generalized path geometry is involutive, then this generalized path geometry locally arises from the construction described above. For this reason, we will refer to a generalized path geometry with involutive $V$ as a path geometry. Moreover, the defining rank-$n$-distribution has been shown to always be involutive if $n\neq2$; see, e.g., \cite{Book}*{Proposition 4.4.4}. In particular, a generalized path geometry $(M,E,V)$ with $\dim(M)\neq 5$ is always a path geometry.


Assume now that $M^{p\,q}$ is a $p\,q$-dimensional manifold, $p,q\geq2$. We follow the convention that an almost Grassmann structure $(M,E,F)$ of type $(p,q)$ on $M$ is defined by two auxiliary vector bundles $E\to M$ and $F\to M$ of rank $p$ and of rank $q$, respectively, and isomorphisms
\begin{align}\label{AG isomorphisms}
    E^*\otimes F\cong TM\quad\text{and}\quad\wedge^pE^*\cong\wedge^q\,F.
\end{align}

Recall from, e.g., \cite{Book}*{Subsections 4.1.3, 4.2.3 and 4.4.3} the following background. Each of the above geometric structures carries a canonical parabolic geometry, which in turn comes with its invariants. The types of fundamental invariants splits into different cases. For the general case of a $(2\,n+1)$-dimensional generalized path geometry with either $n=2$ and $V$ involutive or $n\geq 3$ (where $V$ is automatically involutive), the fundamental invariants consist of a harmonic torsion and a harmonic curvature; see \eqref{harmonic components path} below. In the case where $n=2$ and $V$ not involutive, there is an extra invariant torsion in addition to the general case. For $n=1$, the fundamental invariants consists of two Cotten-York-type tensors; see \cref{LC} below. On the other hand, for the general case of an almost Grassmann structure of type $(2,n)$, $n\geq 3$, the fundamental invariants consist of a harmonic torsion and a harmonic curvature; see \eqref{harmonic components AG} below. The fundamental invariants for an almost Grassmann structure of type $(2,2)$ is a pair of harmonic curvatures; see \cref{LC}. Finally, the fundamental invariants for an almost Grassmann structure of type $(p,q)$, $p,q\geq 3$ is a pair of harmonic torsions. 

In this present article, we only deal with the general situation of the two geometric structures described above, with the exception of \cref{LC} where $3$-dimensional generalized path geometries and almost Grassmann structures of type $(2,2)$ are discussed. In particular, the harmonic torsion $\tau\in\Omega^2(M,TM)$ of an almost Grassmann structure $(M,E,F)$ of type $(2,n)$, $n>2$, coincides with the torsion of any Weyl connection on $(M,E,F)$. Indeed, the Weyl connections on $(M,E,F)$ are exactly pairs of connections on $(E,F)$ which are compatible with the second isomorphism in \eqref{AG isomorphisms} and, via the first isomorphism in \eqref{AG isomorphisms}, induce a connection on $TM$ whose torsion is $\partial^*$-closed and hence equal to $\tau$; see, e.g., \cite{Book}*{Theorem 5.2.3\,(1)}. Since $\tau$ is $\partial^*$-closed, it vanishes upon any contraction. Hence, up to a sign, there is a unique nontrivial contraction of $\tau\otimes\tau$ resulting in a $(0,2)$-tensor field. This contraction is clearly symmetric. We will denote it by
\begin{align}\label{torsion contraction def}
    tr(\iota_\tau\tau)\in\Gamma(\mathrm{Sym}^2\,T^*M).
\end{align}

We introduce a Fefferman-type functor that associates to every path geometry $(M^{2\,n+1},E,V)$ an almost Grassmann structure of type $(2,n+1)$. This is motivated by the work of Crampin and Saunders \cite{CrSa}, who associate to every path geometry $(M^{2\,n+1},E,V)$ a Segre structure of type $(2,n+1)$. Recall that, by definition, every almost Grassmann structure induces a Segre structure of the same type; see, e.g., \cite{Met13}.

The framework of Fefferman-type functors is developed from the pioneering work of Fefferman \cite{Fef76}, who associates to every strictly pseudoconvex real hypersurface in $\mathbb C^n$ a conformal geometry on the total space of a circle bundle. It has been generalized to a construction associating to every abstract CR geometry a conformal geometry by Burns, Diederich and Shnider in \cite{BDS}. As a generalization of the latter, the Fefferman-type functors on Cartan geometries have been introduced by \v Cap in \cite{Cap06}. The framework of Fefferman-type functors has been extended to a slightly broader class of functors with the same properties as the original one by \v Cap and \v Z\'adn\'ik in \cite{CZ09}. This is based on the more general framework of so-called extension functors on Cartan geometries.

The Fefferman-type functor we specify is applied to the canonical parabolic geometry associated to every path geometry $(M^{2\,n+1},E,V)$ together with an auxiliary orientation on $E\oplus V$ or, equivalently, on $E^*\otimes\wedge^n\,V$, since $\otimes^2\, E$ is canonically oriented. The result is a parabolic geometry that induces an almost Grassmann structure of type $(2,n+1)$ on the principal $\mathbb R_+$-bundle $GL_+(E^*\otimes\wedge^n\,V)$. If the initial path geometry arises from a collection of paths on $N^{n+1}$ as described in the first paragraph in this section, the base space of the resulting almost Grassmann structure is the weighted slit tangent bundle $TN\otimes\mathcal E(-1)\setminus\{0\}$. It is precisely the base space of the Segre structure that Crampin and Saunders associate to a path geometry on $N$; see \cite{CrSa}*{Section 3}.

We prove that the resulting parabolic geometry coincides with the parabolic geometry canonically associated to the resulting almost Grassmann structure. In other words, the Fefferman-type construction we specify is normal. In particular, the functor induces an explicit relation between the canonical Cartan connection of a path geometry and the canonical Cartan connection of the resulting almost Grassmann structure. From this relation, it is straightforward to see that the harmonic torsion $\tau$ of every resulting almost Grassmann structure has the property that $tr(\iota_\tau\tau)=0$. Moreover, the resulting almost Grassmann structure has vanishing harmonic torsion if and only if the initial path geometry has vanishing harmonic torsion.

The relations we find also enable us to characterize all almost Grassmann structures that are locally obtained from the construction. The characterization of conformal geometries that locally arise from the construction of Burns, Diederich and Shnider in \cite{BDS} has been given by Sparling; see \cite{Graham}*{Theorem 3.1}. Sparling's condition characterizing such structures is the existence of a certain isotropic conformal Killing field. In the context of parabolic geometry, such a characterization has been obtained by \v Cap and Gover in terms of the existence of a parallel tractor; see \cite{CG08}*{Theorem 3.1} together with \cite{CG06}*{Theorem 2.5}. We also refer the reader to the following articles for some analogous characterizations of geometric structures arising from the respective Fefferman-type constructions.

\begin{table}[h!]
  \begin{center}
    \label{tab:table1}
    \begin{tabular}{l|c}
      \textbf{Article} & \textbf{Fefferman-type constructions} \\
\hline
\cite{Arm} & $(n,n\,(n+1)/2)$-distribution-to-almost spinorial \\
      &$(3,6)$-distribution-to-conformal\\
       &integrable split signature CR-to-$(3,6)$-distribution\\
        &integrable Lagrangian contact-to-$(3,6)$-distribution\\
\hline\cite{HS09} & $(2,3,5)$-distribution-to-conformal \\
\hline\cite{Alt10} & quaternionic contact-to-conformal \\
\hline
\cite{HS11} & $(2,3,5)$-distribution-to-conformal\\
      &$(3,6)$-distribution-to-conformal\\
\hline\cite{LNS17}&$(2,3,5)$-distribution-to-Lie contact\\
\hline\cite{Ham+17}&projective-to-conformal\\
    &integrable Lagrangian contact-to-conformal\\
\end{tabular}
  \end{center}
\end{table}

We also extend our Fefferman-type functor to define a construction from almost Grassmann structures of type $(2,n)$ to almost Grassmann structures of type $(2,n+1)$. We prove that the resulting almost Grassmann structure of type $(2,n+1)$ has the property $tr(\iota_{\tilde\tau}\tilde\tau)=0$ for the harmonic torsion $\tilde\tau$ if and only if the initial almost Grassmann structure of type $(2,n)$ has the same property on its harmonic torsion. Moreover, the resulting almost Grassmann structure of type $(2,n+1)$ is torsion-free if and only if the initial almost Grassmann structure of type $(2,n)$ is torsion-free. We note that, among all the above mentioned characterizations of Fefferman-type constructions, that of projective-to-conformal \cite{Ham+17} is the only one that deals with so-called non-normal Fefferman-type construction. Our Fefferman-type construction from almost Grassmann structures of type $(2,n)$ to almost Grassmann structures of type $(2,n+1)$ is also of this kind. We follow a procedure similar to that in \cite{Ham+17} for characterizing the construction.

The general theory of Fefferman-type constructions prior to the present work has only been well-established on the Cartan level: Given two Cartan geometries related by a Fefferman-type construction, the relation of their Cartan connections and the relation of their Cartan curvatures have been known; moreover, there is an intrinsic characterization in terms of tractor connection and Cartan curvature. These results are reviewed in \cref{section abstract Fefferman} and have originated in \cite{Cap06}, \cite{CGH} and \cite{Cap05}. The difficulty of passing these results to a Fefferman-type construction between geometric structures is due to the lack of general theory when the relevant Cartan geometries are parabolic and need to be considered to vary within certain equivalent classes. In particular, essential questions such as normality and normalization have to be investigated on an individual basis. Our technical solution to the above difficulty is to provide case-by-case descriptions on the Kostant codifferentials applied to relevant related tensors; see the proofs of \cref{path Fef is normal}, \cref{charpath}, \cref{AG Fef normality} and \cref{charAG}. Moreover, the proofs of our two characterization results \cref{charpath} and \cref{charAG} involve normalization in the reduction procedure. In the situation of \cref{charpath} the normalization procedure is straightforward, whereas \cref{charAG} involves a normalization homogeneity-by-homogeneity, using a general result \cref{modified correspondence} originated in \cite{Ham+17}. In particular, the current work together with previously mentioned works on Fefferman-type constructions provide ingredients on extending the formal theory to the situation of parabolic geometries.








\section{A generalized correspondence space theorem}
In this section, we assume that $G$ is a Lie group and $H\subseteq P\subseteq G$ are closed Lie subgroups. Denote their Lie algebras by $\mathfrak h\subseteq\mathfrak p\subseteq\mathfrak g$, respectively.

\begin{definition}\label[definition]{Cartan geometry definition}(See, e.g., \cite{Book}*{Subsection 1.5.1}) A Cartan geometry $(\mathcal G,\omega)$ of type $(G,P)$ is a principal $P$-bundle $\mathcal G\to M$ together with a Cartan connection $\omega\in\Omega^1(\mathcal G,\mathfrak g)$, i.e., a $\mathfrak g$-valued one-form on $\mathcal G$ such that (i) $\omega$ is $P$-equivariant with respect to the principal right action and the adjoint action of $G$ on $\mathfrak g$ restricted to $P\subseteq G$, (ii) $\omega(\zeta_X)=X$ where $\zeta_X\in\mathfrak X(\mathcal G)$ is the fundamental vector field generated by $X\in\mathfrak p$ and (iii) $T\mathcal G\xrightarrow[\cong]{(\pi,\omega)}\mathcal G\times\mathfrak g$ where $\pi:T\mathcal G\twoheadrightarrow\mathcal G$ is the natural projection.

The curvature of the Cartan connection is the two-form $\kappa\in\Omega^2(\mathcal G,\mathfrak g)$ given by 
\begin{align*}
\kappa(\xi,\eta)&=d\omega(\xi,\eta)+[\omega(\xi),\omega(\eta)]\\
&=\xi.\omega(\eta)-\eta.\omega(\xi)-\omega([\xi,\eta])+[\omega(\xi),\omega(\eta)]    
\end{align*}
for all $\xi,\eta\in\mathfrak X(\mathcal G)$.
\end{definition}
It is easy to see that the Cartan connection $\omega:T\mathcal G\to\mathfrak g$ induces an isomorphism
\[TM\cong\mathcal G\times_P\mathfrak g/\mathfrak p.\]
Recall from, e.g., \cite{Book}*{Lemma 1.5.1} that the curvature $\kappa$ is $P$-equivariant and horizontal. In particular, $\kappa$ may also be expressed as a $P$-equivariant map $\kappa:\mathcal G\to L(\wedge^2\,\mathfrak g/\mathfrak p,\mathfrak g)$ and as a two-form $\kappa\in\Omega^2(M,\mathcal G\times_P\mathfrak g)$.

The orbit space $\tilde M=\mathcal G/H$ is called the correspondence space of the Cartan geometry $(\mathcal G\to M,\omega)$. It canonically carries the Cartan geometry 
\begin{align*}
    (\mathcal G\to\tilde M,\omega)
\end{align*}
of type $(G,H)$, whose curvature is also $\kappa\in\Omega^2(\mathcal G,\mathfrak g)$. In particular, as a section of $\Omega^2(\tilde M,\mathcal G\times_H\mathfrak g)$, $\kappa$ vanishes upon inserting any section of the distribution
\[\mathcal G\times_H\mathfrak p/\mathfrak h\subseteq \mathcal G\times_H\mathfrak g/\mathfrak h \cong T\tilde M;\]
see, e.g., \cite{Book}*{Proposition 1.5.13}. Conversely, the correspondence space theorem \cite{Book}*{Theorem 1.5.14} states that if $(\mathcal G\to\tilde M,\omega)$ is a Cartan geometry of type $(G,H)$ and its curvature $\kappa$ satisfies
\[\kappa(\mathcal G\times_H\mathfrak p/\mathfrak h,T\tilde M)=0,\]
then $(\mathcal G\to\tilde M,\omega)$ is locally isomorphic to the result of the correspondence space construction; see the first part of \cref{locally correspondence} below.

\begin{definition}\label[definition]{locally correspondence}
Let $(\mathcal G\to\tilde M,\omega)$ be a Cartan geometry of type $(G,H)$ and $\mathfrak n\subseteq \mathfrak g$ a $P$-submodule. Denote by $p:\mathcal G\to\tilde M$ the natural projection. We say that $(\mathcal G\to\tilde M,\omega)$ is locally isomorphic to the result of the correspondence space construction if there exists an open cover $\tilde M=\bigcup_{i\in \mathcal I} \tilde U_i$ and a collection of Cartan geometries $\{(\mathcal G_i\to M_i,\omega_i):i\in\mathcal I\}$
of type $(G,P)$ such that the following holds. Denote by $\{q_i:\mathcal G_i\to \mathcal G_i/H\}_{i\in\mathcal I}$ the natural projections. For every $i\in\mathcal I$, there exists an open subset $\tilde V_i\subseteq\mathcal G_i/H$ and an isomorphism
\[\Phi_i:(p^{-1}(\tilde U_i)\to\tilde U_i,\omega|_{p^{-1}(\tilde U_i)})
\xrightarrow{\cong}
(q_i^{-1}(\tilde V_i)\to\tilde V_i,\omega_i|_{q_i^{-1}(\tilde V_i)})\]
of Cartan geometries of type $(G,H)$.

More generally, we say that $(\mathcal G\to\tilde M,\omega)$ is locally isomorphic to the result of the correspondence space construction modulo $L(\mathfrak g/\mathfrak p,\mathfrak n)$ if there exists an open cover $\tilde M=\bigcup_{i\in \mathcal I} \tilde U_i$ and a collection of Cartan geometries 
$\{(\mathcal G_i\to M_i,\omega_i):i\in\mathcal I\}$ of type $(G,P)$ such that the following holds. Denote by $\{q_i:\mathcal G_i\to \mathcal G_i/H\}_{i\in\mathcal I}$ the natural projections. For every $i\in\mathcal I$, there exists an open subset $\tilde V_i\subseteq\mathcal G_i/H$ and an isomorphism
\[\Phi_i:(p^{-1}(\tilde U_i)\to\tilde U_i)
\xrightarrow{\cong}
(q_i^{-1}(\tilde V_i)\to\tilde V_i)\]
of principal $H$-bundles such that $\Phi_i^*\omega_i\equiv\omega$
modulo $L(\mathfrak g/\mathfrak p,\mathfrak n)$, i.e.,
\begin{equation*}
   \begin{aligned}
   \begin{cases}
\Phi_i^*\omega_i-\omega\in\Omega^1(p^{-1}(\tilde U_i),\mathfrak n)\quad\text{and}\\
(\omega|_{p^{-1}(\tilde U_i)})^{-1}(X)=(\Phi_i^*\omega_i)^{-1}(X)\quad\text{for all }X\in\mathfrak p.
\end{cases}
\end{aligned}
\end{equation*}
\end{definition}


In \cite{Ham+17}*{Proposition 4.13}, the authors have developed a generalization of the classical correspondence space theorem \cite{Book}*{Theorem 1.5.14} recalled just before \cref{locally correspondence} in their specific setting. The proof of \cref{modified correspondence} below expands the general line of their argument. In the special case where $\mathfrak n=0$, it is just the classical correspondence space theorem. 
\begin{proposition}\label{modified correspondence}
Let $H\subseteq P\subseteq G$ be as at the beginning of this section. Let $\mathfrak n\subseteq\mathfrak g$ be a $P$-submodule. Given a Cartan geometry $(\mathcal G\to\tilde M,\omega)$ of type $(G,H)$ with curvature $\kappa$, consider the 
distribution
\[\mathcal V\tilde M=\mathcal G\times_H\mathfrak p/\mathfrak h\subseteq T\tilde M.\]
The following are equivalent.
\begin{enumerate}
\item [(i)] $\kappa(\mathcal V\tilde M,\mathcal V\tilde M)=0$ and $\kappa(\mathcal V\tilde M,T\tilde M)\subseteq\mathcal G\times_H\mathfrak n$.
\item [(ii)] $(\mathcal G\to\tilde M,\omega)$ is locally the result of the correspondence space construction modulo $L(\mathfrak g/\mathfrak p,\mathfrak n)$. 
\end{enumerate}
\end{proposition}
\begin{proof}
Suppose that $(\mathcal G\to M,\omega')$ is a Cartan geometry of type $(G,P)$, $\tilde M=\mathcal G/H$ and $\omega'\equiv\omega$ modulo $L(\mathfrak g/\mathfrak p,\mathfrak n)$; see \cref{locally correspondence}. Then for the curvature $\kappa'$ of $\omega'$, there holds $\kappa'(\mathcal V\tilde M,T\tilde M)=0$. To prove that (ii) implies (i), it suffices to show that
\begin{align*}
    (\kappa'-\kappa)(\omega^{-1}(X),\omega^{-1}(Y))&=0\quad\text{and}\\
    (\kappa'-\kappa)(\omega^{-1}(X),\omega^{-1}(A))&\subseteq\mathfrak n
\end{align*}
for all $X,Y\in\mathfrak p$ and $A\in\mathfrak g$. Indeed, by assumption, $\omega$ agrees with $\omega'$ on the vertical bundle of $\mathcal G\to M$. Since $\mathfrak n\subseteq\mathfrak g$ is a $P$-submodule, there holds $[\mathfrak p,\mathfrak n]\subseteq\mathfrak n$. The above two relations follow from the formula for curvature in \cref{Cartan geometry definition} and a brief computation.

Conversely, we prove that (i) implies (ii) in two steps. The first step is precisely the first half of the proof of the classical correspondence space theorem; see, e.g., \cite{Book}*{Theorem 1.5.14}. Indeed, the assumption
\[\kappa(\mathcal V\tilde M,\mathcal V\tilde M)=0\]
means that $\omega^{-1}:\mathfrak p\to\mathfrak X(\mathcal G)$ is a Lie algebra homomorphism. Moreover, it implies that $\mathcal V\tilde M\subseteq T\tilde M$ is an involutive distribution and thus around every point of $\tilde M$ we may choose a sufficiently small local leaf space 
\[\tilde M\supseteq\tilde W\twoheadrightarrow U\]
such that there exists a global section of $\tilde W\twoheadrightarrow U$ and a local section of $p:\mathcal G\to\tilde M$ defined on $\tilde W$. Denote their composition by 
\[\sigma:U\to\mathcal G.\]
We choose an arbitrary inner product on $\mathfrak p$, thus 
\[\mathfrak p=\mathfrak h\oplus\mathfrak h^\perp.\]
Let $B_{\mathfrak h^\perp}\subseteq\mathfrak h^\perp$ and $B_{\mathfrak h}\subseteq\mathfrak h$ be sufficiently small open balls such that the maps
\begin{align*}
        B_{\mathfrak h^\perp}\times B_{\mathfrak h}\to P,\quad
        (X,Y)\mapsto\exp\,X\,\exp\,Y
\end{align*}
and
\begin{align*}
    B_{\mathfrak h^\perp}\to P/H,\quad X\mapsto\exp\,X \bmod H
\end{align*}
are both diffeomorphisms onto their images. In particular, the group operation $P\times P\to P$ restricts to an injective map
\begin{align}\label{injective group operation}
\exp\,B_{\mathfrak h^\perp}\times H\to P.    
\end{align}

It follows from Lie's second fundamental theorem that, shrinking $U$ further, the Lie algebra homomorphism $\omega^{-1}:\mathfrak p\to\mathfrak X(\mathcal G)$ integrates to a local $P$-action
\[\mathcal G\times P\supseteq\sigma(U)\times(\exp\,B_{\mathfrak h^\perp}\times\exp\,B_{\mathfrak h})\xrightarrow{F} p^{-1}(\tilde W)\subseteq\mathcal G.\]
Observe that the restriction to $(\exp\,B_{\mathfrak h^\perp}\times\exp\,B_{\mathfrak h})$ of the left invariant vector field generated by $X\in\mathfrak p$ is related to $\omega^{-1}(X)$ by $F$ and that the tangent map of $F$ is bijective at every point. By shrinking $B_{\mathfrak h^\perp}$ and $B_{\mathfrak h}$, we may assume that $F$ is a diffeomorphism onto its image. Now we restrict $F$ to $\sigma(U)\times\exp\,B_{\mathfrak h^\perp}\times\{e\}$ and extend it to a local $P$-action
\[\tilde F:\sigma(U)\times(\exp\,B_{\mathfrak h^\perp})H\to
p^{-1}(\tilde W)\subseteq\mathcal G\]
that is a diffeomorphism onto its image. This is well-defined in view of the injectivity of \eqref{injective group operation}. Now let $\mathcal G'=U\times P$. Denote by $q:\mathcal G'\to \mathcal G'/H$ the natural projection and by $\mathrm{im}\,\tilde F$ the image of $\tilde F$. Let
\begin{align*}
\tilde V=U\times\frac{(\exp\,B_{\mathfrak h^\perp})H}{H}\subseteq \mathcal G'/H
\quad\text{and}\quad
\tilde U=p(\mathrm{im}\,\tilde F)\subseteq\tilde W\subseteq\tilde M.
\end{align*}
Note that
\begin{align*}
q^{-1}(\tilde V)=U\times (\exp\,B_{\mathfrak h^\perp})H\subseteq\mathcal G'
\quad\text{and}\quad
p^{-1}(\tilde U)=\mathrm{im}\,\tilde F.
\end{align*}
Identifying $\sigma(U)$ with $U$, we obtain an isomorphism
\[\Phi=\tilde F:(q^{-1}(\tilde V)\to\tilde V)\to (p^{-1}(\tilde U)\to\tilde U)\]
of principal $H$-bundles. By construction, $\Phi^*\omega$ is a Cartan connection of type $(G,H)$ that in addition reproduces the generators in $\mathfrak p$ of the fundamental vector fields restricted to the subspace $q^{-1}(\tilde V)$. Hence there exists a unique Cartan connection $\omega'\in\Omega^1(\mathcal G',\mathfrak g)^P$ of type $(G,P)$ such that
\begin{align*}  
\omega'|_{U\times H}=\Phi^*\omega|_{U\times H}.\end{align*} In addition, there holds
\begin{align}\label{P-equiv}
    (\omega'|_{q^{-1}(\tilde V)})^{-1}(X)=(\Phi^*\omega)^{-1}(X)
\end{align}
for all $X\in\mathfrak p$.

The second step is a minor modification of the second half of the proof of the classical correspondence space theorem. Assume in addition that
\[\kappa(\mathcal V\tilde M,T\tilde M)\subseteq\mathcal G\times_H\mathfrak n.\]
It remains to show that $\Phi^*\omega-\omega'\in\Omega^1(q^{-1}(\tilde V),\mathfrak n)$. Recall from the first step that $\Phi^*\omega$ coincides with $\omega'$ on $U\times \{e\}$. Moreover, recall that by assumption, every element in $(\exp\,B_{\mathfrak h^\perp})H$ may be expressed as $(\exp\,X)\,h$ for a unique choice of $X\in B_{\mathfrak h^\perp}$ and $h\in H$. Hence it suffices to verify that
\begin{align*}
    (r^{\exp{(X)}})^*(\Phi^*\omega)-Ad(\exp{(-X)})\circ \Phi^*\omega\in\Omega^1(q^{-1}(\tilde V),\mathfrak n)
\end{align*}
for all $X\in B_{\mathfrak h^\perp}$. Indeed, denote by $\mathcal L$ the Lie derivative and by the lower dot the directional derivative. Then for all $\xi,\eta\in\mathfrak X(\mathcal G)$, there holds
\[(\mathcal L_\xi\omega)(\eta)=\xi.\omega(\eta)-\omega([\xi,\eta])\]
and so
\[\kappa(\xi,\eta)=(\mathcal L_\xi\omega)(\eta)-\eta.\omega(\xi)+[\omega(\xi),\omega(\eta)].\]
Substituting $\xi=\omega^{-1}(A)$ for any $A\in\mathfrak g$, we obtain
\begin{align*}
    \kappa(\omega^{-1}(A),\,\cdot\,)=\mathcal L_{\omega^{-1}(A)}\omega+ad(A)\circ\omega.
\end{align*}
This relation together with the assumption $\kappa(\mathcal V\tilde M,T\tilde M)\subseteq\mathcal G\times_H\mathfrak n$ imply that
\[\mathcal L_{\omega^{-1}(X)}\omega+ad(X)\circ\omega\in\Omega^1(\mathcal G,\mathfrak n)\]
for all $X\in B_{\mathfrak h^\perp}$. Since $\Phi^*(\omega^{-1}(X))$ is the restriction to $q^{-1}(\tilde V)$ of the fundamental vector field $\zeta_X\in\mathfrak X(\mathcal G')$ generated by $X$, the pullback of the above expression integrates to
\[(r^{\exp{(tX)}})^*(\Phi^*\omega)-Ad(\exp{(-tX)})\circ \Phi^*\omega\in\Omega^1(q^{-1}(\tilde V),\mathfrak n)\]
for all $t\in [0,1]$. Putting $t=1$, (ii) follows as required.
\end{proof}
Compared to the classical correspondence theorem, this variation requires less information on the curvature of the Cartan geometry of interest and is particularly useful in settings where we are flexible to modify the Cartan connections, as is the case when studying Fefferman-type constructions.

\section{A characterization theorem for abstract Fefferman-type constructions}\label{section abstract Fefferman}
We follow the original setup of the Fefferman-type functors for Cartan geometries in \cite{Cap06}*{Subsection 4.5} and establish a characterization result for Cartan geometries arising from this construction. We refer the reader also to \cite{CZ09}*{Subsection 3.2}, which generalizes the original construction to a slightly broader class of functors and establishes a relation between the curvature of the initial Cartan geometry and the curvature of the resulting Cartan geometries; see \cite{CZ09}*{Subsection 3.3}.

In this section, we assume that $G$ and $\tilde G$ are Lie groups, $\tilde P\subseteq\tilde G$ is a closed Lie subgroup and \[i:G\to\tilde G\] is an injective Lie group homomorphism such that the $G$-orbit of $e\tilde P$ in $\tilde G/\tilde P$ is open. Then for any closed Lie subgroup $P\subseteq G$ containing $H=i^{-1}(\tilde P)$, $i$ defines a Fefferman-type functor from Cartan geometries of type $(G,P)$ to Cartan geometries of type $(\tilde G,\tilde P)$. We describe the resulting Cartan geometry of the functor applied to a Cartan geometry $(\mathcal G\to M,\omega)$ of type $(G,P)$. Let $\tilde M=\mathcal G/H$ and
\[j:\mathcal G\to\mathcal G\times_H\tilde P=\tilde{\mathcal G}\]
be the extension of the principal $H$-bundle $\mathcal G\to\tilde M$ based on the group homomorphism $i|_H:H\to\tilde P$. We obtain $(\tilde{\mathcal G}\to\tilde M,\tilde\omega)$ where $\tilde\omega$ is the unique Cartan connection of type $(\tilde G,\tilde P)$ on $\tilde{\mathcal G}\to\tilde M$ such that \[j^*\tilde\omega=i'\circ\omega.\]

Denote the Lie algebras of the Lie groups 
\begin{align*}
H\subseteq P\subseteq G\quad\text{and}\quad\tilde P\subseteq\tilde G    
\end{align*} by $\mathfrak h\subseteq\mathfrak p\subseteq\mathfrak g$ and $\tilde{\mathfrak p}\subseteq\tilde{\mathfrak g}$, respectively. Recall that the $G$-orbit of $e\tilde P$ in $\tilde G/\tilde P$ is open if and only if $G$ acts locally transitively on $\tilde G/\tilde P$, i.e., if $i'(\mathfrak g)+\tilde{\mathfrak p}=\tilde{\mathfrak g}$. Observe that the composition map $\mathfrak g\xrightarrow{i'} \tilde{\mathfrak g}\twoheadrightarrow\tilde{\mathfrak g}/\tilde{\mathfrak p}$ is surjective and descends to an isomorphism $\mathfrak g/\mathfrak h\cong\tilde{\mathfrak g}/\tilde{\mathfrak p}$ of $H$-modules. We denote by
\begin{align*}
 \pi:\tilde{\mathfrak g}/\tilde{\mathfrak p}\twoheadrightarrow\mathfrak g/\mathfrak p   
\end{align*}
the composition of the natural projection $\mathfrak g/\mathfrak h\twoheadrightarrow\mathfrak g/\mathfrak p$ with the inverse of the above isomorphism. Then the curvature $\tilde\kappa$ of $\tilde\omega$ is completely determined by the curvature $\kappa$ of $\omega$ via the relation
\begin{align}\label{tilde kappa}
    \tilde\kappa(j(u))=i'\circ\kappa(u)\circ\wedge^2\, \pi
\end{align}
for all $u\in\mathcal G$.

Conversely, the subclass of all Cartan geometries of type $(\tilde G,\tilde P)$ that is locally the result of the Fefferman-type functor may be characterized in terms of the existence of a certain parallel tractor. Here, being locally the result of a Fefferman-type functor is defined in a way analogous to \cref{locally correspondence}. Assume now that $(\tilde{\mathcal G}\to \tilde M,\tilde\omega)$ is any Cartan geometry of type $(\tilde G,\tilde P)$. Consider the canonical extension $\tilde{\mathcal G}\to \tilde{\mathcal G}\times_{\tilde P}\tilde G$ of $\tilde{\mathcal G}$ to a principal $\tilde G$-bundle. There is a unique principal connection $\tilde\theta\in\Omega^1(\tilde{\mathcal G}\times_{\tilde P}\tilde G,\tilde{\mathfrak g})$ which pulls back to $\tilde\omega$; see, e.g., \cite{Book}*{Theorem 1.5.6}. In particular, given any $\tilde G$-representation $\tilde{\mathbb W}$, we obtain a tractor bundle \[\tilde{\mathcal W}=(\tilde{\mathcal G}\times_{\tilde P}\tilde G)\times_{\tilde G}\tilde{\mathbb W}\cong\tilde{\mathcal G}\times_{\tilde P}\tilde{\mathbb W}\] with tractor connection induced by the principal connection $\tilde\theta$.

Assume in addition that there exists $\tilde w_0\in\tilde{\mathbb W}$ such that 
\[i(G)=\mathrm{Stab}_{\tilde G}(\tilde w_0)
=\{\tilde g\in\tilde G:\tilde g\,\tilde w_0=\tilde w_0\}.\]
The $\tilde P$-orbit of $\tilde w_0$ defines a fiber subbundle
\[\tilde{\mathcal O}=\tilde{\mathcal G}\times_{\tilde P}(\tilde P\,\tilde w_0)\subseteq\tilde{\mathcal W}\]
with standard fiber $\tilde P\,\tilde w_0\cong\tilde P/i(H)$. The following characterization may be considered as a kin of \cite{CGH}*{Theorem 2.6 (ii)} in a simplified setup.
\begin{theorem}\label[theorem]{Hol}
Notation as above. $(\tilde{\mathcal G}\to\tilde M,\tilde\omega)$ is locally the result of the Fefferman-type functor applied to a Cartan geometry of type $(G,P)$ if and only if there exists a parallel section
\[\sigma\in C^\infty(\tilde{\mathcal G},\tilde P\,\tilde w_0)^{\tilde P}=\Gamma(\tilde{\mathcal O})\subseteq\Gamma(\tilde{\mathcal W})\]
of the tractor bundle $\tilde{\mathcal W}$ such that whenever $\sigma(u)=\tilde w_0$ where $u\in\tilde{\mathcal G}$, there holds
\begin{align*}
\tilde\kappa(u)(i'(\mathfrak p),\tilde{\mathfrak g})=0.   
\end{align*}
\end{theorem}

\begin{proof}
Let us identify $H\subseteq G$ with their images under $i$ and identify their Lie algebras with their images under $i'$.

Assume that $(\tilde{\mathcal G}\to\tilde M,\tilde\omega)$ is the result of the Fefferman-type functor applied to a Cartan geometry $(\mathcal G\to M,\omega)$ of type $(G,P)$. Since $H$ fixes $\tilde w_0$, the constant map $\sigma:\mathcal G\to\{\tilde w_0\}$ is $H$-equivariant. It is extended to 
\[\tilde\sigma\in C^\infty(\tilde{\mathcal G},\tilde P\,\tilde w_0)^{\tilde P}=\Gamma(\tilde{\mathcal O})\]
in a unique way. Denote by $p:\mathcal G\to \tilde M$ the natural projection. For every tangent vector in $T\mathcal G$, there holds
\begin{align}\label{tractor connection on the distinguished tractor}
    \nabla^{\tilde{\mathcal W}}_{Tp(\,\cdot\,)}\tilde\sigma=d\,\sigma(\,\cdot\,)+\omega(\,\cdot\,)\,\sigma.
\end{align}
Here, $d\,\sigma$ denotes the directional derivative of the constant map $\sigma:\mathcal G\to\{\tilde w_0\}$ and $\omega(\,\cdot\,)\,\sigma$ denotes the induced Lie algebra action of $\mathfrak g$ on $\tilde w_0$. Both of them are zero here. In particular, $\tilde\sigma$ is a parallel section of $\tilde{\mathcal W}$. Moreover, it follows from the definition of $\tilde{\mathcal W}$ that for $u\in\tilde{\mathcal G}$, there holds $\tilde\sigma(u)=\tilde w_0$ if and only if $u\in\mathcal G$. In this case, it follows from \eqref{tilde kappa} that $\tilde\kappa(u)(i'(\mathfrak p),\tilde{\mathfrak g})=0$.

Conversely, let $\tilde\sigma\in C^\infty(\tilde{\mathcal G},\tilde P\,\tilde w_0)^{\tilde P}=\Gamma(\tilde{\mathcal O})$ be as stated in the theorem and
\[\mathcal G=\{u\in\tilde{\mathcal G}:\tilde\sigma(u)=\tilde w_0\}.\]
Observe that the natural projection $\mathcal G\to\tilde M$ admits smooth local sections and that the preimage of every point in $\tilde M$ is exactly an $H$-orbit of the fiber of $\tilde{\mathcal G}$ over the same point. Thus $\mathcal G$ is a reduction of $\tilde{\mathcal G}$ to a principal $H$-bundle. It remains to show that $\omega=\tilde\omega|_{\mathcal G}$ has values in $\mathfrak g$ so that, in particular, $(\mathcal G\to\tilde M,\omega)$ is a Cartan geometry of type $(G,H)$. Indeed, $\sigma=\tilde\sigma|_{\mathcal G}$ is the $\tilde w_0$-valued constant map. Since $\mathcal G\subseteq\tilde{\mathcal G}$, \eqref{tractor connection on the distinguished tractor} holds where $\nabla^{\tilde{\mathcal W}}\tilde\sigma$ and $d\,\sigma$ are both zero. In particular, the values of $\omega$ annihilate $\tilde w_0$. Since $G=\mathrm{Stab}_{\tilde G}(\tilde W_0)$, there holds $\mathfrak g=\mathrm{Ann}_{\tilde{\mathfrak g}}(\tilde w_0)$. In particular, the images of $\omega$ lie in $\mathfrak g$. This proves our claim. Moreover, it follows from the assumption on $\tilde\kappa$ that $(\mathcal G\to\tilde M,\omega)$ is locally the correspondence space of a Cartan geometry of type $(G,P)$. That is, $(\tilde{\mathcal G}\to\tilde M,\tilde\omega)$ is locally the result of the Fefferman-type functor applied to a Cartan geometry of type $(G,P)$. This completes the proof of the theorem.
\end{proof}
One may want to apply the above theorem recursively to obtain the characterization of a certain Fefferman-type construction. For example, denote by $G\subseteq\hat G\subseteq\tilde G$ the natural inclusions
\[SU(p+1,q+1)\subseteq U(p+1,q+1)\subseteq SO(2\,p+2,2\,q+2).\] The Fefferman-type construction associating to every CR geometry of type $(p,q)$ a conformal geometry of signature $(2\,p+1,2\,q+1)$ is defined by the canonical group homomorphism $G\to\tilde G$. There exists $\mathbb J\in\mathfrak{so}(2\,p+2,2\,q+2)$ such that $\mathrm{Stab}_{\tilde G}(\mathbb J)=\hat G$. Moreover, for any $0\neq\nu\in\wedge^{p+q+2}_{\mathbb C}\,\mathbb C^{p+1,q+1}$, there holds $\mathrm{Stab}_{\hat G}(\nu)=G$. Thus we arrive at the characterizing condition provided in \cite{CG06}*{2.5 Proposition}. Analogously, the Fefferman-type construction associating to every $(2\,n+1)$-dimensional torsion-free Lagrangian contact structure of signature $(n+1,n+1)$ is defined by the natural inclusion $SL(n+2,\mathbb R)\subseteq  SO(n+2,n+2)$. The subgroup of $SO(n+2,n+2)$ fixing a certain point $\mathbb K\in\mathfrak{so}(n+2,n+2)$ is $GL(n+2,\mathbb R)$ and the subgroup of $GL(n+2,\mathbb R)$ fixing any nonzero element of $\wedge^{n+2}\,\mathbb R$ is $SL(n+2,\mathbb R)$. This leads to a characterizing condition that is equivalent to \cite{Ham+17}*{Proposition 3.10}.


\section{Some parabolic background}\label{2}
Assume now that $G$ is a semisimple Lie group whose Lie algebra $\mathfrak g$ comes with a $|k|$-grading, i.e., there is a vector space decomposition $\mathfrak g=\bigoplus_{\ell=-k}^{k}\mathfrak p_\ell$ such that (i) $[\mathfrak p_i,\mathfrak p_j]\subseteq\mathfrak p_{i+j}$ for all $i,j\in\mathbb Z$, (ii) the Lie subalgebra $\mathfrak p_-=\bigoplus_{\ell=-k}^{-1}\mathfrak p_\ell$ is generated by $\mathfrak p_{-1}$ and (iii) $\mathfrak p_{\pm k}\neq 0$. We follow the convention that $\mathfrak p_i=0$ for $|i|>k$ and $\mathfrak p^i=\bigoplus_{\ell=i}^{k}\mathfrak p_\ell$.

Let $\mathfrak p=\mathfrak p^0$ and $P\subseteq G$ be a Lie subgroup with Lie algebra $\mathfrak p$. $P$ is said to be a parabolic subgroup of $G$. Note that the maximal parabolic subgroup of $G$ is
\[\{g\in G:Ad(g)(\mathfrak p^i)\subseteq\mathfrak p^i,-k\leq i\leq k\}.\]
The Levi subgroup of $P$ is given by
\[P_0=\{g\in P:Ad(g)(\mathfrak p_i)\subseteq\mathfrak p_i,-k\leq i\leq k\}.\]
It is the maximal subgroup of $P$ with Lie algebra $\mathfrak p_0$. Moreover, let $\mathfrak p_+=\mathfrak p^1$ and $P_+=\exp\,\mathfrak p_+$. Then $P_+$ is a normal subgroup of $P$ and $P/P_+\cong P_0$.

The Killing form of $\mathfrak g$ induces a $G$-invariant identification $\mathfrak g\cong\mathfrak g^*$, which further induces a $P$-module identification \[(\mathfrak g/\mathfrak p)^*\cong\mathfrak p_+.\] With this identification, the Kostant codifferential 
\begin{align}\label{Kostant codifferential}
\partial^*:\wedge^{n+1}\,(\mathfrak g/\mathfrak p)^*\otimes\mathfrak g\to\wedge^{n}\,(\mathfrak g/\mathfrak p)^*\otimes\mathfrak g
\end{align}
is determined by
\begin{align*}
    \partial^*(Z_0\wedge\dots\wedge Z_n\otimes A)&=\sum_{i=0}^n(-1)^{i+1}\,Z_0\wedge\dots\wedge\hat{Z_i}\wedge\dots\wedge Z_n\otimes[Z_i,A]\\
    &+\sum_{i<j}(-1)^{i+j}\,[Z_i,Z_j]\wedge Z_0\wedge\dots\wedge\hat{Z_i}\wedge\dots\wedge\hat{Z_j}\wedge\dots\wedge Z_n\otimes A
\end{align*}
for $Z_0,\dots Z_n\in\mathfrak p_+,A\in\mathfrak g$; see, e.g., \cite{Book}*{Subsection 3.1.11}. Here, the hat indicates that the term is omitted from a series. Note that $\partial^*$ is $P$-equivariant and $\partial^*\circ\partial^*=0$.

In particular, let $\{X^i\}\subseteq\mathfrak g/\mathfrak p$ be a basis with dual basis $\{Z_i\}\subseteq\mathfrak p_+$. For $\phi\in\wedge^2\,(\mathfrak g/\mathfrak p)^*\otimes\mathfrak g$ and $X\in\mathfrak g/\mathfrak p$, we have
\begin{align}\label{codiff}
    (\partial^*\phi)(X)=2\sum_i[Z_i,\phi(X,X^i)]-\sum_i\phi([Z_i,\tilde X],X^i)
\end{align}
where $\tilde X\in\mathfrak g$ is any lift of $X\in\mathfrak g/\mathfrak p$; see, e.g., \cite{Book}*{Lemma 3.1.11}. Note that if $\mathfrak g$ is $|1|$-graded, the second summand on the right-hand side vanishes as $[\mathfrak g,\mathfrak p_+]\subseteq\mathfrak p$.

Since $P\subseteq G$ is a parabolic subgroup, a Cartan geometry $(\mathcal G\to M,\omega)$ of type $(G,P)$ is called a parabolic geometry of type $(G,P)$. It is said to be regular if its curvature $\kappa:\mathcal G\to L(\wedge^2\,\mathfrak g/\mathfrak p,\mathfrak g)$ has homogeneity $\geq 1$, i.e., if $\kappa(u)(\mathfrak p^i,\mathfrak p^j)\subseteq \mathfrak p^{i+j+1}$ for all $u\in\mathcal G$. Note that this property holds automatically if $\mathfrak g$ is $|1|$-graded. On the other hand, $(\mathcal G,\omega)$ is said to be normal if $\partial^*\kappa=0$, i.e. 
\[\kappa:\mathcal G\to\ker\,\partial^*\subseteq L(\wedge^2\,\mathfrak g/\mathfrak p,\mathfrak g).\]
In this case, $\kappa$ projects to
\[\kappa_H:\mathcal G\to\ker\,\partial^*/\mathrm{im}\,\partial^*.\] 
This projection is called the harmonic curvature of $(\mathcal G\to M,\omega)$.

The Lie algebra homology differential
\begin{align}\label{partial}
    \partial:L(\wedge^n\,\mathfrak p_-,\mathfrak g)\to L(\wedge^{n+1}\,\mathfrak p_-,\mathfrak g)
\end{align}
given by
\begin{align*}(\partial\phi)(X^0,\dots,X^n)&=\sum_{i=0}^n(-1)^i\,[X^i,\phi(X^0,\dots,\hat X^i,\dots,X^n)]\\
&+\sum_{i<j}\phi([X^i,X^j],X^0,\dots,\hat X^i,\dots,\hat X^j,\dots,X^n)
\end{align*}
is $P_0$-equivariant; see, e.g., \cite{Book}*{Subsection 3.1.10}. Here, $X^0,\dots,X^n\in\mathfrak p_-$ and $\phi\in L(\wedge^n\,\mathfrak p_-,\mathfrak g)$. The hat indicates that the object is omitted. Then $\square=\partial\circ\partial^*+\partial^*\circ\partial$ is the Kostant Laplacian; see, e.g., \cite{Book}*{Subsection 3.1.11}. Recall from \cite{Book}*{Proposition 3.1.11} the $P_0$-invariant Hodge decomposition
\begin{align}\label{Hodge}
   L(\wedge^2\,\mathfrak g/\mathfrak p,\mathfrak g)=\mathrm{im}\,\partial^*\oplus\ker\,\square\oplus\mathrm{im}\,\partial 
\end{align}
where in addition, $\mathrm{im}\,\partial^*\oplus\ker\,\square=\ker\,\partial^*$ and $\ker\,\square\oplus\mathrm{im}\,\partial=\ker\,\partial$. In particular, one obtains an identification
\[\kappa_H(u)\in\ker\,\partial^*/\mathrm{im}\,\partial^*
\cong\ker\,\square\subseteq L(\wedge^2\,\mathfrak g/\mathfrak p,\mathfrak g)\]
for every $u\in\mathcal G$. Using this, one can rephrase the improved Bianchi identity as follows.
\begin{lemma}\label[lemma]{Bianchi}
(See \cite{RelBGG}*{Proposition 4.16 (3)})    Let $(\mathcal G,\omega)$ be a normal regular parabolic geometry of type $(G,P)$ with curvature \[\kappa:\mathcal G\to\ker\,\partial^*\subseteq L(\wedge^2\,\mathfrak g/\mathfrak p,\mathfrak g)\] and harmonic curvature $\kappa_H$. Let $\mathbb F$ be a $P$-submodule of $L(\wedge^2\,\mathfrak g/\mathfrak p,\mathfrak g)$ such that, viewing $\kappa_H(u)\in\ker\,\square$, we have
\[\kappa_H(u)\in\mathbb F\cap\ker\,\square\]
for all $u\in\mathcal G$ and such that $\mathbb F$ is stable under $\mathbb F$-insertions, i.e., for every $\varphi,\psi\in\mathbb F$, applying the Kostant codifferential $\partial^*$ to the map 
\[(X,Y,Z)\mapsto
\psi(\varphi(X,Y),Z)+\psi(\varphi(Y,Z),X)+\psi(\varphi(Z,X),Y),\]
where $X,Y,Z\in\mathfrak g/\mathfrak p$, one obtains an element in $\mathbb F$.

Then $\kappa(u)\in\mathbb F$ for all $u\in\mathcal G$.
\end{lemma}

\begin{convention}\label{convention}
From now on through the rest of the paper, we fix the notation for the following Lie groups and Lie algebras. Let
\[H\subseteq Q\subseteq P\subseteq G\] be the groups
\begin{align*}
    \left\{\left(\begin{array}{c|c|ccc}
 	*&*& &*&\\
 	\hline 0&1& &*&\\
 	\hline  & & & &\\
 	0&0&\ &*&\\
 	&&&&
 \end{array}\right)\right\}\subseteq\left\{\left(\begin{array}{c|c|ccc}
 	*&*& &*&\\
 	\hline 0&\mathbb R_+& &*&\\
 	\hline  & & & &\\
 	0&0&\ &*&\\
 	&&&&
 \end{array}\right)\right\}\subseteq\left\{\left(\begin{array}{c|c|ccc}
 	*&*& &*&\\
 	\hline *&*& &*&\\
 	\hline  & & & &\\
 	0&0&\ &*&\\
 	&&&&
 \end{array}\right)\right\}\subseteq SL(n+2,\mathbb R),
\end{align*}
where the block size is $(1,1,n)\times(1,1,n)$. Their Lie algebras are denoted by
\[\mathfrak h\subseteq\mathfrak q\subseteq\mathfrak p\subseteq\mathfrak g,\]
respectively. Note that $Q$ and $P$ are parabolic subgroups of the semisimple Lie group $G$ arising from the grading decompositions
\begin{align*}
\mathfrak g=\left(\begin{array}{c|c|ccc}
	\mathfrak q_0&\mathfrak q_1^E& &\mathfrak q_2&\\
	\hline \mathfrak q_{-1}^E&\mathfrak q_0& &\mathfrak q_1^V&\\
	\hline  & & & &\\
	\mathfrak q_{-2}&\mathfrak q_{-1}^V&\ &\mathfrak q_{0}&\\
	&&&&
\end{array}\right)
\quad\text{and}\quad
\mathfrak g=\left(\begin{array}{c|ccc}
	\mathfrak p_0&&\mathfrak p_1&\\
 \hline&&&\\
 \mathfrak p_{-1}&&\mathfrak p_0&\\
 &&&
\end{array}\right),
\end{align*}
respectively. We also fix the notation for
\[\tilde P\subseteq\tilde G=SL(n+3,\mathbb R)\]
where $\tilde P\subseteq\tilde G$ is given by $P\subseteq G$ with $n$ increased by $1$. Their Lie algebras are $\tilde{\mathfrak p}\subseteq\tilde{\mathfrak g}$, respectively. The grading
\begin{align*}
\tilde{\mathfrak g}=\tilde{\mathfrak p}_{-1}\oplus\tilde{\mathfrak p}_0\oplus\tilde{\mathfrak p}_1.
\end{align*}
is precisely the grading of $(\mathfrak g,\mathfrak p)$ with $n$ increased by $1$.
\end{convention}
\subsection{Almost Grassmann structures}\label{2.3} The following categorical correspondence follows from, e.g., \cite{Book}*{Subsection 4.1.3 and Theorem 3.1.14}.

\begin{lemma}\label[lemma]{AG parabolic}
Notation as in \cref{convention}. Every parabolic geometry $(\mathcal G\to M,\omega)$ of type $(G,P)$ has an underlying almost Grassmann structure 
\[(M,E=\mathcal G\times_{P}\mathbb R^2,F=\mathcal G\times_{P}\mathbb R^{n+2}/\mathbb R^2)\] 
of type $(2,n)$. Conversely, for every almost Grassmann structure $(M,E,F)$ of type $(2,n)$, there is, up to isomorphisms, a unique normal parabolic geometry of type $(G,P)$ whose underlying geometry is $(M,E,F)$.
\end{lemma}
Let $(\mathcal G\to M,\omega)$ be the normal parabolic geometry of type $(G,P)$ canonically associated to a given almost Grassmann structure $(M,E,F)$. Recall from, e.g., \cite{Book}*{Step (E) in Subsection 4.1.3} that for $n>2$, $\kappa_H$ decomposes into 
\begin{equation}\label{harmonic components AG}
\begin{aligned}
    \tau&\in\Gamma\left((\mathrm{Sym}^2\,E\otimes E^*)_o\otimes(\wedge^2\,F^*\otimes F)_o\right)\subseteq\Omega^2(M,TM)\quad\text{and}\\
    \rho&\in\Gamma(\wedge^2\,E\otimes \mathrm{Sym}^2\,F^*\otimes\mathfrak {sl}(F))\subseteq\Omega^2(M,\mathfrak {sl}(F)).
\end{aligned}
\end{equation}
Here, the natural decomposition
\[\wedge^2\,T^*M=\wedge^2\,(E\otimes F^*)=(\mathrm{Sym}^2\,E\otimes\wedge^2\,F^*)\oplus(\wedge^2\,E\otimes \mathrm{Sym}^2\,F^*)\]
is used and
\begin{align*}
(\mathrm{Sym}^2\,E\otimes E^*)_o\subseteq \mathrm{Sym}^2\,E\otimes E^*
\quad\text{and}\quad
(\wedge^2\,F^*\otimes F)_o\subseteq \wedge^2\,F^*\otimes F
\end{align*}
denote the trace-free component. Stretching notation, we call $\tau$ the harmonic torsion of $(M,E,F)$ and $\rho$ the harmonic curvature of $(M,E,F)$. On the other hand, recall from, e.g., \cite{Book}*{Theorem 3.1.12} that the projection of the curvature $\kappa$ of $\omega$ to 
\begin{align}\label{T def}
T\in\Omega^2_{hor}(\mathcal G,\mathfrak g/\mathfrak p)^P\cong\Omega^2(M,TM)    
\end{align}
coincides with $\tau$ and is called the torsion of $(M,E,F)$. If $\tau=0$, we say that $(M,E,F)$ is torsion-free.

Although Weyl structures and their induced objects are defined for every parabolic geometry, see, e.g., \cite{Book}*{Subsections 5.1-5.2}, we only recall them in the setting of almost Grassmann structure. The Lie algebra grading induces the $P_0$-equivariant decompositions
\begin{align*}
    \omega&=\omega_{\mathfrak p_{-1}}+\omega_{\mathfrak p_0}+\omega_{\mathfrak p_1}\\
    \kappa&=\kappa_{\mathfrak p_{-1}}+\kappa_{\mathfrak p_0}+\kappa_{\mathfrak p_1}
\end{align*}
of the Cartan connection $\omega$ as well as its curvature $\kappa$.

Observe that $\mathcal G_0=\mathcal G/P_+$ is a principal bundle with structure group $P_0=P/P_+$. Recall from, e.g., \cite{Book}*{Proposition 3.1.15\,(2)} that $\omega_{\mathfrak p_{-1}}$ descends to a one-form $\theta\in\Omega^1(\mathcal G_0,\mathfrak p_{-1})^{P_0}$. In fact, $(\mathcal G_0,\theta)$ is the $G$-structure of $(M,E,F)$. In particular, the natural projection $\mathcal G\twoheadrightarrow \mathcal G_0$ induces identifications 
\begin{align*}
E&\cong\mathcal G\times_P\mathbb R^2\cong\mathcal G_0\times_{P_0}\mathbb R^2\\
F&\cong\mathcal G\times_P\mathbb R^{n+2}/\mathbb R^2\cong\mathcal G_0\times_{P_0}\mathbb R^n\\
TM&\cong\mathcal G\times_P\mathfrak g/\mathfrak p\cong\mathcal G_0\times_{P_0}\mathfrak p_{-1}\\
T^*M&\cong\mathcal G\times_P\mathfrak p_1\cong\mathcal G_0\times_{P_0}\mathfrak p_1
\end{align*}
so that
\begin{align*}
    \mathcal G_0\times_{P_0}\mathfrak p_0\cong \mathfrak s(L(E,E)\oplus L(F,F))
\end{align*}
where
\begin{align}\label{frak s}
   \mathfrak s(L(E,E)\oplus L(F,F)))=\Big(L(E,E)\oplus L(F,F)\Big)\cap\mathfrak{sl}(E\oplus F).
\end{align}

Recall from, e.g., \cite{Book}*{Subsection 5.1.1} that there exist $P_0$-equivariant sections of the natural projection $\mathcal G\twoheadrightarrow\mathcal G_0$. We call such a section $\sigma:\mathcal G_0\to\mathcal G$ a Weyl structure of $(M,E,F)$.
Observe that $\sigma^*\omega_{\mathfrak p_{-1}}$ coincides with the one-form $\theta$ defining the $G$-structure and that $\sigma^*\kappa_{\mathfrak p_{-1}}$ coincides with the torsion \eqref{T def}. Hence they are independent of the choice of Weyl structures. 
We obtain the following objects associated to a Weyl structure $\sigma$; see, e.g., \cite{Book}*{Subsections 5.1.2 and 5.2.3}.
\begin{equation}\label{induced Weyl objects}
   \begin{aligned}
\sigma^*\omega_{\mathfrak p_{0}}&\in\Omega^1(\mathcal G_0,\mathfrak p_{0})^{P_0}\text{ (the Weyl connection)}\\
\mathrm P=\sigma^*\omega_{\mathfrak p_{1}}&\in\Omega^1_{hor}(\mathcal G_0,\mathfrak p_{1})^{P_0}\cong \Omega^1(M,T^*M)\text{ (the Rho tensor)}\\
(W,W')=\sigma^*\kappa_{\mathfrak p_{0}}&\in\Omega^2_{hor}(\mathcal G_0,\mathfrak p_{0})^{P_0}\\
&\cong \Omega^2(M,\mathfrak s(L(E,E)\oplus L(F,F)))\text{ (the Weyl tensor)}\\
Y=\sigma^*\kappa_{\mathfrak p_{1}}&\in\Omega^2_{hor}(\mathcal G_0,\mathfrak p_{1})^{P_0}\cong \Omega^2(M,T^*M)\text{ (the Cotton-York tensor)}.
\end{aligned}
\end{equation}
The Weyl connection $\sigma^*\omega_{\mathfrak p_{0}}$ is a principal connection on $\mathcal G_0$ and the Weyl tensor $(W,W')$ is decomposed into
\begin{align*}
    W\in\Omega^2(M,L(E,E))\quad\text{and}\quad W'\in\Omega^2(M,L(F,F)).
\end{align*}
Following the convention of Penrose abstract index notation, see, e.g., \cite{Book}*{Subsection 4.1.3}, we use unprimed upper indices for $E$ and lower indices for $E^*$ and primed upper indices for $F$ and lower indices for $F^*$. Recall that a contraction is denoted by writing the same index variable in the two slots that are being contracted. Thus
\begin{align*}
   W=W{}^A_{A'}{}^B_{B'}{}^{C}_{D},\quad W'=W'{}^A_{A'}{}^B_{B'}{}^{C'}_{D'}\quad\text{and}\quad\tau=\tau{}^A_{A'}{}^B_{B'}{}^{C'}_D
\end{align*}
where $\tau$ is the harmonic torsion of $(M,E,F)$.

The following lemma has been stated in \cite{HSSS12B}*{Equation 25}. For an explicit proof, we refer the reader to \cite{Guo}*{Proposition 4}.
\begin{lemma}\label[lemma]{Weyl tensor}
We specify the following contractions.
\begin{align*}
    tr(W){}^A_{A'}{}^B_{B'}=W{}^A_{A'}{}^I_{B'}{}^{B}_{I}\quad tr(W'){}^A_{A'}{}^B_{B'}=W'{}^A_{A'}{}^B_{I'}{}^{I'}_{B'}\quad\text{and}\quad tr(\iota_\tau\tau){}^A_{A'}{}^B_{B'}=\tau{}^I_{I'}{}^A_{A'}{}^{J'}_J\,\tau{}^J_{J'}{}^B_{B'}{}^{I'}_I.
\end{align*}
Then 
\[tr(W)=tr(W')\in\Gamma(\mathrm{Sym}^2\,T^*M)\]
with
\begin{align*}
            tr(\iota_\tau\tau){}^{(A}_{(A'}{}^{B)}_{B')}&=n\,tr(W){}^{(A}_{(A'}{}^{B)}_{B')}&
            tr(\iota_\tau\tau){}^{[A}_{[A'}{}^{B]}_{B']}&=(n+4)\,tr(W){}^{[A}_{[A'}{}^{B]}_{B']}\\
            tr(\iota_\tau\tau){}^{(A}_{[A'}{}^{B)}_{B']}&=tr(W){}^{(A}_{[A'}{}^{B)}_{B']}=0&
            tr(\iota_\tau\tau){}^{[A}_{(A'}{}^{B]}_{B')}&=tr(W){}^{[A}_{(A'}{}^{B]}_{B')}=0.
\end{align*}
Here, $[\,\cdot\,,\,\cdot\,]$ and $(\,\cdot\,,\,\cdot\,)$ are the alternation and the symmetrization of indices, respectively. 
\end{lemma}


Recall from, e.g., \cite{Book}*{Theorem 5.2.3 (1)} that Weyl connections of $(M,E,F)$ are exactly the principal connections on $\mathcal G_0$ which induce a connection on $TM\cong E^*\otimes F$ with $\partial^*$-closed torsion. Note that the principal connections on $\mathcal G_0$ are equivalent to all pairs $(\nabla^E,\nabla^F)$ of connections on $(E,F)$ compatible with the defining isomorphism $\wedge^2\, E^*\cong \wedge^n\,F$. Moreover,  since $(\mathcal G\to M,\omega)$ is $|1|$-graded, there is a one-to-one correspondence between Weyl structures and Weyl connections of $(\mathcal G\to M,\omega)$; see, e.g., \cite{Book}*{Proposition 5.1.1 and Proposition 5.1.6}. Thus we derive the following one-to-one correspondence between Weyl structures and the admissible pairs $(\nabla^E,\nabla^F)$ mentioned above.
\begin{lemma}\label[lemma]{AG connections}
Let $(M,E,F)$ be an almost Grassmann structure of type $(2,n)$. Given a Weyl structure, its Weyl connection $(\nabla^E,\nabla^F)$ on $(E,F)$ satisfies the following properties.
\begin{enumerate}
    \item[$\bullet$] 
    $(\nabla^E,\nabla^F)$ is compatible with the defining isomorphism $\wedge^2\,E^*\cong\wedge^n\,F$ and
    \item[$\bullet$] 
    $(\nabla^E,\nabla^F)$ induces a connection on $TM\cong E^*\otimes F$ whose torsion is exactly the harmonic torsion $\tau$ of $(M,E,F)$; see \eqref{harmonic components AG}.
\end{enumerate}

Conversely, given any pair of linear connections $(\nabla^E,\nabla^F)$ on $(E,F)$ satisfying these two properties, there is a unique Weyl structure of $(M,E,F)$ whose Weyl connection on $(E,F)$ is $(\nabla^E,\nabla^F)$.
\end{lemma}
We may use a Weyl structure and its induced objects to describe tractor connections; see \cref{section abstract Fefferman}. Let $\sigma:\mathcal G_0\to\mathcal G$ be a Weyl structure of $(\mathcal G\to M,\omega)$. First, it decomposes a tractor bundle into simpler bundles. For example, the adjoint representation of $G$ restricted to $P_0$ admits decomposition $\mathfrak g=\mathfrak p_{-1}\oplus\mathfrak p_0\oplus\mathfrak p_1$. Applying the functor $\mathcal G_0\times_{P_0}\,\cdot\,$ to this decomposition results in
\begin{align}\label{adjoint bundle decomposition}
    \mathcal AM\stackrel{\sigma}{\cong}\mathcal G_0\times_{P_0}\mathfrak g=TM\oplus \mathfrak s(L(E,E)\oplus L(F,F))\oplus T^*M;
\end{align}
see \eqref{frak s}.
Similarly, for the standard tractor bundle induced from the standard representation of $G$ on $\mathbb R^{n+2}$, the $P_0$-module decomposition $\mathbb R^{n+2}=\mathbb R^2\oplus\mathbb R^n$ induces
\begin{align}\label{standard tractor splitting}
    \mathcal T\stackrel{\sigma}{\cong}\mathcal G_0\times_{P_0}\mathbb R^{n+2}=E\oplus F.
\end{align}

Second, consider the the horizontal lift $\mathfrak X(M)\ni\xi\mapsto\xi^h\in\mathfrak X(\mathcal G_0)$ of the Weyl connection $\nabla$ as well as the Rho tensor $\mathrm{P}$. Consider the tractor bundle $\mathcal W=\mathcal G\times_PW$ induced by a representation $\rho:G\to GL(W)$ with derivative $\rho':\mathfrak g\to L(W, W)$. Then the tractor connection $\nabla^{\mathcal W}$ on $\mathcal W$ can be expressed as
\begin{align}\label{tractor connection priliminary}
    \nabla^{\mathcal W}_{\xi} s
    &=\xi^h.f_s+
    \rho'(\sigma^*\omega_{\mathfrak p_{-1}}(\xi^h))(f_s)+\rho'(\sigma^*\omega_{\mathfrak p_1}(\xi^h))(f_s)\nonumber\\
    &=\nabla_\xi s +\xi\bullet s+\mathrm{P}(\xi)\bullet s;
\end{align}
see, e.g., \cite{Book}*{Proposition 5.1.10 (2)}.
Here, $s\in\Gamma(\mathcal W)$ corresponds to $f_s\in C^\infty(\mathcal G_0,W)^{P_0}$, $\xi\in\mathfrak X(M)$ and $\bullet:TM\times\mathcal W\to \mathcal W$ and $\bullet:T^*M\times\mathcal W\to \mathcal W$ are the $\mathcal G_0$-associated maps of $\rho'|_{\mathfrak p_{-1}}$ and $\rho'|_{\mathfrak p_{1}}$, respectively.

\subsection{Path geometries}\label{2.4}
Let $\mathcal D$ be a distribution on a manifold $M$. The Levi bracket \[\mathcal L:\mathcal D\times\mathcal D\to TM/\mathcal D\] is the skew-symmetric $C^\infty(M)$-bilinear map descended from the Lie bracket of sections of $\mathcal D$. Assume now that $M^{2\,n+1}$ is $(2\,n+1)$-dimensional. Recall that a generalized path geometry on $M$ is a distribution $E\oplus V\subseteq TM$ on $M$ where $E$ and $V$ are of rank $1$ and rank $n$, respectively, such that the Levi bracket
\[\mathcal L:\wedge^2(E\oplus V)\to TM/(E\oplus V)\]
has kernel $\wedge^2\,V$ and hence induces an isomorphism 
\[E\otimes V\xrightarrow[\cong]{\mathcal L}TM/(E\oplus V);\]
see, e.g., \cite{Book}*{Definition 4.4.3}.

Notice that an orientation on $E\oplus V$ may be encoded as an orientation on $E^*\otimes\wedge^n\,V$ via the canonical orientation on $\otimes^2\, E$.

\begin{lemma}\label[lemma]{path parabolic}
Notation as in \cref{convention}. Let $(\mathcal G\to M,\omega)$ be a regular parabolic geometry of type $(G,Q)$. Then $M$ has dimension $2\,n+1$ and 
\[(M,E=\mathcal G\times_Q(\mathfrak q_{-1}^E\oplus\mathfrak q)/\mathfrak q,V=\mathcal G\times_Q(\mathfrak q_{-1}^V\oplus\mathfrak q)/\mathfrak q)\]
is a generalized path geometry with oriented rank-$(n+1)$-distribution $E\oplus V\subseteq TM$ or, equivalently, with oriented line bundle $E^*\otimes\wedge^n\,V$.
    
Conversely, for any generalized path geometry on a $(2\,n+1)$-dimensional manifold with oriented rank-$(n+1)$-distribution or, equivalently, with oriented line bundle $E^*\otimes\wedge^n\,V$, there is a unique normal regular parabolic geometry $(\mathcal G\to M,\omega)$ of type $(G,Q)$ with such an underlying geometry.
\end{lemma}
\begin{proof}
Denote by $Aut_{gr}(\mathfrak q_-)$ the group of Lie algebra automorphisms on $\mathfrak q_-=\mathfrak q_{-2}\oplus\mathfrak q_{-1}$ preserving the grading decomposition. Using this identification, one observe that the adjoint representation restricted to the Levi subgroup $Q_0$ of $Q$ induces an identification
\begin{align*}
    Q_0\xrightarrow[\cong]{Ad}&\{\varphi\in Aut_{gr}(\mathfrak q_-):\varphi(\mathfrak q_{-1}^E)\subseteq\mathfrak q_{-1}^E,\varphi(\mathfrak q_{-1}^V)\subseteq\mathfrak q_{-1}^V\text{ and }\varphi|_{\mathfrak q_{-1}}\in GL_+(\mathfrak q_{-1})\}.
\end{align*}
The fact that $\varphi|_{\mathfrak q_{-1}}$ is orientation-preserving is the only difference to the case dealt in \cite{Book}*{Proposition 4.4.3}, so the proof can be completed as there.
\end{proof}

Let $(\mathcal G\to M,\omega)$ be the normal regular parabolic geometry of type $(G,Q)$ canonically associated to a given path geometry $(M,E,V)$ with oriented $E\oplus V$. Denote by $\kappa_H$ its harmonic curvature. Since $V$ is involutive for a path geometry $(M,E,V)$, recall from, e.g., \cite{Book}*{Subsection 4.4.3} that for $n\geq 2$, $\kappa_H$ decomposes to
\begin{equation}\label{harmonic components path}\begin{aligned}
    \tau&\in\Gamma(E^*\otimes(TM/(E\oplus V))^*\otimes V)\quad\text{and}\\
    \rho&\in\Gamma(V^*\otimes(TM/(E\oplus V))^*\otimes L(V,V)).
\end{aligned}
\end{equation}

Stretching notation, we call $\tau$ the harmonic torsion of $(M,E,V)$ and $\rho$ the harmonic curvature of $(M,E,V)$. Analogous to the setting of almost Grassmann structure, $\tau$ vanishes if and only if the torsion $T$ of $(\mathcal G\to M,\omega)$ vanishes. In this case we say that the path geometry $(M,E,V)$ is torsion-free. Here, $T$ is defined by the projection of $\kappa\in\Omega^2_{hor}(\mathcal G,\mathfrak g)^Q$ to 
\[T\in\Omega^2_{hor}(\mathcal G,\mathfrak g/\mathfrak q)^Q\cong\Omega^2(M,TM).\] 

\section{The Fefferman-type construction on path geometries}\label{The Fefferman-type constructions}
We follow \cref{convention} for the notation of Lie groups, Lie algebras and gradings. Denote by $e_1,\dots,e_{n+3}$ the standard basis of $\mathbb R^{n+3}$. The decomposition of $\mathbb R^{n+3}$ into the line $\mathbb R e_3$ and the hyperplane 
\[\mathrm{span}\{e_1,e_2+e_3,e_4,\dots,e_{n+3}\}\] induces a representation of the Lie group $G$ on $\mathbb R^{n+3}$. We write this representation as the injective group homomorphism
\begin{align}\label{i}
i:G&\to \tilde G,\\
	\left(\begin{array}{c|c|ccc}
		x_{11}&x_{12}& &x_{13}&\\
		\hline x_{21}&x_{22}& &x_{23}&\\
		\hline  & & & &\\
		x_{31}&x_{32}&\ &x_{33}&\\
		&&&&
	\end{array}\right)&\mapsto\left(\begin{array}{c|c|c|ccc}
		x_{11}&x_{12}&0& &x_{13}&\\
		\hline x_{21}&x_{22}&0& &x_{23}&\\
		\hline x_{21}&x_{22}-1&1& &x_{23}&\\
		\hline &&&&\\
		x_{31}&x_{32}&0& &x_{33}&\\
		&&&&&
	\end{array}\right)\nonumber
\end{align}
where the block sizes are $(1,1,n)\times(1,1,n)$ and $(1,1,1,n)\times(1,1,1,n)$, respectively. Observe that $G$ acts locally transitively on $\tilde G/\tilde P$, that $H=i^{-1}(\tilde P)$ and that $H\subseteq Q$. Recall that every path geometry \[(M,E,V)\] where $M$ is $(2\,n+1)$-dimensional and $E\oplus V$ is oriented is canonically associated to a normal regular parabolic geometry \[(\mathcal G\to M,\omega)\] of type $(G,Q)$. The Fefferman-type functor defined by $i:G\to\tilde G$ associates to $(\mathcal G\to M,\omega)$ a parabolic geometry \[(\tilde{\mathcal G}\to\tilde M,\tilde\omega)\] of type $(\tilde G,\tilde P)$ with underlying almost Grassmann structure \[(\tilde M,\tilde E,\tilde F)\] of type $(2,n+1)$. The construction
\begin{align}\label{functor1}
(M,E,V)\mapsto(\tilde M,\tilde E,\tilde F)
\end{align}
is clearly functorial.

We denote by
\[p:\tilde M\to M\]
the natural projection and by $\kappa$ and $\tilde\kappa$ the curvatures of $\omega$ and $\tilde\omega$, respectively. Denote by 
\[j:\mathcal G\to\tilde{\mathcal G}\]
the natural map defining $\tilde{\mathcal G}$ and by
\[\pi:\tilde{\mathfrak g}/\tilde{\mathfrak p}\twoheadrightarrow\mathfrak g/\mathfrak q\]
the projection canonically induced by the Lie algebra map $i'$ of \eqref{i}; see \cref{section abstract Fefferman}.

For the statement of the following theorem, recall \cref{path parabolic}.

\begin{theorem}\label[theorem]{bundles}
Notation as above. Then
\begin{align*}
\tilde M= GL_+(E^*\otimes\wedge^n\, V)   
\end{align*}
as a principal $\mathbb R_+$-bundle over $M$ and
\begin{align}\label{aux vec bundles}
    \tilde E^*=(Tp)^{-1}(E)\quad\text{and}\quad
     \tilde F=(Tp)^{-1}(V)
\end{align}
as vector bundles over $\tilde M$.

Moreover, the construction \eqref{functor1} is independent of the choice of orientation on $E^*\otimes\wedge^n\, V$. In particular, it leads to a construction associating to every path geometry $(M',E',V')$ without auxiliary orientation an almost Grassmann structure $(\tilde M',\tilde E',\tilde F')$ where $\tilde M'=GL(E'^{*}\otimes\wedge^nV')/\mathbb Z_2$ and $\tilde E'^*,\tilde F'\subseteq T\tilde M'$ are preimages of $E',V'\subseteq TM'$, respectively. Here, $\mathbb Z_2$ denotes the group generated by the negation map $v\mapsto-v$ on the line bundle $E'^{*}\otimes\wedge^nV'$. 
\end{theorem}
\begin{proof}
Observe that the adjoint representation induces a surjective group homomorphism
\[Q\twoheadrightarrow GL_+\left(\left((\mathfrak q_{-1}^E\oplus\mathfrak q)/\mathfrak q\right)^*\otimes\wedge^n\,\left((\mathfrak q_{-1}^V\oplus\mathfrak q)/\mathfrak q\right)\right)=\mathbb R_+\]
with kernel $H$. Hence 
\[\tilde M=\mathcal G/H=\mathcal G\times_HQ/H= GL_+(E^*\otimes\wedge^n\,V)\]
as a principal $\mathbb R_+$-bundle over $M$.

Next, we construct the identifications in \eqref{aux vec bundles}. Notice that
\begin{equation}\label{lifts of E and V}
\begin{aligned}
\begin{cases}
    (Tp)^{-1}(E)=\mathcal G\times_H (\mathfrak q_{-1}^E\oplus\mathfrak q)/\mathfrak h\quad\text{and}\\
    (Tp)^{-1}(V)=\mathcal G\times_H(\mathfrak q_{-1}^V\oplus\mathfrak q)/\mathfrak h.
\end{cases}
\end{aligned}
\end{equation}
On the other hand, consider the filtration $\mathbb R^{n+2}\supseteq\mathbb R^2\supseteq\mathbb R$ induced by the standard representation of $H$ and the filtration $\mathbb R^{n+3}\supseteq\mathbb R^2$ induced by the standard representation of $\tilde P$. Observe that the principal bundle map $j:\mathcal G\to\tilde{\mathcal G}$ in the Fefferman-type functor canonically induces
\begin{equation}\label{assoc H bundles}
\begin{aligned}
\begin{cases}
\tilde E^*=\tilde{\mathcal G}\times_{\tilde P}\mathbb R^{2*}&=\mathcal G\times_H\mathbb R^{2*}
    \quad\text{and}\\
\tilde F=\tilde{\mathcal G}\times_{\tilde P}\mathbb R^{n+3}/\mathbb R^2&=\mathcal G\times_H\mathbb R^{n+2}/\mathbb R.    
\end{cases}
\end{aligned}
\end{equation}


Using the $H$-module isomorphisms $(\mathfrak q_{-1}^E\oplus\mathfrak q)/\mathfrak h\cong\mathbb R^{2*}$ and
$(\mathfrak q_{-1}^V\oplus\mathfrak q)/\mathfrak h\cong\mathbb R^{n+2}/\mathbb R$, we obtain \eqref{aux vec bundles} as desired. This proves the first assertion of the theorem.

Assume now that $(M,E,V)$ arises from a path geometry $(M',E',V')$ without auxiliary orientation where $M\to M'$ is the two-fold covering encoding the pair of orientations on $E'^{*}\otimes\wedge^nV'$, $E$ is the pullback of $E'$ and $V$ is the pullback of $V'$. Applying \eqref{functor1} to $(M,E,V)$ results in $(\tilde M,\tilde E,\tilde F)$ where $\tilde M=GL_+(E^*\otimes\wedge^n\,V)=GL(E'^*\otimes\wedge^n\,V')$. It suffices to construct an almost Grassmann structure $(\tilde M',\tilde E',\tilde F')$ where $\tilde M'=GL(E'^{*}\otimes\wedge^nV')/\mathbb Z_2$ and $(\tilde M,\tilde E,\tilde F)$ is the pullback of $(\tilde M',\tilde E',\tilde F')$.

Indeed, denote by $\hat Q\subseteq PGL(n+2,\mathbb R)$ the maximal subgroup with Lie algebra $\mathfrak q$. We lift $\hat Q$ along the natural projection $GL(n+2,\mathbb R)\twoheadrightarrow PGL(n+2,\mathbb R)$ to the block upper triangular subgroup of $GL(n+2,\mathbb R)$ with block diagonal entries $\mathrm{diag}(a,b,C)$ such that $b>0$ and $a\,b\,\det C=\pm 1$. This choice provides a unique representative for each element of $\hat Q$. This way, we obtain an identification
\[\hat Q\cong \mathrm{diag}(\pm1,1,\mathbb I_n)\,Q=Q\, \mathrm{diag}(\pm1,1,\mathbb I_n)\subseteq GL(n+2,\mathbb R).\]
Now let $\hat H\trianglelefteq\hat Q$ be the kernel of
\[\hat Q\twoheadrightarrow GL\left(\left((\mathfrak q_{-1}^E\oplus\mathfrak q)/\mathfrak q\right)^*\otimes\wedge^n\,\left((\mathfrak q_{-1}^V\oplus\mathfrak q)/\mathfrak q\right)\right).\]
The identification above restricted to $\hat H$ induces
\[\hat H\cong \mathrm{diag}(\pm1,1,\mathbb I_n)\,H=H\, \mathrm{diag}(\pm1,1,\mathbb I_n).\]
Recall from, e.g., \cite{Book}*{Proposition 4.4.3} that $(M',E',V')$ comes with a unique normal regular parabolic geometry $(\mathcal G'\to M',\omega')$ of type $(PGL(n+2,\mathbb R),\hat Q)$. Then
\[(\mathcal G'\to\mathcal G'/Q,\omega')\]
is the normal regular parabolic geometry of type $(G,Q)$ corresponding to $(M,E,V)$. Similar to the previous step, the Cartan geometry
\[(\mathcal G'\to\mathcal G'/\hat H,\omega')\]
has the desired underlying almost Grassmann structure $(\tilde M',\tilde E',\tilde F')$ which clearly pulls back to the underlying geometry $(\tilde M,\tilde E,\tilde F)$ of
\[(\mathcal G'\to\mathcal G'/H,\omega').\]
This completes the proof of the theorem. Note that the second half of the proof can be rephrased as an extension functor extending the Fefferman-type construction \eqref{i}; see \cite{CZ09}.
\end{proof}
\begin{remark}\label{distinguished sections}
Although we have constructed $\tilde E^*$ and $\tilde F$ as distributions $(Tp)^{-1}(E)$ and $(Tp)^{-1}(V)$ of $\tilde M$, respectively, it offers more clarity to think of $\tilde E^*$ and $\tilde F$ as auxiliary vector bundles. Observe that for all $a\in\mathbb R$, $(0,a)\in\mathbb R^{2*}$ and $(a,0^n)^t\in\mathbb R^{n+2}/\mathbb R$ are fixed by $H$. We specify the sections 
\begin{align*}
\varphi=\mathcal G\times_H(0,1)\in\Gamma(\tilde E^*)
\quad\text{and}\quad
\eta=\mathcal G\times_H(-1,0^n)^t\in\Gamma(\tilde F);
\end{align*}
see \eqref{assoc H bundles}. Denote by $\zeta_a\in\mathfrak X(\tilde M)$ the fundamental vector field of $p:\tilde M\to M$ generated by $a\in\mathbb R$. Then the defining isomorphism $\tilde E^*\otimes\tilde F\cong T\tilde M$ of $(\tilde M, \tilde E,\tilde F)$ restricts to
\begin{align*}
\tilde E^*\otimes\eta=(Tp)^{-1}(E),\quad\varphi\otimes\tilde F=(Tp)^{-1}(V)
\quad\text{and}\quad
\varphi\otimes\eta=\zeta_{-1}.
\end{align*}
\end{remark}


Let $N$ be an $(n+1)$-dimensional manifold with full frame bundle $\mathcal PN$. For $s\in\mathbb R$, recall that the $(-\frac{s}{n+2})$-density bundle 
\[\mathcal E(s)=\mathcal PN\times_{\rho(s)}\mathbb R\]
is associated to the representation
\begin{align*}
    \rho(s):GL(n+1,\mathbb R)&\to GL_+(\mathbb R)\\
    C&\mapsto|\det(C)|^{s/(n+2)};
\end{align*}
see, e.g., \cite{Book}*{Subsection 4.1.5}. For example, if $N$ is an oriented projective structure, the standard tractor bundle $\mathcal T$ of the geometry comes with a canonical short exact sequence
\[0\to\mathcal E(-1)\to\mathcal T\to TN\otimes\mathcal E(-1)\to 0;\]
see, e.g., \cite{Book}*{Subsection 5.2.6}. Now let $N$ be any $(n+1)$-dimensional manifold. Denote by $\mathcal PTN\to N$ the projectivized tangent bundle of $N$ and \[V\subseteq\mathcal H\subseteq T\mathcal PTN\] the vertical distribution and the tautological distribution, respectively. Recall that any line subbundle $E\subseteq\mathcal H$ such that $\mathcal H=E\oplus V$ defines a full path geometry $(\mathcal PTN,E,V)$.

In \cite{CrSa}*{Section 3}, the authors associate to every full path geometry $(\mathcal PTN,E,V)$ with oriented $N$ a Segre structure on the weighted slit tangent bundle $TN\otimes\mathcal E(-1)\setminus\{0\}$. We show that this space can be identified with our Fefferman-type space. Assume now that $N$ is not necessarily oriented. We denote by
\[\widehat{\mathcal PTN}=(\wedge^{n+1}\,\mathcal H\setminus\{0\})/\mathbb R_+\to\mathcal PTN\]
the two-fold covering encoding the pair of orientations on $\mathcal H\subseteq T\mathcal PTN$. Then every distribution of $\mathcal PTN$ canonically lifts to a distribution of $\widehat{\mathcal PTN}$, which we denote by the same notation. In particular, every full path geometry $(\mathcal PTN,E,V)$ canonically lifts to the path geometry $(\widehat{\mathcal PTN},E,V)$ where \[E\oplus V=\mathcal H\subseteq T\widehat{\mathcal PTN}\] is oriented.
\begin{theorem}\label[theorem]{Compare with CrSa}
Given a full path geometry $(\mathcal PTN,E,V)$, the Fefferman-type space of $(\widehat{\mathcal PTN},E,V)$ is isomorphic to $TN\otimes\mathcal E(-1)\setminus\{0\}$, viewed as a fiber bundle over $N$ with standard fiber $\mathbb R^{n+1}\setminus\{0\}$.
\end{theorem}

\begin{proof}
We first prove the statement of the theorem for every full path geometry on $N$ that arises from a certain correspondence space. Indeed, up to shrinking $N$, we may assume without loss of generality that there exists a normal parabolic geometry
\[(\mathcal G\to N,\omega)\]
of type $(SL(n+2,\mathbb R),R)$. Here, $R\subseteq SL(n+2,\mathbb R)=G$ is the stabilizer of $(\mathbb R, 0^{n+1})^t\subseteq\mathbb R^{n+2}$. We denote the Lie algebra of $R$ by $\mathfrak r$. It is well-known that $\mathcal E(s)\cong\mathcal G\times_R\mathbb R$ where the representation of $R$ on $\mathbb R$ is given by
\begin{align*}
    R\ni \left(\begin{array}{c|c|ccc}
 	a&*& &*&\\
 	\hline 0&*& &*&\\
 	\hline  & & & &\\
 	0&*&\ &*&\\
 	&&&&
 \end{array}\right)\mapsto |a|^{-s}\in GL_+(\mathbb R).
\end{align*}
The block size is $(1,1,n)\times(1,1,n)$. Note that the above homomorphism is the result of the composition
\[R\xrightarrow{\underline{Ad}} GL(\mathfrak g/\mathfrak r)=GL(n+1,\mathbb R)\xrightarrow{\rho(s)}GL_+(\mathbb R).\]

On the other hand, consider the representation of $R$ on
\begin{align*}
    \mathbb W=(\mathfrak g/\mathfrak r,\underline{Ad})\otimes(\mathbb R,\rho(-1)).
\end{align*}
To gain intuition, note that $\mathbb W\cong\mathbb R^{n+2}/\mathbb R$ as $R$-modules. Observe that $H$
is the isotropic subgroup of some nonzero base point whose $R$-orbit contains all nonzero elements in the module. Hence $R/H\cong\mathbb W\setminus\{0\}$
and so
\[\mathcal G/H=\mathcal G\times_RR/H=\mathcal G\times_R\mathbb W\setminus\{0\}\cong TN\otimes\mathcal E(-1)\setminus\{0\}\]
as fiber bundles over $N$.

Analogous to, e.g., \cite{Book}*{Theorem 4.4.3}, the representation of $R$ on $\mathfrak g/\mathfrak r$ gives rise to the identification
\[\mathcal G/Q=\mathcal G\times_RR/Q=\widehat{\mathcal PTN}\]
of the correspondence space. Moreover,
\[(\mathcal G\to\mathcal G/Q,\omega)\]
is the normal regular parabolic geometry canonically associated to the canonical lifting of a full path geometry $(\mathcal PTN,E,V)$ to $\widehat{\mathcal PTN}$. Its Fefferman-type space $\mathcal G/H$ admits the identification as desired.

It has been proven in \cref{bundles} that given any full path geometry $(\mathcal PTN,E,V)$, the Fefferman-type space of $(\widehat{\mathcal PTN},E,V)$ is isomorphic to the principal $\mathbb R_+$-bundle $GL_+((\mathcal H/V)^*\otimes\wedge^n\,V)$ over $\widehat{\mathcal PTN}$. We conclude from the above that $GL_+((\mathcal H/V)^*\otimes\wedge^n\,V)$, viewed as a fiber bundle over $N$, is isomorphic to  $TN\otimes\mathcal E(-1)\setminus\{0\}$. This proves the statement of the theorem.
\end{proof}
We need the following technical result to study the construction using Cartan geometry.
\begin{lemma}\label[lemma]{path}
Let $(\mathcal G\to M,\omega)$ be a normal regular parabolic geometry of type $(G,Q)$ that is canonically associated to a path geometry $(M,E,V)$ where $M$ is $(2\,n+1)$-dimensional and $E\oplus V$ is oriented. Consider the $Q$-submodules 
\[A=\left\{\left(\begin{array}{c|c|ccc}
		0&0& &*&\\
		\hline 0&0& &*&\\
		\hline  & & & &\\
		0&0&\ &*&\\
		&&&&
	\end{array}\right)\right\}\subseteq B=
\left\{\left(\begin{array}{c|c|ccc}
	0&*& &*&\\
	\hline 0&*& &*&\\
	\hline  & & & &\\
	0&*&\ &*&\\
	&&&&
\end{array}\right)\right\}\subseteq\mathfrak g=\mathfrak {sl}(n+2,\mathbb R)\]
where the block size is $(1,1,n)\times(1,1,n)$. The following properties hold for the curvature $\kappa$ of the Cartan connection $\omega$.
\begin{enumerate}
\item[(i)] $\kappa\in\Omega^2(M,\mathcal G\times_QB)$, $\kappa(V,TM)\subseteq\Gamma(\mathcal G\times_QA)$ and $\kappa(E\oplus V,E\oplus V)=0$.
\item[(ii)] If $(M,E,V)$ is torsion-free, then in addition to (i), we have that $\kappa\in\Omega^2(M,\mathcal G\times_QA)$ and $\kappa(E,TM)=0$.
\end{enumerate}
\end{lemma}
\begin{proof}
We use the improved Bianchi identity given in \cref{Bianchi}. Recall in \cref{convention} the $|2|$-grading decomposition of $\mathfrak g$ and from \cite{Book}*{Subsection 4.4.3} that the harmonic curvature 
\[\kappa_H:\mathcal G\to\ker\,\partial^*/\mathrm{im}\,\partial^*\cong\ker\,\square\subseteq\wedge^2\, \mathfrak q_+\otimes\mathfrak g\]
is given by $\kappa_H=(\tau,\rho)$ where
\begin{align*}
    \tau:\mathcal G\to(\mathfrak q_1^E\wedge\mathfrak q_2\otimes\mathfrak q_{-1}^V)\cap\ker\square\subseteq \mathfrak (\mathfrak q_1^E\oplus\mathfrak q_2)\wedge\mathfrak q_2\otimes B
\end{align*}
and
\begin{align*}
\rho:\mathcal G\to(\mathfrak q_1^V\wedge\mathfrak q_2\otimes\mathfrak q_{0})\cap\ker\square\subseteq \mathfrak q_1^V\wedge\mathfrak q_2\otimes A.
\end{align*}
Here, we have used that
\[(\mathfrak q_1^V\wedge\mathfrak q_2\otimes\mathfrak q_{0})\cap\ker\square\subseteq(\mathfrak q_1^V\wedge\mathfrak q_2\otimes\mathfrak q_{0})\cap\ker\partial^*\subseteq\mathfrak q_1^V\wedge\mathfrak q_2\otimes\mathfrak q^{ss}_{0},\]
where $\mathfrak q^{ss}_{0}=[\mathfrak q_0,\mathfrak q_0]\subseteq A$.

Consider the $Q$-module 
\[\mathbb F=(\mathfrak q_1^V\wedge\mathfrak q_2\otimes A)\oplus(\mathfrak q_1^E\wedge\mathfrak q_2\otimes B)
\oplus(\mathfrak q_2\wedge\mathfrak q_2\otimes B).\]
Then the images of $\kappa_H$ lie in $\mathbb F$. We observe that for $\varphi,\psi\in\mathbb F$, $\iota_\varphi\psi\in (\mathfrak q_1^E\oplus\mathfrak q_2)\wedge\mathfrak q_2\wedge\mathfrak q_2\otimes A$. Applying $\partial^*$ on both sides we get $\partial^*(\iota_\varphi\psi)\in(\mathfrak q_1^E\oplus\mathfrak q_2)\wedge\mathfrak q_2\otimes\mathfrak q_2\subseteq\mathbb F$. Hence $\mathbb F$ is stable under $\mathbb F$-insertions and so $\kappa$ is a section of $\Gamma(\mathcal G\times_Q\mathbb F)$, proving the claim in (i).

Now suppose that $(M,E,V)$ is torsion-free, i.e., $\kappa_H=\rho$. Then the images of $\kappa_H$ lie in the $Q$-module 
\[\mathbb E=(\mathfrak q_1^V\oplus\mathfrak q_2)\wedge\mathfrak q_2\otimes A.\] Observe that $\mathbb E$ is trivially stable under $\mathbb E$-insertions because for every $\varphi,\psi\in\mathbb E$, $\iota_\varphi\psi=0$. Hence $\kappa$ is a section of $\Gamma(\mathcal G\times_Q\mathbb E)$, proving the claim in (ii).
\end{proof}

\begin{proposition}\label{path Fef is normal}
Notation as at the beginning of the section. Then the Cartan connection $\tilde\omega$ is automatically normal. That is, the Fefferman-type construction from path geometries to almost Grassmann structures is normal.
\end{proposition}
\begin{proof}
Denote by $\partial^*$ and $\tilde\partial^*$ the Kostant codifferentials of $(\mathfrak g,\mathfrak q)$ and of $(\tilde{\mathfrak g},\tilde{\mathfrak p})$, respectively. Consider
\begin{align*}
    \phi\in\wedge^2\mathfrak q_+\otimes\mathfrak g
    \quad\text{and}\quad
    \tilde\phi\in\wedge^2\tilde{\mathfrak p}_+\otimes\tilde{\mathfrak g}
\end{align*}
such that
\[\tilde\phi=i'\circ\phi\circ\wedge^2\,\pi;\]
where $\pi$ is as at the beginning of the section.

Let $\alpha:\mathfrak g\to\tilde{\mathfrak g}$ be the embedding of $\mathfrak g=\mathfrak{sl}(n+2,\mathbb R)$ to the Lie subalgebra of $\tilde{\mathfrak g}=\mathfrak{sl}(n+3,\mathbb R)$ with zeros in the third row and the third column and $B\subseteq\mathfrak g$ be the $Q$-submodule defined in \cref{path}. We show that if $\phi\in\wedge^2\,\mathfrak q_+\otimes B$ and $\phi(\mathfrak q_{-1}^E,\mathfrak q_{-1}^V)=0$, then \[\tilde\partial^*\tilde\phi=\alpha\circ(\partial^*\phi)\circ\pi.\] Indeed, observe that the dual of $\pi:\tilde{\mathfrak g}/\tilde{\mathfrak p}\twoheadrightarrow\mathfrak g/\mathfrak q$ is $\pi^*:\mathfrak q_+\to\tilde{\mathfrak p}_+$ where, as matrices,
\begin{align}\label{dual pi}
\pi^*(Z)_{ij}=Z_{i,j+1},\quad Z\in\mathfrak q_+.
\end{align}
That is, we have $\langle Z,\pi( \tilde X)\rangle=\langle \pi^*( Z),\tilde X\rangle$ for all $Z\in\mathfrak q_+$ and $\tilde X\in\tilde{\mathfrak g}/\tilde{\mathfrak p}$. Observe that for all $Z\in \mathfrak q_+$ and $W\in B$, there holds
\begin{align}\label{pi and alpha}
[\pi^*(Z),i'(W)]=\alpha([Z,W]).
\end{align}
Moreover, let $X\in\mathfrak g$ be any element and $\{X^i\}$ a basis of $\mathfrak g/\mathfrak q$ with dual basis $\{Z_i\}$ of $\mathfrak q_+$. Then
\begin{align*}
    (\partial^*\phi)(X)=2\sum_i[Z_i,\phi(X,X^i)]-\sum_i\phi([Z_i,X],X^i)=2\sum_i[Z_i,\phi(X,X^i)];
\end{align*}
see \eqref{codiff}. Note that the assumption $\phi(\mathfrak q_{-1}^E,\mathfrak q_{-1}^V)=0$ implies that $\phi([Z_i,X],X^i)=0$ for all pairs $(X^i,Z_i)$.

Let $(\{\tilde X^i\},\tilde Y)$ be a basis of $\tilde{\mathfrak g}/\tilde{\mathfrak p}$ such that $\pi(\tilde X^i)=X^i$ for all $i$ and $\tilde Y\in\ker\,\pi$. Then its dual basis is $(\{\pi^*Z_i\},\tilde W)\in\tilde{\mathfrak p}_+$ where $\tilde W$ is uniquely determined. For any $\tilde X\in\tilde{\mathfrak g}/\tilde{\mathfrak p}$, write $X=\pi(\tilde X)$. Then, using \eqref{codiff}, the assumption $\pi(\tilde Y)=0$ and \eqref{pi and alpha},
\begin{align*}
   (\tilde\partial^*\tilde\phi) (\tilde X)&=2\left([\tilde W,i'\circ\phi(X,0)]+
   \sum_i[\pi^*(Z_i),i'\circ\phi(X,X^i)]\right)\\
   &=2\sum_i\alpha([Z_i,\phi(X,X^i)])\\
   &=(\partial^*\phi)(X).
\end{align*}
This proves the claim.

Using \cref{path} (i) and the above claim, we conclude that
\[(\tilde\partial^*\tilde\kappa)(j(u))=\alpha\circ(\partial^*\kappa)(u)\circ\pi=0\]
for every $u\in\mathcal G$; where $j,\kappa$ and $\tilde\kappa$ are as at the beginning of the section. This completes the proof of the proposition.
\end{proof}
Consequently, the natural map
\[j:\mathcal G\to\tilde{\mathcal G}=\mathcal G\times_H\tilde P\]
defining $\tilde{\mathcal G}$ relates Cartan objects of $(M,E,V)$ to those of $(\tilde M,\tilde E,\tilde F)$ in a canonical way. For example, the two equalities at the beginning of \cref{section abstract Fefferman} are relations between the normal regular Cartan connection $\omega$ of the initial path geometry and the normal Cartan connection $\tilde\omega$ of the resulting almost Grassmann structure and between the curvature $\kappa$ of $\omega$ and the curvature $\tilde\kappa$ of $\tilde\omega$, respectively. A tractor bundle on $(\tilde M,\tilde E,\tilde F)$ associated to a representation of $\tilde G$ is the pullback of the tractor bundle on $(M,E,V)$ associated to the representation induced by $i:G\to\tilde G$. Moreover, recall that the canonical curves on $(M,E,V)$ are projections of the flows $Fl^{\omega^{-1}(X)}(u)$ along $\mathcal G\to M$ where $X\in\mathfrak g$ and $u\in\mathcal G$. Observe that 
\begin{align*}
j\circ Fl\,^{\omega^{-1}(X)}(u)
=Fl\,^{{\tilde\omega}^{-1}(i(X))}
(j(u)).
\end{align*} 
It follows that every canonical curve of $(M,E,V)$ is lifted along $\tilde M\twoheadrightarrow M$ to a canonical curve of $(\tilde M,\tilde E,\tilde F)$.

Additionally, the following properties hold for the  Fefferman-type construction we specified. Notation as at the beginning of the paragraph. Then
\begin{align*}
\tilde{\mathcal T}=\tilde{\mathcal G}\times_{\tilde{P}}\mathbb R^{n+3}
\quad\text{and}\quad
\tilde{\mathcal G}\times_{\tilde{P}}\tilde{\mathfrak g}=\mathfrak{sl}(\tilde{\mathcal T})
\end{align*}
are the standard tractor bundle and the adjoint tractor bundle of $(\tilde M,\tilde E,\tilde F)$, respectively. Then
\begin{align*}
    \tilde\kappa\in\Omega^2(\tilde M,\mathfrak{sl}(\tilde{\mathcal T})).
\end{align*}
Recall the natural short exact sequences
\begin{align*}
    0\to\tilde E\to\tilde{\mathcal T}\to\tilde F\to 0\quad\text{and}\quad 0\to \tilde F^*\to\tilde{\mathcal T}^*\to\tilde E^*\to 0.
\end{align*}
They follow easily from \cref{AG parabolic}. Consider the nowhere vanishing sections
\begin{align*}
\varphi=\mathcal G\times_H(0,1)\in\Gamma(\tilde E^*)
\quad\text{and}\quad
\eta=\mathcal G\times_H(-1,0^n)^t\in\Gamma(\tilde F)
\end{align*}
defined in \cref{distinguished sections}. Moreover, denote the torsion of $(\tilde M,\tilde E,\tilde F)$ by $\tilde\tau$. The contraction of $\tilde\tau\otimes\tilde\tau$ specified in \cref{Weyl tensor} results in \[tr(\iota_{\tilde\tau}\tilde\tau)\in\Gamma(\mathrm{Sym}^2\,T^*\tilde M).\]

\begin{theorem}\label[theorem]{properties of Fef on path}
Consider the Fefferman-type construction
\[(M,E,V)\mapsto(\tilde M,\tilde E,\tilde F)\]
of almost Grassmann structures from path geometries. The following properties hold.
\begin{itemize}
\item [(i)]
$tr(\iota_{\tilde\tau}\tilde\tau)=0$.
\item [(ii)] $(\tilde M,\tilde E,\tilde F)$ is torsion-free if and only if $(M,E,V)$ is torsion-free.
\item [(iii)] There are parallel sections $\mu\in\Gamma(\tilde{\mathcal T})$ and $\nu\in\Gamma(\tilde{\mathcal T}^*)$ with projections $\eta\in\Gamma(\tilde F)$ and $\varphi\in\Gamma(\tilde E^*)$, respectively. Moreover, there holds $\mu(\nu)=1$.
\item [(iv)] Denote by \[\ker\,\varphi\subseteq\tilde E\subseteq\mathrm{span}(\eta)+\tilde E\subseteq\tilde{\mathcal T}\] the natural filtration of ranks $1\leq 2\leq 3\leq n+3$ and $\tilde{\mathcal D}=\mathrm{span}(\varphi\otimes\tilde F,
    \tilde E^*\otimes\eta)\subseteq T\tilde M$. Then $\tilde\kappa$ has the following properties.
\begin{equation*}
\begin{aligned}
\begin{cases}
\tilde\kappa(\varphi\otimes\eta,T\tilde M)&=0\\
\tilde\kappa(\tilde{\mathcal D},\tilde{\mathcal D})&=0\\
\tilde\kappa(\varphi\otimes\tilde F,T\tilde M)&\subseteq\{\phi\in\Gamma(\mathfrak{sl}(\tilde{\mathcal T})):\phi(\mathrm{span}(\eta)+\tilde E)=0\}\quad\text{and}\\
\tilde\kappa(T\tilde M,T\tilde M)&\subseteq\{\phi\in\Gamma(\mathfrak{sl}(\tilde{\mathcal T})):\phi(\ker\,\varphi)=0\}.
\end{cases}
\end{aligned}
\end{equation*}
If $(M,E,V)$ is torsion-free, then in addition to the above we have that
\begin{equation*}
\begin{aligned}
\begin{cases}
\tilde\kappa(\tilde E^*\otimes\eta,T\tilde M)&=0
\quad\text{and}\\
\tilde\kappa(T\tilde M,T\tilde M)&\subseteq\{\phi\in\Gamma(\mathfrak{sl}(\tilde{\mathcal T})):\phi(\mathrm{span}(\eta)+\tilde E)=0\}.
\end{cases}
\end{aligned}
\end{equation*}
\end{itemize} 
\end{theorem}
\begin{proof}
Notation as in the first paragraph of the proof of \cref{path Fef is normal}. Recall that $\tilde\kappa$ is completely determined by the curvature $\kappa$ of $(\mathcal G\to M,\omega)$ via the relation
\[\tilde\kappa(j(u))=
i'\circ\kappa(u)\circ\wedge^2\,\pi\]
for all $u\in\mathcal G$. Thus the properties of $\kappa$ given in \cref{path} (i) give rise to the properties of $\tilde\kappa$ in (iv) in the statement of the theorem. Recall from \cref{distinguished sections} that $\varphi\otimes\eta$ spans the vertical bundle of $p:\tilde M\to M$.

Let $A\subseteq\mathfrak g$ be the $Q$-submodule defined in \cref{path}. That $i'^{-1}(\tilde{\mathfrak p})=\mathfrak h\subseteq\mathfrak q$ immediately implies that if $\tilde\phi\in\wedge^2\,\tilde{\mathfrak p}_+\otimes\tilde{\mathfrak p}$, then $\phi\in\wedge^2\,\mathfrak q_+\otimes \mathfrak q$. Moreover, that $i'(A)\subseteq\tilde{\mathfrak p}$ immediately implies that if  $\phi\in\wedge^2\,\mathfrak q_+\otimes A$, then $\tilde\phi\in\wedge^2\,\tilde{\mathfrak p}_+\otimes\tilde{\mathfrak p}$. Using \cref{path} and these two relations, we conclude that (ii) holds.

Observe that \cref{Hol} applies to our Fefferman-type functor from Cartan geometries of type $(G,Q)$ to Cartan geometries of type $(\tilde G,\tilde P)$ defined by $i:G\to\tilde G$ where we choose the tractor bundle $\tilde{\mathcal W}=\tilde{\mathcal T}\oplus\tilde{\mathcal T}^*$ and the fiber subbundle $\tilde{\mathcal O}=\tilde{\mathcal G}\times_{\tilde P}(\tilde P\,\tilde w_0)$ determined by the point
\[\tilde w_0=\left((0,0,-1,0^n)^t,(0,1,-1,0^n)\right)\in\mathbb R^{n+3}\oplus\mathbb R^{(n+3)*}\] fixed by $H$. As explained in the proof of \cref{Hol}, the constant map $\mathcal G\to \{\tilde w_0\}$ uniquely extends to a pair
\begin{align*}
(\mu,\nu)\in     
C^\infty(\tilde{\mathcal G},\mathbb R^{n+3}\oplus\mathbb R^{(n+3)*})^{\tilde P}
=\Gamma(\tilde{\mathcal T}\oplus\tilde{\mathcal T}^*)
\end{align*}
of parallel sections. Clearly, $(\mu,\nu)$ projects to $(\eta,\varphi)$ and $\mu(\nu)=1$. This proves (iii).

It remains to prove (i). For every $\tilde x\in\tilde M$, we may choose $u\in\mathcal G$ such that $j(u)\in\tilde{\mathcal G}$ projects to $\tilde x$. Let $\tilde\sigma:\tilde{\mathcal G}/\tilde P_+\to\tilde{\mathcal G}$ be a Weyl structure such that $j(u)$ lies in its image. Consider the component
\[\tilde W
\in\Omega^2(\tilde M,L(\tilde E,\tilde E))\]
of the Weyl tensor of $\tilde\sigma$; see \eqref{induced Weyl objects}. It follows from the fourth property of $\tilde\kappa$ in (iv) that
\[\tilde W_{\tilde x}(T_{\tilde x}\tilde M,T_{\tilde x}\tilde M)\subseteq \varphi(\tilde x)\otimes\tilde E_{\tilde x}\]
and from the third property of $\tilde\kappa$ in (iv) that 
\[\tilde W_{\tilde x}(\varphi(\tilde x)\otimes\tilde F_{\tilde x},T_{\tilde x}\tilde M)=0.\]
Recall that in \cref{Weyl tensor} we fix a Ricci-type contraction $tr(\tilde W)$ of $\tilde W$ and provide a relation between $tr(\tilde W)$ and $tr(\iota_{\tilde\tau}\tilde\tau)(\tilde x)$, which implies that $tr(\iota_{\tilde\tau}\tilde\tau)(\tilde x)=0$ if and only if $tr(\tilde W_{\tilde x})=0$. In fact, the above two properties on $\tilde W_{\tilde x}$ together imply that $tr(\tilde W_{\tilde x})=0$. This proves (i).
\end{proof}
From (iv) in \cref{properties of Fef on path} we conclude that
\begin{align}
\tilde\tau\in\Gamma\big(L\big(\wedge^2(T\tilde M/(\varphi\otimes\tilde F)),\varphi\otimes\tilde F\big)\big).  
\end{align}
Observe that $\tilde\tau$ is completely determined by the projection of the regular normal Cartan curvature $\kappa\in\Omega^2(M,\mathcal G\times_Q\mathfrak g)$ of the path geometry $(M,E,V)$ in $\Omega^2(M,\mathcal G\times_Q\mathfrak g/\mathfrak h)$; see \cref{convention}. If $(M,E,V)$ is torsion-free, then $\tilde\tau=0$ by \cref{properties of Fef on path} (ii) and the harmonic curvature $\tilde\rho\in\Omega^2(\tilde M,\mathfrak{sl}(\tilde F))$ of $(\tilde M,\tilde E,\tilde F)$ equals the projection of $\tilde\kappa\in\Omega^2(\tilde M,\tilde{\mathcal G}\times_{\tilde P}\tilde{\mathfrak p})$ to $\Omega^2(\tilde M,\tilde{\mathcal G}\times_{\tilde P}\tilde{\mathfrak p}/\tilde{\mathfrak p}_+)$; see \cref{convention} and, e.g., \cite{Book}*{Theorem 3.1.12}. In particular, 
\begin{align}
\tilde\rho\in\Gamma\big(L\big(\wedge^2(T\tilde M/(\tilde E^*\otimes\eta)),\mathfrak{sl}(\tilde F)\big)\big)
\end{align}
and is completely determined by the projection of $\kappa\in\Omega^2(M,\mathcal G\times_Q\mathfrak q)$ to $\Omega^2(M,\mathcal G\times_Q\mathfrak q/\mathfrak q_2)$.

Combining (i) and (ii) in \cref{properties of Fef on path}, we see that there are plenty of almost Grassmann structures with non-vanishing torsion $\tau$ such that $tr(\iota_\tau\tau)=0$. Thus the property $tr(\iota_\tau\tau)=0$ is a nontrivial weakening of torsion-freeness.

The tensor field $tr(\iota_\tau\tau)$ turns out to be important to almost Grassmann structures: We prove in \cref{AG Fef normality} below that a Fefferman-type construction from an almost Grassmann structure of type $(2,n)$ to an almost Grassmann structure of type $(2,n+1)$ is normal if and only if the initial almost Grassmann structure of type $(2,n)$ has the property that $tr(\iota_\tau\tau)=0$. It is proven in \cite{Guo}*{Theorem 8 and Proposition 11} that parallel standard tractors (respectively, parallel standard cotractors) of an almost Grassmann structure $(M,E,F)$ are in one-to-one correspondence with sections of $F$ (respectively, of $E^*$) which are parallel for some Weyl connection and annihilate $tr(\iota_\tau\tau)$.

Conversely, we want to characterize all almost Grassmann structures $(\tilde M,\tilde E,\tilde F)$ that are locally the result of the Fefferman-type construction. That is, there is an open cover of $\tilde M$ such that the restriction of $(\tilde M,\tilde E,\tilde F)$ to every open set in the cover is isomorphic to an open subspace of an almost Grassmann structure resulted from the Fefferman-type construction. Here, the restriction of the almost Grassmann structure
\[(\tilde M,\tilde E,\tilde F)=(\tilde M,p:\tilde E\to\tilde M,q:\tilde F\to\tilde M)\]
to an open subset $\tilde U\subseteq\tilde M$ refers to the almost Grassmann structure $(\tilde U,p^{-1}(\tilde U),q^{-1}(\tilde U))$, which is said to be an open subspace of $(\tilde M,\tilde E,\tilde F)$.

Let $(\tilde M,\tilde E,\tilde F)$ be an almost Grassmann structure of type $(2,n+1)$ with canonically associated normal parabolic geometry $(\tilde{\mathcal G}\to\tilde M,\tilde\omega)$ of type $(\tilde G,\tilde P)$. Denote by $\tilde \kappa\in\Omega^2(\tilde M,\tilde{\mathcal G}\times_{\tilde P}\tilde{\mathfrak g})$ the curvature of $\tilde\omega$ and by $\tilde{\mathcal T}=\tilde{\mathcal G}\times_{\tilde P}\mathbb R^{n+3}$ the standard tractor bundle of $(\tilde M,\tilde E,\tilde F)$. Then $\tilde{\mathcal T}^*$ is the standard cotractor bundle of $(\tilde M,\tilde E,\tilde F)$. Recall the natural projections $\tilde{\mathcal T}\twoheadrightarrow\tilde F$ and $\tilde{\mathcal T}^*\twoheadrightarrow\tilde E^*$. They may be obtained using \cref{AG parabolic} in a straightforward way.

\begin{theorem}\label[theorem]{charpath}
Notation as in the preceding paragraph.
$(\tilde M,\tilde E,\tilde F)$ is locally the result of the Fefferman-type construction on a path geometry if and only if the following hold.
\begin{enumerate}
\item [(i)] There are parallel sections $\mu\in\Gamma(\tilde{\mathcal T})$ and $\nu\in\Gamma(\tilde{\mathcal T}^*)$ with $\mu(\nu)=1$. Moreover, the projection $\eta\in\Gamma(\tilde F)$ of $\mu$ and $\varphi\in\Gamma(\tilde E^*)$ of $\nu$ are both nowhere vanishing.
    \item[(ii)] $\tilde\kappa(\varphi\otimes\eta,
    \tilde E^*\otimes\tilde F)=0$,  $\tilde\kappa(\varphi\otimes\tilde F,
    \tilde E^*\otimes\eta)\subseteq\tilde{\mathcal G}\times_{\tilde P}\tilde{\mathfrak p}_+$ and 
    $\varphi\otimes\tilde F\subseteq T\tilde M$ is an involutive distribution.
\end{enumerate}
Moreover, if $(\tilde M,\tilde E,\tilde F)$ is torsion-free, then (i) implies (ii) already.
\end{theorem}
\begin{proof}
That all almost Grassmann structures resulting from our Fefferman-type construction satisfy (i) and (ii) has been proven in \cref{properties of Fef on path}. We prove that the converse direction holds.

As mentioned in the proof of \cref{properties of Fef on path}, \cref{Hol} applies to our Fefferman-type functor from Cartan geometries of type $(G,Q)$ to Cartan geometries of type $(\tilde G,\tilde P)$ defined by $i:G\to\tilde G$ where we choose the tractor bundle $\tilde{\mathcal W}=\tilde{\mathcal T}\oplus\tilde{\mathcal T}^*$ and the fiber subbundle $\tilde{\mathcal O}=\tilde{\mathcal G}\times_{\tilde P}(\tilde P\,\tilde w_0)$ determined by the point
\[\tilde w_0=\left((0,0,-1,0^n)^t,(0,1,-1,0^n)\right)\in\mathbb R^{n+3}\oplus\mathbb R^{(n+3)*}\] fixed by $H$. One could easily check that the sections of $\tilde{\mathcal O}$ are precisely pairs of sections $\mu'\in\Gamma(\tilde{\mathcal T})$ and $\nu'\in\Gamma(\tilde{\mathcal T}^*)$ such that $\nu'(\mu')=1$ and that the projections $\eta'\in\Gamma(\tilde F)$ and $\varphi'\in\Gamma(\tilde E^*)$ are both nowhere vanishing. Moreover, the distribution $\tilde{\mathcal G}\times_{\tilde P}(i'(\mathfrak q)+\tilde{\mathfrak p})/\tilde{\mathfrak p}$ is spanned by $\varphi'\otimes\eta'$.

Assume that (i) and (ii) hold. It follows from \cref{Hol} that $(\tilde{\mathcal G}\to\tilde M,\tilde\omega)$ is locally the result of the Fefferman-type functor applied to a parabolic geometry of type $(G,Q)$. Without loss of generality, assume that $(\tilde{\mathcal G}\to\tilde M,\tilde\omega)$ is the result of the Fefferman-type functor applied to a parabolic geometry $(\mathcal G\to M,\omega)$ of type $(G,Q)$. Consider the distinguished distributions 
\begin{align*}
E=\mathcal G\times_Q(\mathfrak q_{-1}^E\oplus\mathfrak q)/\mathfrak q
\quad\text{and}\quad
V=\mathcal G\times_Q(\mathfrak q_{-1}^V\oplus\mathfrak q)/\mathfrak q
\end{align*}
on $M$. Denote by $p:\tilde M\to M$ the natural projection. All relations in \cref{bundles} and \cref{distinguished sections} hold. In particular, we have that
\begin{align*}
(Tp)^{-1}(E)=\tilde E^*\otimes\eta
\quad\text{and}\quad
(Tp)^{-1}(V)=\varphi\otimes\tilde F.
\end{align*}
Thus the assumptions on $\tilde\kappa$ in (ii) imply that $(\mathcal G\to M,\omega)$ is a regular parabolic geometry with underlying path geometry $(M,E,V)$ with oriented $E\oplus V$.

Note that the characterizing properties in (i) and (ii) of the theorem is a subcollection of the properties listed in \cref{properties of Fef on path}. This subcollection of conditions is sufficient to prove that $(\tilde M,\tilde E,\tilde F)$ is the result of the Fefferman-type construction on $(M,E,V)$ but is insufficient to prove that $\omega$ is normal. In fact, let $\omega^{nor}\in\Omega^1(\mathcal G,\mathfrak g)$
be the normal regular Cartan connection of type $(G,Q)$ canonically associated to the geometry $(M,E,V)$; see, e.g., \cite{Book}*{Proposition 3.1.13}. Recall that $\omega^{nor}$ is unique up to pullback by an automorphism on the principal bundle $\mathcal G\to M$; see, e.g., \cite{Book}*{Proposition 3.1.14}. We show that
\begin{enumerate}
    \item [$\bullet$] there holds $\omega^{nor}-\omega\in\Omega^1(\mathcal G,\mathfrak q^1)$ where $\mathfrak q^1\subseteq\mathfrak h=i'^{-1}(\tilde{\mathfrak p})$; and 
    \item [$\bullet$] if in addition $\tilde\kappa(\varphi\otimes\tilde F,\tilde E^*\otimes\eta)=0$, we may choose $\omega^{nor}=\omega$.
\end{enumerate}
Note that the assumption $\tilde\kappa(\varphi\otimes\tilde F,\tilde E^*\otimes\eta)=0$ in the second claim above is also taken from the properties listed in \cref{properties of Fef on path}. Indeed, we use the following notation from the proof of \cref{path Fef is normal}. Let $\partial^*$ and $\tilde{\partial}^*$ be the Kostant differentials of $(\mathfrak g,\mathfrak q)$ and $(\tilde{\mathfrak g},\tilde{\mathfrak p})$, respectively. Let $\pi^*:\mathfrak q_+\to\tilde{\mathfrak p}_+$ be the dual map of $\pi:\tilde{\mathfrak g}/\tilde{\mathfrak p}\twoheadrightarrow\mathfrak g/\mathfrak q$, $\{X^i\}\subseteq\mathfrak g/\mathfrak q$ any basis with dual basis $\{Z_i\}\subseteq\mathfrak q_+$, $(\{\tilde X^i\},\tilde Y)\subseteq\tilde{\mathfrak g}/\tilde{\mathfrak p}$ a basis such that $\pi(\tilde X_i)=X_i$ for all $i$ and $0\neq\tilde Y\subseteq\ker\,\pi$ with dual basis $(\{\pi^*(Z_i)\},\tilde W)\subseteq\tilde{\mathfrak p}_+$. Moreover, let 
\begin{align*}
\beta:\tilde{\mathfrak g}\twoheadrightarrow\mathfrak g
\end{align*}
be the linear map given by deleting the third row and summing up the second and the third column of $\tilde{\mathfrak g}=\mathfrak {sl}(n+3,\mathbb R)$. Then we have the relation
\[\beta([\pi^*(Z),i'(W)])=[Z,W]\]
for all $Z\in\mathfrak q_+$ and $W\in\mathfrak g$. Let $\phi\in\wedge^2\,\mathfrak q_+\otimes\mathfrak g$ and 
\[\tilde\phi
=i'\circ\phi\circ\wedge^2\,\pi\in\wedge^2\,\tilde{\mathfrak p}_+\otimes\tilde{\mathfrak g}.\]
Recall that
\begin{align*}
    (\partial^*\phi)(X)
    =2\sum_i[Z_i,\phi(X,X^i)]-\sum_i\phi([Z_i,X],X^i).
\end{align*}
for all $X\in\mathfrak g$; see \eqref{codiff}. Let $\tilde X=i'(X)\in \tilde{\mathfrak g}$. Since $\pi(\tilde X+\tilde{\mathfrak p})=X+\mathfrak q$ and the pair $(\tilde W,\tilde Y)$ of bases satisfies $[\tilde W,\tilde\phi(\tilde X,\tilde Y)]=[\tilde W,i'\circ\phi(X,0)]=0$, we conclude that
\begin{align}\label{Kostant relation for beta}
    2\sum_i[Z_i,\phi(X,X^i)]=2\sum_i \beta([\pi^*(Z_i),\tilde\phi(\tilde X,\tilde X^i)])
    &=\beta\circ(\tilde\partial^*\tilde\phi)(\tilde X).
\end{align}
On the other hand, we show that
\begin{equation}\label{nontrivial summand}
-\sum_i\phi([Z_i,X],X^i)=
\begin{aligned}   
\begin{dcases}   
0& X\in\mathfrak q^{-1}\\   
2\,\phi(X_E,X_V)& X=[X_E,X_V],\quad X_E\in\mathfrak q_{-1}^E,X_V\in\mathfrak q_{-1}^V.
\end{dcases}
\end{aligned}  
\end{equation}
Indeed, we may assume without loss of generality that
\begin{align*}    Z_1\in\mathfrak q_1^E,\quad Z_{1+j}\in\mathfrak q_1^V\quad\text{and}\quad Z_{n+1+j}=[Z_1,Z_{1+j}]\in\mathfrak q_2\end{align*} 
for all $1\leq j\leq n$. Then
\begin{align*}
    X^{n+1+j}=-[X^1,X^{1+j}].
\end{align*}
If $X\in\mathfrak q^{-1}$, then $\phi([Z_i,X],X^i)\subseteq\phi(\mathfrak q,X^i)=0$ for all $1\leq i\leq 2n+1$; and if 
\begin{align*}    X+\mathfrak q=X^{n+1+j_0}+\mathfrak q=-[X^1,X^{1+j_0}]+\mathfrak q
\end{align*}
for some $1\leq j_0\leq n$, then
\begin{align*}    
-\sum_i\phi([Z_i,[X^1,X^{1+j_0}]],X^i)&=\sum_i\phi([Z_i,X],X^i)  \\  &=\phi([Z_1,X],X^1)+\phi([Z_{1+j_0},X],X^{1+j_0})\\    &=\phi(-X^{1+j},X^1)+\phi(X^1,X^{1+j_0})\\    
&=2\,\phi(X^1,X^{1+j_0}).
\end{align*}
This completes the proof of \eqref{nontrivial summand}.

Now substitute $\phi=\kappa(u)$ and $\tilde\phi=\tilde\kappa(j(u))$ in the above where $\kappa$ is the curvature of $\omega$, $\tilde\kappa$ is the curvature of $\tilde\omega$, $j:\mathcal G\to\tilde{\mathcal G}$ the inclusion map and $u\in\mathcal G$ is any point. Since $\tilde\kappa(\tilde E^*\otimes\eta,\varphi\otimes\tilde F)\subseteq\tilde{\mathcal G}\times_{\tilde P}\tilde{\mathfrak{\mathfrak p}}_+$ and $i'^{-1}(\tilde{\mathfrak p}_+)=\mathfrak q_{1}^V\oplus\mathfrak q_2\subseteq\mathfrak q^1$, we conclude that
\begin{align*}
(\partial^*\kappa)(E\oplus V)=0
\quad\text{and}\quad
    (\partial^*\kappa)\circ\mathcal L|_{E\times V}=2\kappa|_{E\times V}\subseteq\mathcal G\times_Q\mathfrak q^1
\end{align*}
where $\mathcal L:\wedge^2(E\oplus V)\to\frac{TM}{E\oplus V}$ is the Levi bracket. In particular, $\partial^*\kappa$ has homogeneity $\geq 3$. It follows from, e.g., \cite{Book}*{Proposition 3.1.13} that $\omega^{nor}-\omega$ also has homogeneity $\geq 3$ and thus lies in $\Omega^1(\mathcal G,\mathfrak q^1)$. If in addition $\tilde\kappa(\tilde E^*\otimes\eta,\varphi\otimes\tilde F)=0$, the above procedure implies that $\partial^*\kappa=0$ and so $(\mathcal G\to M,\omega)$ is already the normal regular parabolic geometry canonically associated to $(M,E,V)$. This completes the proof of the above two claims.

Since $i'(\mathfrak q^1)\subseteq\tilde{\mathfrak p}$, the Fefferman-type functor applied to $(\mathcal G\to M,\omega^{nor})$ results in a parabolic geometry $(\tilde{\mathcal G}\to\tilde M,\widetilde{\omega^{nor}})$ such that $\widetilde{\omega^{nor}}-\tilde\omega\in\Omega^1(\tilde{\mathcal G},\tilde{\mathfrak p})$. In particular, the underlying geometry of $(\tilde{\mathcal G}\to\tilde M,\widetilde{\omega^{nor}})$ is the same as the underlying geometry of $(\tilde{\mathcal G}\to\tilde M,\tilde\omega)$; see, e.g., \cite{Book}*{Proposition 3.1.10 (1)}. Hence $(\tilde M,\tilde E,\tilde F)$ is the result of the Fefferman-type construction on $(M,E,V)$. This proves the converse direction.

Finally, if $(\tilde M,\tilde E,\tilde F)$ is torsion-free and satisfies (i), then $\tilde\kappa(\varphi\otimes\tilde F,\tilde E^*\otimes\eta)$ and $\tilde\kappa(\varphi\otimes\tilde F,\varphi\otimes\tilde F)$ take values in $\tilde{\mathcal G}\times_{\tilde P}\tilde{\mathfrak p}_+$. In particular, $\varphi\otimes\tilde F$ is an involutive distribution. Moreover, observe that the trace-free part of $\mu\otimes\nu\in\mathfrak{gl}(\tilde{\mathcal T})$ is a parallel adjoint tractor. By \cite{Cap08}*{3.5 Corollary}, there holds $\tilde\kappa(\varphi\otimes\eta,\tilde E^*\otimes\tilde F)=0$ and so we obtain (ii). This completes the prove of the theorem.
\end{proof}

Thus we have provided a full characterization of the Fefferman-type spaces we have introduced in a style analogous to the classical result in \cite{CG06}*{Theorem 2.5} characterizing conformal structures that are locally constructed from CR structures by the existence of a parallel adjoint tractor. Recall from, e.g., \cite{Book}*{Subsection 4.4.5} that torsion-free path geometries are locally the correspondence space of a Grassmann structure. Hence given a Grassmann structure $(\tilde M,\tilde E,\tilde F)$ of type $(2,n+1)$, there holds \cref{charpath} (i) if and only if $(\tilde M,\tilde E,\tilde F)$ arises locally as the Fefferman-type space of a Grassmann structure of type $(2,n)$; see \cref{properties of Fef on path} (ii).

Recall from \cref{AG connections} the explicit description of Weyl connections of an almost Grassmann structure and the one-to-one correspondence between Weyl structures and Weyl connections of an almost Grassmann structure. In particular, every Weyl connection is canonically associated to a collection of Weyl objects listed in  \eqref{induced Weyl objects}. Now let $(M,E,F)$ be an almost Grassmann structure of type $(2,n)$, $n\geq 3$. Recall the convention of Penrose abstract index notation below \eqref{induced Weyl objects}. Given a Weyl connection $\nabla$ with canonically associated Weyl structure $\sigma$, denote by $R(\nabla^{TM}),R(\nabla^E)$ and $R(\nabla^F)$ the curvatures of $\nabla$ on $TM,E$ and $F$, respectively. Denote by $\tau$ the harmonic torsion and $\rho$ the harmonic curvature of $(M,E,F)$. Recall that $\tau$ agrees with the torsion of $\nabla$ on $TM$, see \cref{AG connections}, and
\begin{align*}
\rho{}^A_{A'}{}^B_{B'}{}^{C'}_{D'}=R(\nabla^F){}_{(}{}^A_{A'}{}^B_{B'}{}^{C'}_{D'}{}_{)}-\frac{1}{n+2}\,R(\nabla^F){}_{(}{}^A_{A'}{}^B_{B'}{}_{|}{}^{I'}_{I'}{}_{|}\delta{}^{C'}_{D'}{}_{)}
-\frac{2}{n+2}\,R(\nabla^F){}^{[}{}^{A}_{I'}{}_{(}{}^{B}_{A'}{}^{]}
{}^{I'}_{B'}\delta{}^{C'}_{D'}{}_{)};
\end{align*}
see, e.g., the end of Step (E) in \cite{Book}*{Subsection 4.1.3}. Moreover, denote by $Ric(\nabla)$ the Ricci curvature of $\nabla$ and denote by $\mathrm{P}$ and $(W,W')$ the Rho tensor and the Weyl tensor of $\sigma$, respectively. We prove the following relations.
\begin{lemma}\label[lemma]{Weyl tensor Rho Ricci}
Notation as in the preceding paragraph. Then there holds
\begin{equation}\label{Weyl tensor expanded}
\begin{aligned}
\begin{dcases}
W{}^A_{A'}{}^B_{B'}{}^{C}_{D}&=R(\nabla^E)^A_{A'}{}^B_{B'}{}^{C}_{D}+\delta^{B}_{D}\,\mathrm P{}^{A}_{A'}{}^{C}_{B'}-\delta^{A}_{D}\,\mathrm P{}^{B}_{B'}{}^{C}_{A'}\\
W'{}^A_{A'}{}^B_{B'}{}^{C'}_{D'}&=R(\nabla^F)^A_{A'}{}^B_{B'}{}^{C'}_{D'}-\delta^{C'}_{B'}\,\mathrm P{}^{A}_{A'}{}^{B}_{D'} +\delta^{C'}_{A'}\,\mathrm P{}^{B}_{B'}{}^{A}_{D'}   
\end{dcases}
\end{aligned}
\end{equation}
where
\begin{align}
    \mathrm P{}^{A}_{A'}{}^{B}_{B'}&=\frac 1 n\,Ric(\nabla)^{(A}_{(A'}{}^{B)}_{B')}+\frac{1}{n+4}\,Ric(\nabla)^{[A}_{[A'}{}^{B]}_{B']}\nonumber\\
    &+\frac{1}{n+2}\,Ric(\nabla)^{(A}_{[A'}{}^{B)}_{B']}+\frac{1}{n+2}\,Ric(\nabla)^{[A}_{(A'}{}^{B]}_{B')}.\label{Rho Ric}
\end{align}
\end{lemma}
\begin{proof}
Denote by $R(\nabla^{TM}),R(\nabla^E)$ and $R(\nabla^F)$ the curvatures of $\nabla$ on $TM,E$ and $F$, respectively. Since $TM\cong E^*\otimes F$, there holds
\begin{align*}
        R(\nabla^{TM})^A_{A'}{}^B_{B'}{}^{C'}_{C}{}^D_{D'}&=\delta^D_C\, R(\nabla^F)^A_{A'}{}^B_{B'}{}^{C'}_{D'}-\delta^{C'}_{D'}\,R(\nabla^E)^A_{A'}{}^B_{B'}{}^{D}_{C}
    \end{align*}
so that
\begin{align*}
Ric(\nabla)^A_{A'}{}^D_{D'}&=R(\nabla^{TM})^A_{A'}{}^I_{I'}{}^{I'}_{I}{}^D_{D'}\nonumber\\
&=\delta^D_I\,R(\nabla^F)^A_{A'}{}^I_{I'}{}^{I'}_{D'}-\delta^{I'}_{D'}\,R(\nabla^E)^A_{A'}{}^I_{I'}{}^{D}_{I}\nonumber\\
    &=R(\nabla^F)^A_{A'}{}^D_{I'}{}^{I'}_{D'}-R(\nabla^E)^A_{A'}{}^I_{D'}{}^{D}_{I}.\label{Ric}
\end{align*}



Denote by $(W,W')$ the Weyl tensor of $\sigma$ and by $\partial$ the Lie algebra homology differential; see \eqref{partial}. Recall from, e.g., \cite{Book}*{Theorem 5.2.3 (3)} that 
\[(W,W')=(R(\nabla^E),R(\nabla^F))+\partial\mathrm P\in\Omega^2(M,\mathfrak s(L(E,E)\oplus L(F,F))).\]
We write $\partial\mathrm P=(\Phi,\Psi)$ where $\Phi\in\Omega^2(M,L(E,E))$ and $\Psi\in\Omega^2(M,L(F,F))$.
A computation shows that 
\begin{align*}
(\Phi{}^{A}_{A'}{}^{B}_{B'}{}^{C}_{D},\Psi{}^{A}_{A'}{}^{B}_{B'}{}^{C'}_{D'})=(\delta^{B}_{D}\,\mathrm P{}^{A}_{A'}{}^{C}_{B'}-\delta^{A}_{D}\,\mathrm P{}^{B}_{B'}{}^{C}_{A'},-\delta^{C'}_{B'}\,\mathrm P{}^{A}_{A'}{}^{B}_{D'} +\delta^{C'}_{A'}\,\mathrm P{}^{B}_{B'}{}^{A}_{D'}).
\end{align*}
Hence there holds \eqref{Weyl tensor expanded} and so
\begin{equation}\label{e}
\begin{aligned}
\begin{dcases}
W{}^{A}_{A'}{}^{I}_{B'}{}^{B}_{I}
    &=R(\nabla^E)^{A}_{A'}{}^{I}_{B'}{}^{B}_{I}+2\,\mathrm P{}^A_{A'}{}^B_{B'}-\mathrm P{}^A_{B'}{}^B_{A'}\\
W'{}^{A}_{A'}{}^{B}_{I'}{}^{I'}_{B'}
    &=R(\nabla^F)^{A}_{A'}{}^{B}_{I'}{}^{I'}_{B'}-n\,\mathrm P{}^A_{A'}{}^B_{B'}+\mathrm P{}^B_{A'}{}^A_{B'}.
\end{dcases}
\end{aligned}
\end{equation}
Recall from \cref{Weyl tensor} that $W{}^{A}_{A'}{}^{I}_{B'}{}^{B}_{I}=W'{}^{A}_{A'}{}^{B}_{I'}{}^{I'}_{B'}$. Thus from \eqref{e} we obtain
\begin{align*}
Ric(\nabla)^A_{A'}{}^B_{B'}=(n+2)\,\mathrm P{}^A_{A'}{}^B_{B'}-\mathrm P{}^A_{B'}{}^B_{A'}-\mathrm P{}^B_{A'}{}^A_{B'}.
\end{align*}
This equality is equivalent to \eqref{Rho Ric} as each of them is equivalent to the following collection of relations.
\begin{align*}
Ric(\nabla)^{(A}_{(A'}{}^{B)}_{B')}&=n\,\mathrm P{}^{(A}_{(A'}{}^{B)}_{B')}&Ric(\nabla)^{[A}_{[A'}{}^{B]}_{B']}&=(n+4)\,\mathrm P{}^{[A}_{[A'}{}^{B]}_{B']}\\
            Ric(\nabla)^{(A}_{[A'}{}^{B)}_{B']}&=(n+2)\,\mathrm P{}^{(A}_{[A'}{}^{B)}_{B']}&Ric(\nabla)^{[A}_{(A'}{}^{B]}_{B')}&=(n+2)\,\mathrm P{}^{[A}_{(A'}{}^{B]}_{B')};        
\end{align*}
cf. the end of \cite{CM}*{Subsection 4.3}
\end{proof}
We provide an equivalent description to \cref{charpath}.  
\begin{corollary}\label[corollary]{altermative char path}
Let $(M,E,F)$ be an almost Grassmann structure of type $(2,n)$, $n\geq 3$. Denote by $\tau$ the harmonic torsion of $(M,E,F)$ and denote by $\rho$  the harmonic curvature of $(M,E,F)$. Given nowhere vanishing sections $\eta\in\Gamma(F)$ and $\varphi\in\Gamma(E^*)$, each of the following conditions is well-defined.
\begin{enumerate}
\item [$\bullet$] $(\nabla\eta)_o=0$ for one Weyl connection $\nabla$ and hence for all Weyl connections. Equivalently, $\eta$ is parallel for some Weyl connection.
\item [$\bullet$]
$tr(\iota_{\tau}\tau)(\eta)=0$ where we view $tr(\iota_{\tau}\tau)\in\Gamma(F,T^*M\otimes E)$.
\item [$\bullet$] $(\nabla\varphi)_o=0$ for one Weyl connection $\nabla$ and hence for all Weyl connections. Equivalently, $\varphi$ is parallel for some Weyl connection.
\item [$\bullet$]
$tr(\iota_{\tau}\tau)(\varphi)=0$ where we view $tr(\iota_{\tau}\tau)\in\Gamma(E^*,T^*M\otimes F^*)$.
\item [$\bullet$]
$\frac 1 2\,\eta^{A'}\,\nabla^A_{A'}\varphi_{A}-\frac 1 n\,\varphi_A\,\nabla^A_{A'}\eta^{A'}=1$.
\item [$\bullet$] $\rho(\varphi\otimes\eta,T M)=0$.
\item [$\bullet$] $\tau(\varphi\otimes F,\varphi\otimes F)\subseteq\varphi\otimes F$.
\item [$\bullet$] $\tau(\varphi\otimes\eta,T M)=0$ and, in addition, $(\varphi_A\,\eta^{B'}\,W{}^A_{A'}{}^B_{B'}{}^{C}_{D},\varphi_A\,\eta^{B'}\,
W'{}^A_{A'}{}^B_{B'}{}^{C'}_{D'})=0$ for the Weyl tensor $(W,W')$ of a Weyl structure and hence for the Weyl tensors of all Weyl structures.
\end{enumerate}
Moreover, $(M,E,F)$ is locally the result of the Fefferman-type construction on a path geometry if and only if there exists nowhere vanishing sections $\eta\in\Gamma(F)$ and $\varphi\in\Gamma(E^*)$ such that all the above conditions hold.
\end{corollary}
\begin{proof}
Denote by $\mathcal T$ the standard tractor of $(M,E,F)$. Then its dual $\mathcal T^*$ is the standard cotractor of $(M,E,F)$. Recall from \cite{Guo}*{Proposition 5} that the trace-free part $(\nabla\eta)_o$ of $\nabla\eta\in\Gamma(E\otimes F^*\otimes F)$ is independent of the choice of Weyl connection $\nabla$ and the assignment $\eta\mapsto(\nabla\eta)_o$ is precisely the first Bernstein-Gelfand-Gelfand (BGG) operator of $\mathcal T$ and, similarly, the trace-free part $(\nabla\varphi)_o$ of $\nabla\varphi\in\Gamma(E\otimes F^*\otimes E^*)$ is independent of the choice of Weyl connection $\nabla$ and the assignment $\varphi\mapsto(\nabla\varphi)_o$ is precisely the first BGG operator of the standard cotractor bundle. Moreover, recall from \cite{Guo}*{Proposition 5} that if $(\nabla\eta)_o=0$, then $\eta$ is parallel for some Weyl connection and, similarly, $(\nabla\varphi)_o=0$, then $\varphi$ is parallel for some Weyl connection. Thus we see that the first and the third conditions in the corollary are well-defined.

The symmetry \eqref{harmonic components AG} of $\tau$ implies that
\begin{align*}
\tau(\varphi\otimes\eta,TM)=0\quad\iff\quad\tau(\varphi\otimes F,E^*\otimes\eta)=0.
\end{align*}
In this case, it is easy to see that $(W,W')(\varphi\otimes F,E^*\otimes\eta)=0$ for the Weyl tensor $(W,W')$ of a Weyl structure implies that the same equality holds for the Weyl tensor $(W,W')$ of every Weyl structure. This shows that the last condition in the corollary is well-defined. That the other conditions make sense is obvious.

Recall from \cite{Guo}*{Remark 6 and Theorem 8} the following. The natural projection $\mathcal T\twoheadrightarrow F$ induces a one-to-one correspondence
\begin{align*}
\{\mu\in\Gamma(\mathcal T)\text{ parallel}\}
\longleftrightarrow\{\eta\in\Gamma(F):(\nabla\eta)_o=0\text{ and } tr(\iota_{\tau}\tau)(\eta)=0\}.
\end{align*}
The inverse of this correspondence is the splitting operator $L:\Gamma(F)\to\Gamma(\mathcal T)$. Similarly, the natural projection $\mathcal T^*\twoheadrightarrow E^*$ induces a one-to-one correspondence
\begin{align*}
\{\nu\in\Gamma(\mathcal T^*)\text{ parallel}\}
\longleftrightarrow\{\varphi\in\Gamma(E^*):(\nabla\varphi)_o=0\text{ and } tr(\iota_{\tau}\tau)(\varphi)=0\}.
\end{align*}
The inverse of this correspondence is the splitting operator $L:\Gamma(E^*)\to\Gamma(\mathcal T^*)$. Using the formulae of the two splitting operators in \cite{Guo}*{Proposition 5}, one computes that
\[(L\,\varphi)(L\,\eta)=\frac 1 2\,\eta^{A'}\,\nabla^A_{A'}\varphi_{A}-\frac 1 n\,\varphi_A\,\nabla^A_{A'}\eta^{A'}.\] In particular, we see that the first five conditions in the corollary are equivalent to the conditions in \cref{charpath} (i).

Denote by $\kappa$ the curvature of the normal Cartan geometry $(\mathcal G\to M,\omega)$ of type $(G,P)$ canonically associated to $(M,E,F)$. Recall from \cite{Cap06}*{Subsection 3.3\ Proposition} that there holds $\kappa(\varphi\otimes\eta,TM)=0$ if and only if both $\tau(\varphi\otimes\eta,TM)=0$ and $\rho(\varphi\otimes\eta,TM)=0$. In this case, there holds $\tau(\varphi\otimes F,E^*\otimes\eta)=0$ so that, in particular, there holds $\kappa(\varphi\otimes\eta,TM)\subseteq\mathcal G\times_P\mathfrak p_+$ if and only if $(\varphi_A\,\eta^{B'}\,W{}^A_{A'}{}^B_{B'}{}^{C}_{D},\varphi_A\,\eta^{B'}\,
W'{}^A_{A'}{}^B_{B'}{}^{C'}_{D'})=0$ for the Weyl tensor $(W,W')$ of a Weyl structure and hence for the Weyl tensors of all Weyl structures. Finally, it is well-known that $\tau(\varphi\otimes F,\varphi\otimes F)\subseteq\varphi\otimes F$ if and only if $\varphi\otimes F$ is involutive; see, e.g., \cite{Book}*{Lemma 1.5.14}. We conclude that the last two conditions in the corollary are equivalent to the conditions in \cref{charpath} (ii). This completes the proof of the corollary.
\end{proof}
In view of \cref{Weyl tensor Rho Ricci} and the paragraph preceding it, \cref{altermative char path} is a characterization in terms of a Weyl connection and its associated tensors; it does not involve Cartan geometry. Moreover, each condition there is independent of the choice of Weyl connection.
\subsection{The case $n=1$}\label{LC}
The group homomorphism $i:G\to\tilde G$ given in \eqref{i} can be used to define a functor associating to every $(2\,n+1)$-dimensional generalized path geometry $(M,E,V)$ an almost Grassmann structure of type $(2,n+1)$. We have studied such constructions for all $n>2$ as well as the restrictions of the construction over all geometries $(M,E,V)$ where $n=2$ and $V$ is involutive, i.e., where $(M,E,V)$ is a path geometry. The cases where $n=1$ and where $n=2$ with $V$ non-involutive have been excluded from our previous discussion due to their different types of fundamental invariants; see the beginning of the article. We discuss the case where $n=1$ below.

First, we recall the following background from, e.g., \cite{Book}*{Subsections 4.1.3, 4.1.4 and 4.2.3}. In dimension $2\cdot 1+1=3$, the concepts of generalized path geometry and of Lagrangian contact structure coincide. Such a geometry $(M,E,V)$ is automatically torsion-free. Its fundamental invariants are given by two Cotten-York-type tensors in
\begin{align}\label{Cotten-York}
    \Gamma(E^*\otimes(TM/(E\oplus V))^*\otimes E^*)\quad\text{and}\quad\Gamma(V^*\otimes(TM/(E\oplus V))^*\otimes V^*),
\end{align}
respectively. On the other hand, a nonzero volume form on $\mathbb R^4$ gives rise to a natural identification $\wedge^2\,\mathbb R^4\cong\mathbb R^{(3,3)}$, which in turn induces an isomorphism
\[\tilde G=SL(4,\mathbb R)\cong Spin(3,3).\]
It restricts to an isomorphism between the structure groups
\begin{align*}
\tilde P_0\cong CSpin(2,2).
\end{align*}
Here, $\tilde P_0$ denotes the Levi subgroup of $\tilde P$, that is, the block diagonal subgroup of $SL(4,\mathbb R)$ where the block size is $(2,2)\times (2,2)$. In particular, an almost Grassmann structure $(\tilde M,\tilde E,\tilde F)$ of type $(2,2)$ is the same as a conformal spin structure of signature $(2,2)$. Moreover, $(\tilde M,\tilde E,\tilde F)$ is automatically torsion-free. Its fundamental invariants are a pair of harmonic curvatures in
\begin{align*}
    \Gamma(\wedge^2\,\tilde E\otimes \mathrm{Sym}^2\,\tilde F^*\otimes\mathfrak{sl}(\tilde F))\quad\text{and}\quad\Gamma(\mathrm{Sym}^2\,\tilde E\otimes \wedge^2 \tilde F^*\otimes\mathfrak{sl}(\tilde E)),
\end{align*}
respectively.

Recall in addition that the spin representation of $Spin(3,3)$ corresponds to the natural representation of $SL(4,\mathbb R)$ on $\mathbb R^4\oplus\mathbb R^{4*}$ and the Clifford multiplication decomposes into the maps
\begin{align*}
    \wedge^2\,\mathbb R^4\times\mathbb R^4\to\mathbb R^{4*}
    \quad\text{and}\quad
    \wedge^2\,\mathbb R^4\times\mathbb R^{4*}\to\mathbb R^4
\end{align*}
canonically induced by the volume form on $\mathbb R^4$ chosen before. In particular, every spinor $0\neq s\in\mathbb R^4\oplus\mathbb R^{4*}$ is pure, that is, it is annihilated by a maximal isotropic subspace of $\wedge^2\,\mathbb R^4$. Moreover, the almost Grassmann standard tractor bundle $\tilde{\mathcal T}$ is the positive conformal spin tractor bundle and its dual $\tilde{\mathcal T}^*$ is the negative conformal spin tractor bundle.

The authors of \cite{Ham+17} introduced a Fefferman-type construction based on a group homomorphism 
\[SL(n+2,\mathbb R)\to Spin(n+2,n+2)\] 
that associates to every $(2\,n+1)$-dimensional Lagrangian contact structure a conformal spin structure of signature $(n+1,n+1)$; see \cite{Ham+17}*{Subsection 3.2}. In the case where $n=1$, the group homomorphism introduced in \cite{Ham+17}*{Subsection 3.2} coincides with our homomorphism \[i:SL(3,\mathbb R)\to SL(4,\mathbb R)\cong Spin(3,3);\] see \eqref{i}. In particular, the two Fefferman-type constructions are the same in this dimension up to minor differences on the choice of parabolic subgroups $Q\subseteq G$ and $\tilde P\subseteq\tilde G$.

We obtain the following results, which agree with \cite{Ham+17}*{Proposition 3.7 and Proposition 3.10} under the setup of \cite{Ham+17}*{Subsection 3.2} with $n=2$ in their work. As mentioned above, every tractor spinor is pure in this specific dimension.

\begin{theorem}
The Fefferman-type construction from $3$-dimensional generalized path geometries to almost Grassmann structures of type $(2,2)$ determined by \eqref{i} is normal.

Moreover, an almost Grassmann structure $(\tilde M,\tilde E,\tilde F)$ of type $(2,2)$ is locally the result of the Fefferman-type construction on a $3$-dimensional generalized path geometry if and only if the conditions in \cref{charpath} (i) are satisfied. In terms of conformal spin structure, these conditions reads as follows:

There are pure parallel tractor spinors $\mu\in\Gamma(\tilde{\mathcal T})$ and $\nu\in\Gamma(\tilde{\mathcal T}^*)$ with $\mu(\nu)=1$. Moreover, the projection $\eta\in\Gamma(\tilde F)$ of $\mu$ and $\varphi\in\Gamma(\tilde E^*)$ of $\nu$ to the underlying twistor spinors are both nowhere vanishing.
\end{theorem}

\begin{proof}
Since the two Cotten-York-type invariants \eqref{Cotten-York} above are still sections of the $Q$-module $\mathbb F$ defined in the proof of \cref{path} (i), the statement of \cref{path} (i) holds for $n=1$ as well. Since the proof of \cref{path Fef is normal} is a computation based on \cref{path} (i), it applies to the case $n=1$. Hence the statement of \cref{path Fef is normal} holds true for $n=1$ too. This proves the first assertion of the theorem.

The proof of \cref{charpath} can be directly applied to the case $n=1$. Since $(\tilde M,\tilde E,\tilde F)$ is automatically torsion-free, \cref{charpath} gives the second assertion of the theorem.
\end{proof}

The Fefferman-type construction induced by the related homomorphism \[SL(n+2,\mathbb R)\to SO(n+2,n+2)\] associating to every Lagrangian contact structure a conformal structure has been studied in \cite{Ma+}. Note that \[Spin(n+2,n+2)\to SO(n+2,n+2)\] is a two-fold covering. In particular, we obtain the following from the last relation in the proof of \cite{Ma+}*{Theorem 3.3}. Given a $3$-dimensional generalized path geometry $(M,E,V)$ with canonically associated normal regular parabolic geometry $(\mathcal G\to M,\omega)$, there exists local coordinates $(x,u,p)=U\subseteq M$ and a local section $\phi\in\Gamma_U(\mathcal G)$ such that
\begin{align*}
\phi^*\omega=\left(\begin{array}{ccc}
		*&*&*\\
		\theta&\alpha&*\\
		\sigma&\pi&*
	\end{array}\right)\in\Omega^1(U,\mathfrak {sl}(3,\mathbb R)).
\end{align*}
Here,
\begin{align*}
    \sigma=du-p\,dx,\quad\theta=dx,\quad\pi=dp-f\,dx,\quad\alpha=\frac 1 6\frac{\partial^2f}{\partial p^2}\sigma+\frac 2 3\frac{\partial f}{\partial p}\theta
\end{align*}
and $f=f(x,u,p)$ is the defining function of the geometry with respect to the given local coordinates; see \cite{DMT}*{(2.1)}. On the other hand, the projection of $\phi\in\Gamma_U(\mathcal G)$ to $\underline\phi\in\Gamma_U(\tilde M)$ defines a local trivialization of the principal bundle $\tilde M=GL_+(E^*\otimes V)$. Since the isomorphism between $Q/H$ and the structure group $\mathbb R_+$ of $\tilde M\to M$ is defined by the surjective homomorphism
\begin{align*}
    Q\to\mathbb R_+,\quad\left(\begin{array}{ccc}
		*&*&*\\
		0&a&*\\
		0&0&*
	\end{array}\right)\mapsto a^{-3}
\end{align*}
with kernel $H$, the isomorphism
\begin{align*}
\mathbb R\setminus\{0\}\xrightarrow{\cong}\mathbb R_+,\quad X\mapsto e^{-3\,X}  
\end{align*}
defines a local trivialization $(x,u,p,\tilde s)=\underline\phi\times(\mathbb R\setminus\{0\})\subseteq\tilde M$ such that for all local section $\psi$ of $\mathcal G\to\tilde M$ with $\psi\circ\underline\phi=\phi$, there holds
\begin{align*}
\psi^*\omega(\frac{\partial}{\partial s})=\left(\begin{array}{ccc}
		*&*&*\\
		0&1&*\\
		0&0&*
	\end{array}\right).
\end{align*}
In particular,
\begin{align*}
\psi^*\omega=\left(\begin{array}{ccc}
		*&*&*\\
		\theta&\alpha+d\tilde s&*\\
		\sigma&\pi&*
	\end{array}\right)\in\Omega^1(\underline\phi(U)\times\mathbb R_+,\mathfrak{sl}(3,\mathbb R))
\end{align*}
We have continued to use $\sigma,\theta,\pi,\alpha$ for their respective pullbacks to $\Omega^1(\underline\phi(U)\times\mathbb R_+)$. Hence \[\left(\begin{array}{cc}
		\theta&\alpha+d\tilde s\\
		\sigma&\pi
	\end{array}\right)\]
is a local coframe on $\tilde M$ adapted to the almost Grassmann structure $(\tilde M,\tilde E,\tilde F)$. Moreover, the metric in the conformal class of the geometry is given by
\begin{align*}
    \tilde g&=\theta\odot\pi-\sigma\odot(\alpha+d\,\tilde s)\\
    &=\frac 1 2(\theta\otimes\pi+\pi\otimes\theta-\sigma\otimes(\alpha+d\,\tilde s)-(\alpha+d\,\tilde s)\otimes\sigma).
\end{align*}
This agrees with the formula of the metric provided in  \cite{Ma+}*{Theorem 3.3 and (6.4)}. Note that they have chosen a different trivialization on the standard fiber of $\tilde M\to M$; see \cite{Ma+}*{(3.5)}.

\section{The Fefferman-type construction on almost Grassmann structures}\label{A generalization}
We follow \cref{convention} for the notation of Lie groups, Lie algebras and gradings. Recall that every almost Grassmann structure
\[(M,E,F)\]
of type $(2,n)$ is canonically associated to a normal parabolic geometry
\[(\mathcal G\to M,\omega)\]
of type $(G,P)$. Replacing $Q$ by $P$ in the argument in the first paragraph of \cref{The Fefferman-type constructions}, we see that the Fefferman-type functor defined by $i:G\to\tilde G$ in \eqref{i} associates to $(\mathcal G\to M,\omega)$ a parabolic geometry \[(\tilde{\mathcal G}\to\tilde M,\tilde\omega)\] of type $(\tilde G,\tilde P)$ with underlying almost Grassmann structure \[(\tilde M,\tilde E,\tilde F)\] of type $(2,n+1)$. This defines a functorial construction
associating to every almost Grassmann structure of type $(2,n)$ an almost Grassmann structure of type $(2,n+1)$. We denote by $\kappa$ and $\tilde\kappa$ the curvatures of $\omega$ and $\tilde\omega$, respectively. Denote by 
\[j:\mathcal G\to\tilde{\mathcal G}\]
the natural map defining $\tilde{\mathcal G}$ and denote by
\[\pi:\tilde{\mathfrak g}/\tilde{\mathfrak p}\twoheadrightarrow\mathfrak g/\mathfrak p\]
the projection canonically induced by the Lie algebra map $i'$ of $i:G\to\tilde G$; see \cref{section abstract Fefferman}.

We provide a geometric description of $\tilde M,\tilde E$ and $\tilde F$.
\begin{theorem}\label[theorem]{geometry of Fef on AG}
Notation as above. The following canonical identifications hold.
\begin{enumerate}
\item [$\bullet$] The fiber bundle $p:\tilde M\to M$ is isomorphic to $E^*\setminus\{0\}$, i.e., the complement of the zero section in $E^*$.
\item [$\bullet$] $\tilde E$ is the pullback bundle $p^*E$. Additionally, $\tilde E^*$ is the vertical bundle of $p:\tilde M\to M$.
\item [$\bullet$] $\tilde F=(p^*\mathcal T)/\ell$. Here, $p^*\mathcal T$ is the pullback of the standard tractor bundle $\mathcal T$ of $(M,E,F)$ and $\ell$ is the tautological subbundle defined by 
\[\ell(\varphi_x)=\ker\,\varphi_x\subseteq E_x\subseteq\mathcal T_x\]
over every $\varphi_x\in E_x^*\setminus
\{0\}\subseteq \tilde M$ for each $x\in M$.
\end{enumerate}    
\end{theorem}
\begin{proof}
Notation as at the beginning of the section. Consider the natural representation of $P$ on $\mathbb R^{2*}$. The isotropy subgroup fixing the base point $(0,1)\in\mathbb R^{2*}$ is precisely $H$ and the $P$-orbit of the base point is $\mathbb R^{2*}\setminus\{0\}$. Since $E^*=\mathcal G\times_P\mathbb R^{2*}$,
\[\tilde M=\mathcal G/H=\mathcal G\times_PP/H=E^*\setminus\{0\}\]
as fiber bundles over $M$. This proves the first assertion of the theorem. Note that, in contrast to \cref{bundles}, there is no natural principal bundle structure on $p:\tilde M\to M$ because $H$ is not a normal subgroup of $P$.

Using the same argument from the proof of \cref{bundles}, we obtain the relations
\begin{align*}
\tilde E=\mathcal G\times_H\mathbb R^2
\quad\text{and}\quad
\tilde F=\mathcal G\times_H\mathbb R^{n+2}/\mathbb R
\end{align*}
where 
\[\mathbb R=(\mathbb R,0)^t\subseteq\mathbb R^2\subseteq\mathbb R^{n+2}\]
is the filtration determined by the natural representation of $H$ on $\mathbb R^{n+2}$. In particular, $\tilde E=p^*E$. Moreover, the isomorphism $\mathbb R^{2*}\cong\mathfrak p/\mathfrak h$ of $H$-modules implies that $\tilde E^*$ is canonically identified with the vertical distribution $\mathcal G\times_H\mathfrak p/\mathfrak h\subseteq T\tilde M$. This completes the proof of the second assertion of the theorem.

Finally, there holds $\tilde F=(p^*\mathcal T)/\ell$ because $\mathcal G\times_H\mathbb R^{n+2}=p^*\mathcal T$ and $\mathcal G\times_H(\mathbb R,0)^t=\ell$. This proves the final assertion.
\end{proof}

\begin{remark}\label[remark]{distinguished sections AG}
Observe that the sections
\begin{align*}
\varphi=\mathcal G\times_H(0,1)\in\Gamma(\tilde E^*)
\quad\text{and}\quad
\eta=\mathcal G\times_H(-1,0^n)^t\in\Gamma(\tilde F)
\end{align*}
we specified in \cref{distinguished sections} are also well-defined in the above situation. Here, $\varphi$ is the tautological section of $\tilde E^*$ so that its rank-$1$-kernel is precisely the line bundle
\[\ell\subseteq\tilde E\]
defined in \cref{geometry of Fef on AG}. Moreover, the vertical distribution of $p:\tilde M\to M$ is
\[\tilde V=\tilde E^*\otimes\eta\subseteq T\tilde M.\]
\end{remark}

\begin{convention}\label{H0 invariant LA decomp}
We define a decomposition
    \begin{align*}
\tilde{\mathfrak g}=\mathfrak {sl}(n+3,\mathbb R)=\left(\begin{array}{ccc|c|ccc}
	&\tilde{\mathfrak n}^E_{0}&&\tilde{\mathfrak n}_{1}^E&\tilde{\mathfrak n}_{2}&\\
 &&&&&\\
 \hline &\tilde{\mathfrak n}_{-1}^E&&\tilde{\mathfrak n}^\eta_{0}&\tilde{\mathfrak n}_{1}^F&\\
	\hline&&&&&\\
    &\tilde{\mathfrak n}_{-2}&&\tilde{\mathfrak n}_{-1}^F&\tilde{\mathfrak n}^F_{0}&\\
&& &&&\end{array}\right)
\end{align*}
in the block size $(2,1,n)\times(2,1,n)$. Observe that it is a grading decomposition of some parabolic subgroup of $\tilde G$ that contains $H$. In particular, its induced filtration is $H$-invariant and each block component is invariant under the natural representation of
\begin{align}\label{H0}
H_0=P_0\cap H=\tilde P_0\cap i(H).
\end{align}
Observe that
\begin{align*}
\tilde{\mathfrak n}^{1,F}=\tilde{\mathfrak n}_{1}^F\oplus\tilde{\mathfrak n}_{2}
\end{align*}
is an $H$-submodule of $\tilde{\mathfrak g}$. 
\end{convention}
 
We say that the Fefferman-type construction on an almost Grassmann structure $(M,E,F)$ of type $(2,n)$ is normal if the Cartan connection $\tilde\omega$ given at the beginning of the section is normal. We prove that this notion of being normal is equivalent to a certain property of the torsion $\tau$ of $(M,E,F)$. Recall from \cref{Weyl tensor} the contraction $tr(\iota_\tau\tau)\in\Gamma(\mathrm{Sym}^2\,T^*M)$ of $\tau\otimes\tau$.

\begin{proposition}\label{AG Fef normality}
The Fefferman-type construction on $(M,E,F)$ is normal if and only if there holds $tr(\iota_\tau\tau)=0$.
\end{proposition}

\begin{proof}
Notation as at the beginning of the section. For every $\tilde x\in\tilde M$, we may choose $u\in\mathcal G$ such that $j(u)\in\tilde{\mathcal G}$ projects to $\tilde x$. Denote by $x\in M$ the projection of $\tilde x$ along $p:\tilde M\to M$. Choose a Weyl structure $\sigma:\mathcal G_0=\mathcal G/P_+\to\mathcal G$ whose image contains $u$ so that, in particular, we consider the identification between $\mathfrak g$ and the fiber of $\mathcal G\times_P\mathfrak g$ over $x\in M$ defined by $u$. Thus
\begin{align}
\kappa(u)=\left(\begin{array}{cc|cc}
W(x){}^A_{A'}{}^B_{B'}{}^{C}_{D}&&
Y(x){}^A_{A'}{}^B_{B'}{}^{C}_{D'}&\\
 \hline&&&\\
\tau(x){}^A_{A'}{}^B_{B'}{}^{C'}_{D}&&
W'(x){}^A_{A'}{}^B_{B'}{}^{C'}_{D'}&\\
 &&&
\end{array}\right).
\end{align}
Here, $(W,W')$ is the Weyl tensor of $\sigma$ and $Y$ is the Cotten-York tensor of $\sigma$; see \eqref{induced Weyl objects}. These Weyl objects evaluated at $x$ are independent of the choice of the Weyl structure mentioned above. We view
\[\varphi=\mathcal G\times_H(0,1)\in\Gamma(\tilde E^*)\]
defined in \cref{distinguished sections AG} as a section of $p^*E^*=\tilde E^*$; see \cref{geometry of Fef on AG}. Using the identification between $\tilde{\mathfrak g}$ and the fiber of $\tilde{\mathcal G}\times_{\tilde P}\tilde{\mathfrak g}$ over $\tilde x\in\tilde M$ defined by $j(u)\in\tilde{\mathcal G}$, there holds
\begin{align*}
\tilde\kappa(j(u))&=\left(\begin{array}{c|c|ccc}
	W(x){}^A_{A'}{}^B_{B'}{}^{C}_D&0&Y(x){}^A_{A'}{}^B_{B'}{}^{C}_{D'}&\\
 &&&\\
 \hline\varphi(\tilde x)_I\, W(x){}^A_{A'}{}^B_{B'}{}^{I}_D&0&\varphi(\tilde x)_I\, Y(x){}^A_{A'}{}^B_{B'}{}^{I}_{D'}&\\
\hline\tau(x){}^A_{A'}{}^B_{B'}{}^{C'}_{D}
    &0&W'(x){}^A_{A'}{}^B_{B'}{}^{C'}_{D'}&\\
 &&&
\end{array}\right)\circ\wedge^2\,\pi.
\end{align*}
The decomposition of $\tilde{\mathfrak g}$ is as given in \cref{H0 invariant LA decomp}.

Denote by $\tilde\partial^*$ the Kostant codifferential of $(\tilde{\mathfrak g},\tilde{\mathfrak p})$. We show that
\begin{align}\label{AG Fef curvature Kostant}
\frac 1 2\,(\tilde\partial^*\tilde\kappa)(j(u))
&=
\left(\begin{array}{ccc|c|ccc}
	&0&&0&0&\\
 &&&&&\\
 \hline &0&&0&-\varphi_J(\tilde x)\,W(x){}^A_{A'}{}^I_{C'}{}^{J}_{I}&\\
	\hline&0&&0&0&\\
&& &&&\end{array}\right)\circ \pi.
\end{align}


Indeed, the dual of $\pi:\tilde{\mathfrak g}/\tilde{\mathfrak p}\twoheadrightarrow\mathfrak g/\mathfrak p$ is 
\begin{align*}
    \pi^*:\mathfrak p_+\to\tilde{\mathfrak p}_+,\quad \pi^*(Z)_{ij}=Z_{i,j+1},\quad Z\in\mathfrak p_+.
\end{align*}
Note that it coincides with the restriction of $\pi^*$ defined in the proof of \cref{path Fef is normal}. We continue to denote by $\alpha:\mathfrak g\to\tilde{\mathfrak g}$ the embedding of $\mathfrak g$ into the Lie subalgebra of $\tilde{\mathfrak g}$ with zeros in the third row and the third column. Then
\begin{align*}
    \left[\pi^*\left(\begin{array}{cc|cc}
	0&&Z_{12}&\\
 \hline&&&\\
	0&&0&\\
 &&&
\end{array}\right),i'\left(\begin{array}{cc|cc}
	\Phi_{11}&&\Phi_{12}&\\
 \hline&&&\\
	\Phi_{21}&&\Phi_{22}&\\
 &&&
\end{array}\right)\right]
=
\alpha([Z,\Phi])-\left(\begin{array}{ccc|c|ccc}
	&0&&0&0&\\
 &&&&&\\
 \hline &0&&0&(0,1)\,\Phi_{11}Z_{12}&\\
	\hline&0&&0&0&\\
&& &&&\end{array}\right)
\end{align*}
for all $Z=Z_{12}\in\mathfrak p_+$ and $\Phi=(\Phi_{ij})_{1\leq i,j\leq 2}\in\mathfrak g$. The block size on the left-hand side is $(2,n)\times(2,n)$ and the block size on the right-hand side is $(2,1,n)\times(2,1,n)$. Following steps analogous to those succeeding \eqref{pi and alpha}, one obtains the expression above.

It follows from \cref{Weyl tensor} that $\varphi_J(\tilde x)\,W(x){}^A_{A'}{}^I_{C'}{}^{J}_{I}=0$ if and only if $\varphi_J(\tilde x)\,tr(\iota_\tau\tau)(x){}^A_{A'}{}^J_{C'}=0$, that is, there holds $(\tilde\partial^*\tilde\kappa)(\tilde x)=0$ if and only if $\varphi_J(\tilde x)\,tr(\iota_\tau\tau)(x){}^A_{A'}{}^J_{C'}=0$. Since the $P$-orbit of $(0,1)$ is $\mathbb R^{2*}\setminus\{0\}$, given any nonzero element $v\in E_x^*$, there is a unique point $\tilde x\in \tilde M$ over $x$ such that $\varphi(\tilde x)=v$ and hence there holds $\tilde\partial^*\tilde\kappa=0$ if and only if $tr(\iota_\tau\tau)=0$. This proves the statement of the proposition.
\end{proof}
Recall from, e.g., \cite{Book}*{Proposition 3.1.13} that we may normalize $(\tilde{\mathcal G}\to\tilde M,\tilde\omega)$ and thus obtain the normal parabolic geometry $(\tilde{\mathcal G}\to\tilde M,\tilde\omega^{nor})$ canonically associated to $(\tilde M,\tilde E,\tilde F)$. Denote the curvature of $\tilde\omega^{nor}$ by $\tilde\kappa^{nor}$. Let the group $H_0$ and the grading decomposition of $\tilde{\mathfrak g}$ be as in \cref{H0 invariant LA decomp}. Moreover, let $\tilde V\subseteq T\tilde M$ be the vertical distribution of the projection $\tilde M\to M$ as in \cref{distinguished sections AG}.
\begin{lemma}\label{normalization}
Notation as above. Then the following properties hold.
\begin{enumerate}
\item [$\bullet$] $j^*(\tilde\omega^{nor}-\tilde\omega)\in\Omega^1(\mathcal G,\tilde{\mathfrak n}^{1,F})^H$.
    \item [$\bullet$] $(\tilde\omega^{nor}-\tilde\omega)(\tilde V)=0$.
    \item [$\bullet$] $j^*(\tilde\kappa^{nor}-\tilde\kappa)\in\Omega^2(\mathcal G,[\tilde{\mathfrak g},\tilde{\mathfrak n}^{1,F}])^H$.
    \item [$\bullet$] $(j^*\tilde\kappa^{nor})(\tilde V,T\tilde M)\subseteq\mathcal G\times_H\tilde{\mathfrak n}^{1,F}$.
    \item [$\bullet$] $\tilde\kappa^{nor}(\tilde V,\tilde V)=0$.
\end{enumerate}   
\end{lemma}
\begin{proof}
Consider the $H$-modules
\begin{align*}
   \mathbb E&=\tilde{\mathfrak n}_{2}\otimes\tilde{\mathfrak n}^{1,F}\\
    &=\{\phi\in \tilde{\mathfrak p}_+\otimes\tilde{\mathfrak n}^{1,F}:\phi(i'(\mathfrak p))=0\}\quad\text{and}\\
    \mathbb F&=\tilde{\mathfrak n}_{1}^E\wedge\tilde{\mathfrak n}_{2}\otimes\tilde{\mathfrak n}^{1,F}\oplus\wedge^2\,\tilde{\mathfrak n}_{2}\otimes[\tilde{\mathfrak g},\tilde{\mathfrak n}^{1,F}]\\
    &=\{\psi\in \wedge^2\,\tilde{\mathfrak p}_+\otimes[\tilde{\mathfrak g},\tilde{\mathfrak n}^{1,F}]:\phi(i'(\mathfrak p),i'(\mathfrak p))=0\text{ and }\phi(i'(\mathfrak p),i'(\mathfrak g))\subseteq \tilde{\mathfrak n}^{1,F}\}.
\end{align*}
Note that the indicated decomposition of $\mathbb F$ is only $H_0$-invariant. Let $\tilde\partial$ be the Lie algebra differential and $\tilde\partial^*$ the Kostant codifferential of $(\tilde{\mathfrak g},\tilde{\mathfrak p})$. A computation shows that
\begin{align}\label{stability between E and F}
    \tilde\partial\,\mathbb E\subseteq\mathbb F\quad\text{and}\quad
    \tilde\partial^*\,\mathbb F\subseteq \mathbb E
\end{align}
so that, in particular, it follows from the Hodge decomposition \eqref{Hodge} that each of $\tilde\partial$ and $\tilde\partial^*$ induces a linear isomorphism
\begin{align}\label{isomorphism in Hodge}
    \mathrm{im}\,\tilde\partial^*\cap\mathbb E\xrightleftharpoons[\tilde\partial^*]{\tilde\partial}
\mathrm{im}\,\tilde\partial\cap\mathbb F.
\end{align}
in opposite directions. Here, $\tilde\partial$ is $H_0$-equivariant and $\tilde\partial^*$ is $H$-equivariant.

On the other hand, the expression \eqref{AG Fef curvature Kostant} in the proof of \cref{AG Fef normality} implies that
\begin{align}\label{space of nonnormality}    j^*(\tilde\partial^*\tilde\kappa)\in\Gamma(\mathcal G\times_{H}\mathbb E).
\end{align}
Since $\mathbb E^{(1)}=\mathbb E\subseteq\tilde{\mathfrak p}_1\otimes\tilde{\mathfrak p}^0$, we define 
\[\mathbb E^{(2)}=\mathbb E\cap(\tilde{\mathfrak p}_1\otimes\tilde{\mathfrak p}_1)=\tilde{\mathfrak n}_2\otimes\tilde{\mathfrak n}_2,\]
which is an $H$-submodule of $\mathbb E$. We follow the normalizing procedure in the proof of \cite{Book}*{Proposition 3.1.13} in our specific setting. Indeed, in view of \eqref{isomorphism in Hodge}, for every local trivialization $\tilde U\times H\subseteq\mathcal G$ with $\tilde U\subseteq\tilde M$ open, we may choose an $H$-equivariant map $\phi:\tilde U\times H\to\mathbb E$ such that the image of the $H$-equivariant map 
\[\tilde\partial^*\tilde\partial\phi
+j^*(\tilde\partial^*\tilde\kappa):
\tilde U\times H\to\mathbb E\]
is contained in $\mathbb E^{(2)}$. Using partition of unity, we obtain $\phi\in\Gamma(\mathcal G\times_H\mathbb E)$ such that
\[\tilde\partial^*\tilde\partial\phi
+j^*(\tilde\partial^*\tilde\kappa)\in\Gamma(\mathcal G\times_H\mathbb E^{(2)}).\]
$\phi$ can be extended uniquely to $\Gamma(\tilde{\mathcal G}\times_{\tilde P}(\tilde{\mathfrak p}_1\otimes\tilde{\mathfrak p}^0))$. Now define a Cartan connection
\[\tilde\omega'=\tilde\omega+\phi\circ\tilde\omega\] with curvature $\tilde\kappa'$. A computation using the formula of Cartan curvature in \cref{Cartan geometry definition} shows that
\[j^*(\tilde\kappa'-\tilde\kappa)\in\Gamma(\mathcal G\times_H\mathbb F).\]
Moreover, there holds
\[j^*(\tilde\partial^*\tilde\kappa')\in\Gamma(\mathcal G\times_H\mathbb E^{(1)}).\]
Repeating the above procedure once more, we obtain the normalized connection $\tilde\omega^{nor}$ which satisfies the following properties. 
\begin{align}\label{subspace for the difference}
j^*(\tilde\omega^{nor}-\tilde\omega)\in\Gamma(\mathcal G\times_H\mathbb E)
\quad\text{and}\quad
j^*(\tilde\kappa^{nor}-\tilde\kappa)\in\Gamma(\mathcal G\times_H\mathbb F).
\end{align}
We obtain all assertions of the lemma by putting \eqref{subspace for the difference}, the properties
\begin{align*}
j^*\tilde\omega\in\Omega^1(\mathcal G,i'(\mathfrak g))
\quad\text{and}\quad
\tilde\kappa(\tilde V,T\tilde M)=0    
\end{align*}
of the Fefferman-type construction together with the relations
\begin{align*}
   \tilde{\mathfrak n}^{1,F}\subseteq [\tilde{\mathfrak g},\tilde{\mathfrak n}^{1,F}]
   \quad\text{and}\quad
   [i'(\mathfrak p),\tilde{\mathfrak n}^{1,F}]\subseteq\tilde{\mathfrak n}^{1,F}.
\end{align*}
This completes the proof of the lemma.
\end{proof}
Consider the standard tractor bundle $\tilde{\mathcal T}=\tilde{\mathcal G}\times_{\tilde P}\mathbb R^{n+3}$. Let $\eta\in\Gamma(\tilde F)$, $\varphi\in\Gamma(\tilde E^*)$ and $\tilde V\subseteq T\tilde M$ be as in \cref{distinguished sections AG}. We obtain a refined filtration
\begin{align*}
\tilde E\subseteq\mathrm{span}(\eta)+\tilde E\subseteq \tilde{\mathcal T}
\end{align*}
of ranks $2\leq 3\leq n+3$. Its dual filtration is
\begin{align*}\mathrm{Ann}(\eta)\subseteq \tilde F^*\subseteq \tilde{\mathcal T}^*
\end{align*}
with ranks $n\leq n+1\leq n+3$. Here, $\mathrm{Ann}(\eta)$ denotes the annihilators of $\eta$ in $\tilde F^*$. Both the standard tractor bundle $\tilde{\mathcal T}$ and the standard cotractor bundle $\tilde{\mathcal T}^*$ are equipped with a tractor connection induced by $\tilde\omega^{nor}$. We denote them by $\nabla^{\tilde\omega^{nor}}$.
\begin{corollary}\label{tractor property AG}
There is a unique standard tractor $\mu\in\Gamma(\tilde{\mathcal T})$ with projection $\eta\in\Gamma(\tilde F)$ and a unique standard cotractor $\nu\in\Gamma(\tilde{\mathcal T}^*)$ with projection $\varphi\in\Gamma(\tilde E^*)$ such that $\nu(\mu)=1$, $\nabla^{\tilde\omega^{nor}}\mu=0$ and that
\[\nabla^{\tilde\omega^{nor}}\nu\in\{\Phi\in\Omega^1(\tilde M,\mathrm{Ann}(\eta))\subseteq\Omega^1(\tilde M,\tilde F^*):\Phi(\tilde V)=0\}.\]

Moreover, there holds
\[\tilde\kappa^{nor}(\tilde V,T\tilde M)\subseteq\tilde{\mathcal N}^{1,F}\]
where
\begin{align*}\tilde{\mathcal N}^{1,F}=\{\Psi\in\Gamma(\mathfrak{sl}(\tilde{\mathcal T})):\Psi(\mathrm{span}(\eta)+\tilde E)=0\text{ and }\Psi(\tilde{\mathcal T})\subseteq\mathrm{span}(\eta)+\tilde E\}.\end{align*}
\end{corollary}
\begin{proof}
The uniqueness of $\mu$ and $\nu$ follows from general theory; see \cref{BGG background} below. We now prove the existence of $\mu$ and $\nu$. To see this, observe that the argument in the proof of \cref{properties of Fef on path} (iii) applies here. Precisely, the sections 
\begin{align*}
\mu=\mathcal G\times_H(0,0,-1,0^n)^t\in\Gamma(\tilde{\mathcal T})
\quad\text{and}\quad
\nu=\mathcal G\times_H(0,1,-1,0^n)\in\Gamma(\tilde{\mathcal T}^*)
\end{align*}
are well-defined, projects to $\eta$ and $\varphi$, respectively, and such that $\nu(\mu)=1$. Moreover, denote by $\nabla^{\tilde\omega}$ the tractor connections on $\tilde{\mathcal T}$ and on $\tilde{\mathcal T}^*$ induced by the Cartan connection $\tilde\omega$ as given at the beginning of the section. We have proven that
\begin{align*}
\nabla^{\tilde\omega}\mu=0
\quad\text{and}\quad
\nabla^{\tilde\omega}\nu=0.
\end{align*}
On the other hand, observe that the relation \eqref{tractor connection on the distinguished tractor} in the proof of \cref{Hol} applies. Together with the properties of $\tilde\omega^{nor}-\tilde\omega$ in \cref{normalization}, we compute that
\begin{equation*}
    \begin{aligned}
(\nabla^{\tilde\omega^{nor}}_\xi\mu-\nabla^{\tilde\omega}_\xi\mu)(u)&\subseteq
\tilde{\mathfrak n}^{1,F}\,(0,0,-1,0^n)^t=0\\
(\nabla^{\tilde\omega^{nor}}_\xi\nu-\nabla^{\tilde\omega}_\xi\nu)(u)&\subseteq
\tilde{\mathfrak n}^{1,F}\,(0,1,-1,0^n)\subseteq(0,0,0,\mathbb R^n)\quad\text{and}\\
(\nabla^{\tilde\omega^{nor}}_{\tilde V}\nu-
\nabla^{\tilde\omega}_{\tilde V}\nu)(u)&=0           
    \end{aligned}
\end{equation*}
for all $\xi\in\mathfrak X(\tilde M)$ and $u\in\mathcal G$. Note that
\[\mathcal G\times_H(0,0,0,\mathbb R^n)=\mathrm{Ann}(\eta)\subseteq \tilde F^*\subseteq \tilde{\mathcal T}^*.\]
Combining all above relations, we conclude that 
\begin{align*}
\nabla^{\tilde\omega^{nor}}\mu=0,\quad
\nabla^{\tilde\omega^{nor}}\nu\in\Omega^1(\tilde M,\mathrm{Ann}(\eta))
\quad\text{and}\quad
\nabla^{\tilde\omega^{nor}}_{\tilde V}\nu=0,
\end{align*}
as desired.

Moreover, recall from \cref{normalization} that $(j^*\tilde\kappa^{nor})(\tilde V,T\tilde M)\subseteq\mathcal G\times_H\tilde{\mathfrak n}^{1,F}$. Observe that $\mathcal G\times_H\tilde{\mathfrak n}^{1,F}=\tilde{\mathcal N}^{1,F}$. This completes the proof of the corollary.
\end{proof}

With the consideration of \cref{modified correspondence} and our knowledge about $\tilde{\partial}^*\tilde\kappa$ from the proof of \cref{AG Fef normality}, we may modify \cref{charpath} carefully to obtain a characterization of all almost Grassmann structures of type $(2,n+1)$ that are locally the result of the Fefferman-type construction on an almost Grassmann structures of type $(2,n)$ via our Fefferman-type construction.

Now let $(\tilde M,\tilde E,\tilde F)$ be an almost Grassmann structure  of type $(2,n+1)$ and $(\tilde{\mathcal G}\to\tilde M,\tilde\omega^{nor})$ be the normal parabolic geometry of type $(\tilde G,\tilde P)$ canonically associated to it. Denote the curvature of $\tilde\omega^{nor}$ by $\tilde\kappa^{nor}$. Recall the short exact sequences
\begin{align*}
    0\to\tilde E\to\tilde{\mathcal T}\to\tilde F\to 0\quad\text{and}\quad 0\to \tilde F^*\to\tilde{\mathcal T}^*\to\tilde E^*\to 0
\end{align*}
where $\tilde{\mathcal T}$ is the standard tractor bundle of $(\tilde{\mathcal G}\to\tilde M,\tilde\omega^{nor})$ with dual $\tilde{\mathcal T}^*$. Each of the standard tractor bundle $\tilde{\mathcal T}$ and the standard cotractor bundle $\tilde{\mathcal T}^*$ are equipped with the tractor connection induced by $\tilde\omega^{nor}$. We denote them by $\nabla^{\tilde\omega^{nor}}$. Moreover, every nowhere vanishing section 
\[\eta\in\Gamma(\tilde F)\]
identifies $\tilde E^*$ with the distribution
\[\tilde V=\tilde E^*\otimes\eta\subseteq T\tilde M.\]
We also obtain a refined filtration
\begin{align*}
\tilde E\subseteq\mathrm{span}(\eta)+\tilde E\subseteq \tilde{\mathcal T}
\end{align*}
of ranks $2\leq 3\leq n+3$. Its dual filtration is
\begin{align*}\mathrm{Ann}(\eta)\subseteq \tilde F^*\subseteq \tilde{\mathcal T}^*
\end{align*}
of ranks $n\leq n+1\leq n+3$. Here, $\mathrm{Ann}(\eta)$ denotes the annihilators of $\eta$ in $\tilde F^*$. Let
\begin{align*}\tilde{\mathcal N}^{1,F}=\{\Psi\in\mathfrak{sl}(\tilde{\mathcal T}):\Psi(\mathrm{span}(\eta)+\tilde E)=0\text{ and }\Psi(\tilde{\mathcal T})\subseteq\mathrm{span}(\eta)+\tilde E\}.\end{align*}

\begin{theorem}\label[theorem]{charAG}
Notation as above. Then $(\tilde M,\tilde E,\tilde F)$ is locally the result of the Fefferman-type construction on an almost Grassmann structure of type $(2,n)$ if and only if the following hold.
    \begin{enumerate}
    \item [(i)]
    \begin{enumerate}
        \item [$\bullet$] There is a parallel standard tractor $\mu\in\Gamma(\tilde{\mathcal T})$ with nowhere vanishing projection $\eta\in\Gamma(\tilde F)$.
        \item [$\bullet$] There is a standard cotractor $\nu\in\Gamma(\tilde{\mathcal T}^*)$ such that
        \begin{align*}
            \nabla^{\tilde\omega^{nor}}\nu\in\{\Phi\in\Omega^1(\tilde M,\tilde F^*):\Phi(\tilde V)=0\text{ and }\Phi(T\tilde M)\subseteq \mathrm{Ann}(\eta)\}.
        \end{align*}
        Moreover, there holds $\nu(\mu)=1$ and the projection $\varphi\in\Gamma(\tilde E^*)$ of $\nu$ is nowhere vanishing.
    \end{enumerate}
    \item[(ii)]\begin{enumerate}
        \item [$\bullet$] $\tilde\kappa^{nor}(\tilde V,\tilde V)=0$ and
        \item [$\bullet$] $\tilde\kappa^{nor}(\tilde V,T\tilde M)\subseteq\tilde{\mathcal N}^{1,F}$.
        \end{enumerate} 
\end{enumerate}
\end{theorem}
\begin{proof}
That (i) and (ii) hold for a given Fefferman-type construction follows from \cref{normalization} together with \cref{tractor property AG}. It remains to show the converse direction. 

Indeed, consider the two grading decompositions of $\tilde{\mathfrak g}$ given in \cref{convention} and
\cref{H0 invariant LA decomp} and the $H$-modules $\mathbb E$ and $\mathbb F$ given in the proof of \cref{normalization}. Let $(\tilde M,\tilde E,\tilde F)$ be an almost Grassmann structure of type $(2,n+1)$, $(\tilde{\mathcal G}\to \tilde M,\tilde\omega^{nor})$ the normal parabolic geometry of type $(\tilde G,\tilde P)$ canonically associated to it and $\mu$ a tractor and $\nu$ a cotractor of $(\tilde{\mathcal G}\to \tilde M,\tilde\omega)$ satisfying the properties in (i) and (ii) of the statement of the theorem. We show the following.
\begin{enumerate}
    \item[$\bullet$] There is a reduction $j:\mathcal G\to\tilde{\mathcal G}$ to structure group $H$ and a Cartan geometry 
    \[(\mathcal G\to\tilde M,\omega)\]
    of type $(G,H)$ such that
    \begin{align*}
        i'\circ\omega-j^*\tilde\omega^{nor}=\phi\circ(j^*\tilde\omega^{nor})
    \end{align*}
    for some $\phi\in\Gamma(\mathcal G\times_H\mathbb E)$ where $\mathbb E$ is as given in the proof of \cref{normalization}. Moreover,
     \item[$\bullet$] $(\mathcal G\to\tilde M,\omega)$ is locally the result of the correspondence space modulo $L(\mathfrak g/\mathfrak p,\mathfrak p_+)$ of a normal parabolic geometry of type $(G,P)$; see \cref{locally correspondence}.
\end{enumerate}
Note that two Cartan connections of type $(\tilde G,\tilde P)$ on $\tilde{\mathcal G}\to\tilde M$ with difference in $\tilde{\mathfrak p}$ have the same underlying almost Grassmann structure; see, e.g., \cite{Book}*{Proposition 3.1.10 (1)}. In particular, since $\mathbb E\subseteq\tilde{\mathfrak p}_+\otimes\tilde{\mathfrak n}^{1,F}$ and $i'(\mathfrak p_+)\subseteq\tilde{\mathfrak n}^{1,F}\subseteq\tilde{\mathfrak p}$, the above two assertions imply that $(\tilde M,\tilde E,\tilde F)$ is locally isomorphic to the result of the Fefferman-type construction.

Indeed, note that a section $\mu$ of the standard tractor bundle and a section $\nu$ of the standard cotractor bundle such that $\nu(\mu)=1$ and that their projections $\eta$ and $\varphi$, respectively, are nowhere vanishing uniquely determine a reduction 
\[j:\mathcal G\to\tilde{\mathcal G}\]
to structure group $H$ such that 
\begin{align*}
    \mu=\mathcal G\times_H(0,0,-1,0^n)^t\quad\text{and}\quad \nu=\mathcal G\times_H(0,1,-1,0^n).
\end{align*}
This is proven in the third paragraph in the proof of \cref{charpath}. In particular, there holds
\begin{align}\label{H modules geometric meaning}
    \tilde V=\tilde E^*\otimes\eta=\mathcal G\times_H\frac{\tilde{\mathfrak n}_{-1}^E\oplus\tilde{\mathfrak p}}{\tilde{\mathfrak p}}\quad\text{and}\quad\mathcal G\times_H\tilde{\mathfrak n}^{1,F}=\tilde{\mathcal N}^{1,F}.
\end{align}
Moreover, the assumptions
\begin{align*}
    \nabla^{\tilde\omega^{nor}}\mu=0,
    \quad\nabla^{\tilde\omega^{nor}}_{\tilde V}\nu=0
    \quad\text{and}\quad
    \nabla^{\tilde\omega^{nor}}\nu\in\Omega^1(\tilde M,\mathrm{Ann}(\eta))
\end{align*}
in the statement of the theorem imply that 
\[j^*\tilde\omega^{nor}\in
\Omega^1(\mathcal G,\tilde{\mathfrak g})\] 
takes values in 
\begin{align*}
&\{A\in\tilde{\mathfrak g}:A\cdot(0,0,-1,0^n)^t=0\text{ and }(0,1,-1,0^n)\cdot A\in(0,0,0,\mathbb R^n)\}\\
&\quad=i'(\mathfrak g)+\tilde{\mathfrak n}^{1,F}
\end{align*}
where $\tilde{\mathfrak n}^{1,F}$ is given in \cref{H0 invariant LA decomp}. Moreover, the restriction of $j^*\tilde\omega^{nor}$ to the maximal distribution projecting to $\tilde V=\tilde E^*\otimes\eta\subseteq T\tilde M$ takes values in
\begin{align*}
&\{A\in\tilde{\mathfrak g}:A\cdot(0,0,-1,0^n)^t=0\text{ and }(0,1,-1,0^n)\cdot A=0\}\\
&\quad=i'(\mathfrak g).
\end{align*}
Let
\begin{align*}
q:i'(\mathfrak g)+\tilde{\mathfrak n}^{1,F}\twoheadrightarrow\mathfrak g
\end{align*}
be given by deleting the third row and the third column of $\tilde{\mathfrak g}$. Observe that $q|_{i'(\mathfrak g)}$ and $q|_{\tilde{\mathfrak n}^{1,F}}$ are $H$-equivariant so that, in particular,
\begin{align*}
    \omega=q\circ j^*\tilde\omega^{nor}\in\Omega^1(\mathcal G,\mathfrak g)
\end{align*}
is $H$-equivariant. Since $q\circ i'$ restricts to the identity map on $\mathfrak h$, $\omega$ reproduces the generator of fundamental vector fields. Finally, since $q$ is surjective, $\omega$ restricted to every tangent space is a linear isomorphism. Hence 
\begin{align}\label{converse direction first intermediate Cartan geometry}
(\mathcal G\to \tilde M,\omega)
\end{align}
is a Cartan geometry of type $(G,H)$. By construction,
\begin{align*}
    \omega-j^*\tilde\omega^{nor}=\phi\circ(j^*\tilde\omega^{nor})\in\Omega^1(\mathcal G,\mathfrak g)^H
\end{align*}
for some
\begin{align*}
    \phi\in\Gamma(\mathcal G\times_H\mathbb E).
\end{align*}
Denote by $\kappa$ the curvature of $\omega$. As mentioned in the proof of \cref{normalization}, a computation using the formula of Cartan curvature shows that
\begin{align*}
    i'\circ\kappa-j^*\tilde\kappa^{nor}\in\Gamma(\mathcal G\times_H\mathbb F).
\end{align*}
In particular, it follows from the assumptions in (ii) of the statement of the theorem that
\begin{align*}
    \kappa(\tilde V,\tilde V)&=0\quad\text{and}\\
    \kappa(\tilde V,T\tilde M)&\subseteq\mathcal G\times_H\mathfrak p_+.
\end{align*}
Note that $\mathfrak p_+=(i')^{-1}(\tilde{\mathfrak n}^{1,F})$.
Hence $(\mathcal G\to\tilde M,\omega)$ is locally isomorphic to the correspondence space of some Cartan geometry of type $(G,P)$ modulo $L(\mathfrak g/\mathfrak p,\mathfrak p_+)$; see \cref{modified correspondence}. It remains to show that such a Cartan geometry of type $(G,P)$ may be chosen to be normal. Without loss of generality, we may assume that there is a parabolic geometry
\begin{align*}
    (\mathcal G\to M,\omega')
\end{align*}
of type $(G,P)$ such that $\mathcal G/H=\tilde M$ and $\omega'=\omega$ modulo $L(\mathfrak g/\mathfrak p,\mathfrak p_+)$. Equivalently,
\begin{align*}
    \omega'-\omega=\phi'\circ\omega\in\Omega^1(\mathcal G,\mathfrak g)^H
\end{align*}
for some
\begin{align*}
    \phi'\in\Gamma(\mathcal G\times_HL(\mathfrak g/\mathfrak p,\mathfrak p_+))\subseteq\Gamma(\mathcal G\times_H\mathbb E).
\end{align*}
We have viewed $(\mathfrak g/\mathfrak p)^*\subseteq(\mathfrak g/\mathfrak h)^*\cong(\tilde{\mathfrak g}/\tilde{\mathfrak p})^*$. Note that $i'(\mathfrak p_+)\subseteq\tilde{\mathfrak n}^{1,F}$. Moreover, since $j^*\tilde\omega^{nor}$ and $\omega$ induce the same identification
\[T\tilde M\cong\mathcal G\times_H\mathfrak g/\mathfrak h=\mathcal G\times_H\tilde{\mathfrak g}/\tilde{\mathfrak p},\]
we conclude that
\begin{align*}
    i'\circ\omega'-j^*\tilde\omega^{nor}=(\phi+\phi')\circ(j^*\tilde\omega^{nor})
\end{align*}
and hence $i'\circ\kappa'-j^*\tilde\kappa^{nor}\in\Gamma(\mathcal G\times_H\mathbb F)$. Since $j^*\tilde\kappa^{nor}\in\Gamma(\mathcal G\times_H\mathbb F)$, there holds
\begin{align*}
    i'\circ\kappa'\in\Gamma(\mathcal G\times_H\mathbb F).
    \end{align*}

Now, denote by $\partial^*$ and $\tilde\partial^*$ the Kostant codifferential of $(\mathfrak g,\mathfrak p)$ and of $(\tilde{\mathfrak g},\tilde{\mathfrak p})$, respectively. Let $\beta:\tilde{\mathfrak g}\to\mathfrak g$ be as given in \eqref{Kostant relation for beta} in the proof of \cref{charpath}, i.e., $\beta$ is defined by deleting the third row and summing up the second and the third column of $\tilde{\mathfrak g}$. Since $\partial^*$ is the restriction of the Kostant codifferential on $(\mathfrak g,\mathfrak q)$, the relation \eqref{Kostant relation for beta} between $\tilde\partial^*$ and $\partial^*$ still holds. Hence
\begin{align*}
    \partial^*\kappa'=\beta\circ\tilde{\partial}^*(i'\circ\kappa').
\end{align*}
Since
\[\tilde{\partial}^*(i'\circ\kappa')\in\tilde{\partial}^*\,\mathbb F\subseteq\mathbb E\]
and the composition of an element in $\mathbb E$ with $\beta$ lies in $L(\mathfrak g/\mathfrak p,\mathfrak p_+)$, we conclude that
\[\partial^*\kappa'\in\Omega^1_{hor}(\mathcal G,\mathfrak p_+)^P.\]
Therefore, there exists a normal parabolic geometry
\[(\mathcal G\to M,\omega'_{nor})\]
of type $(G,P)$
such that $\omega'_{nor}=\omega'$ modulo $L(\mathfrak g/\mathfrak p,\mathfrak p_+)$. This completes the proof of the theorem.
\end{proof}

\begin{remark}\label{BGG background}
Notation as in \cref{charAG}. Note that by Bernstein-Gelfand-Gelfand (BGG) theory, $\nabla^{\tilde\omega^{nor}}\nu\in\Omega^1(\tilde M,\tilde{F}^*)$ means that $\nu$ is the image of $\varphi$ under the splitting operator associated to $\tilde{\mathcal T}^*$ and $\varphi$ is a solution of the so-called first BGG operator of $\tilde{\mathcal T}^*$. Similarly, $\nabla^{\tilde\omega^{nor}}\mu=0$ means that $\mu$ is the image of $\eta$ under the splitting operator associated to $\tilde{\mathcal T}$ and $\eta$ is a normal solution of the first BGG operator of $\tilde{\mathcal T}$. In fact, analogous to \cref{altermative char path}, the conditions in \cref{charAG} (i) can be expressed as conditions on sections $\varphi\in\Gamma(\tilde E^*)$ and $\eta\in\Gamma(\tilde F)$ in terms of a Weyl connection. To do so, one applies \cite{Guo}*{Proposition 5 and Lemma 7} in the situation of \cref{charAG} (i).
\end{remark}
%\appendix






% \bib, bibdiv, biblist are defined by the amsrefs package.
\begin{bibdiv}
	\begin{biblist}
		
\bib{Alt10}{article}{		
    AUTHOR = {Alt, Jesse},
     TITLE = {On quaternionic contact {F}efferman spaces},
   JOURNAL = {Differential Geom. Appl.},
    VOLUME = {28},
      YEAR = {2010},
    NUMBER = {4},
     PAGES = {376--394},
      ISSN = {0926-2245,1872-6984},
       DOI = {10.1016/j.difgeo.2010.04.001},
       URL = {https://doi.org/10.1016/j.difgeo.2010.04.001},
}



\bib{Arm}{arXiv}{
      title={Free $n$-distributions: holonomy, sub-Riemannian structures, Fefferman constructions and dual distributions}, 
      author={Stuart Armstrong},
      year={2007},
      eprint={0706.4441},
      url={https://arxiv.org/abs/0706.4441}, 
}

\bib{BDS}{article}{
   author={Burns, D., Jr.},
   author={Diederich, K.},
   author={Shnider, S.},
   title={Distinguished curves in pseudoconvex boundaries},
   journal={Duke Math. J.},
   volume={44},
   date={1977},
   number={2},
   pages={407--431},
   issn={0012-7094},
   review={\MR{0445009}},
}

\bib{Cap05}{article}{
   author={\v Cap, Andreas},
   title={Correspondence spaces and twistor spaces for parabolic geometries},
   journal={J. Reine Angew. Math.},
   volume={582},
   date={2005},
   pages={143--172},
   issn={0075-4102},
   review={\MR{2139714}},
   doi={10.1515/crll.2005.2005.582.143},
}



\bib{Cap06}{article}{
    AUTHOR = {\v Cap, Andreas},
     TITLE = {Two constructions with parabolic geometries},
   JOURNAL = {Rend. Circ. Mat. Palermo (2) Suppl.},
    NUMBER = {79},
      YEAR = {2006},
     PAGES = {11--37},
      ISSN = {1592-9531},
}

\bib{Cap08}{article}{
    AUTHOR = {\v Cap, Andreas},
     TITLE = {Infinitesimal automorphisms and deformations of parabolic
              geometries},
   JOURNAL = {J. Eur. Math. Soc. (JEMS)},
    VOLUME = {10},
      YEAR = {2008},
    NUMBER = {2},
     PAGES = {415--437},
      ISSN = {1435-9855,1435-9863},
       DOI = {10.4171/JEMS/116},
       URL = {https://doi.org/10.4171/JEMS/116},
}



\bib{CG08}{article}{
   author={\v Cap, Andreas},
   author={Gover, A. Rod},
   title={CR-tractors and the Fefferman space},
   journal={Indiana Univ. Math. J.},
   volume={57},
   date={2008},
   number={5},
   pages={2519--2570},
   issn={0022-2518},
   review={\MR{2463976}},
   doi={10.1512/iumj.2008.57.3359},
}

\bib{CG06}{article}{
author={\v Cap, Andreas},
   author={Gover, A. Rod},
     TITLE = {A holonomy characterisation of {F}efferman spaces},
   JOURNAL = {Ann. Global Anal. Geom.},
    VOLUME = {38},
      YEAR = {2010},
    NUMBER = {4},
     PAGES = {399--412},
      ISSN = {0232-704X,1572-9060},
       DOI = {10.1007/s10455-010-9220-6},
       URL = {https://doi.org/10.1007/s10455-010-9220-6},
}


\bib{CGH}{article}{
   author={\v Cap, Andreas},
   author={Gover, A. Rod},
   author={Hammerl, Matthias},
   title={Holonomy reductions of Cartan geometries and curved orbit
   decompositions},
   journal={Duke Math. J.},
   volume={163},
   date={2014},
   number={5},
   pages={1035--1070},
   issn={0012-7094},
   review={\MR{3189437}},
   doi={10.1215/00127094-2644793},
}

\bib{CM}{article}{
    author={\v Cap, Andreas},
    author={Mettler, Thomas},
    title = {Induced almost para-Kähler Einstein metrics on cotangent bundles},
    journal = {The Quarterly Journal of Mathematics},
    volume = {75},
    number = {4},
    pages = {1285-1299},
    year = {2024},
    issn = {0033-5606},
    doi = {10.1093/qmath/haae047},
    url = {https://doi.org/10.1093/qmath/haae047},
    eprint = {https://academic.oup.com/qjmath/article-pdf/75/4/1285/59123021/haae047.pdf},
}


\bib{Book}{book}{
author={\v Cap, Andreas},
author={Slov\'ak, Jan},
     TITLE = {Parabolic geometries. {I}},
    SERIES = {Mathematical Surveys and Monographs},
    VOLUME = {154},
      NOTE = {Background and general theory},
 PUBLISHER = {American Mathematical Society, Providence, RI},
      YEAR = {2009},
     PAGES = {x+628},
      ISBN = {978-0-8218-2681-2},
       DOI = {10.1090/surv/154},
       URL = {https://doi.org/10.1090/surv/154},
}

\bib{RelBGG}{article}{
   author={\v Cap, Andreas},
   author={Sou\v cek, Vladim\'ir},
   title={Relative BGG sequences; II. BGG machinery and invariant operators},
   journal={Adv. Math.},
   volume={320},
   date={2017},
   pages={1009--1062},
   issn={0001-8708},
   review={\MR{3709128}},
   doi={10.1016/j.aim.2017.09.016},
}

\bib{CZ09}{article}{
   author={\v Cap, Andreas},
   author={\v Z\'adn\'ik, Vojt\v ech},
   title={On the geometry of chains},
   journal={J. Differential Geom.},
   volume={82},
   date={2009},
   number={1},
   pages={1--33},
   issn={0022-040X},
   review={\MR{2504769}},
}

\bib{CrSa}{incollection}{
    author = {Crampin, Mike},
    author = {Saunders, David J.},
     TITLE = {Path geometries and almost {G}rassmann structures},
 BOOKTITLE = {Finsler geometry, {S}apporo 2005---in memory of {M}akoto
              {M}atsumoto},
    SERIES = {Adv. Stud. Pure Math.},
    VOLUME = {48},
     PAGES = {225--261},
 PUBLISHER = {Math. Soc. Japan, Tokyo},
      YEAR = {2007},
      ISBN = {978-4-931469-42-6},
       DOI = {10.2969/aspm/04810225},
       URL = {https://doi.org/10.2969/aspm/04810225},
}


\bib{DMT}{article}{
   author={Doubrov, Boris},
   author={Medvedev, Alexandr},
   author={The, Dennis},
   title={Homogeneous integrable Legendrian contact structures in dimension
   five},
   journal={J. Geom. Anal.},
   volume={30},
   date={2020},
   number={4},
   pages={3806--3858},
   issn={1050-6926},
   review={\MR{4167266}},
   doi={10.1007/s12220-019-00219-x},
}


\bib{Fef76}{article}{
   author={Fefferman, Charles L.},
   title={Monge-Amp\`ere equations, the Bergman kernel, and geometry of
   pseudoconvex domains},
   journal={Ann. of Math. (2)},
   volume={103},
   date={1976},
   number={2},
   pages={395--416},
   issn={0003-486X},
   review={\MR{0407320}},
   doi={10.2307/1970945},
}

\bib{Graham}{article}{
   author={Graham, C. Robin},
   title={On Sparling's characterization of Fefferman metrics},
   journal={Amer. J. Math.},
   volume={109},
   date={1987},
   number={5},
   pages={853--874},
   issn={0002-9327},
   review={\MR{0910354}},
   doi={10.2307/2374491},
}

\bib{Guo}{arXiv}{
      title={Parallel (co-)tractors and the geometry of first BGG solutions on almost Grassmannian structures}, 
      author={Zhangwen Guo},
      year={2025},
      eprint={2507.17605},
      url={https://arxiv.org/abs/2507.17605}, 
}



\bib{HS09}{article}{
   author={Hammerl, Matthias},
   author={Sagerschnig, Katja},
   title={Conformal structures associated to generic rank 2 distributions on
   5-manifolds---characterization and Killing-field decomposition},
   journal={SIGMA Symmetry Integrability Geom. Methods Appl.},
   volume={5},
   date={2009},
   pages={Paper 081, 29},
   review={\MR{2529166}},
   doi={10.3842/SIGMA.2009.081},
}

\bib{HS11}{article}{
   author={Hammerl, Matthias},
   author={Sagerschnig, Katja},
   title={The twistor spinors of generic 2- and 3-distributions},
   journal={Ann. Global Anal. Geom.},
   volume={39},
   date={2011},
   number={4},
   pages={403--425},
   issn={0232-704X},
   review={\MR{2776770}},
   doi={10.1007/s10455-010-9240-2},
}

\bib{Ham+17}{article}{
   author={Hammerl, Matthias},
   author={Sagerschnig, Katja},
   author={\v Silhan, Josef},
   author={Taghavi-Chabert, Arman},
   author={\v Z\'adn\'ik, Vojt\v ech},
   title={A projective-to-conformal Fefferman-type construction},
   journal={SIGMA Symmetry Integrability Geom. Methods Appl.},
   volume={13},
   date={2017},
   pages={Paper No. 081, 33},
   review={\MR{3714492}},
   doi={10.3842/SIGMA.2017.081},
}

\bib{HSSS12B}{article}{
author={Hammerl, Matthias},
author={Somberg, Petr},
author={Sou\v cek, Vladim\'ir},
author={\v Silhan, Josef},
     TITLE = {Invariant prolongation of overdetermined {PDE}s in projective,
              conformal, and {G}rassmannian geometry},
   JOURNAL = {Ann. Global Anal. Geom.},
    VOLUME = {42},
      YEAR = {2012},
    NUMBER = {1},
     PAGES = {121--145},
      ISSN = {0232-704X,1572-9060},
       DOI = {10.1007/s10455-011-9306-9},
       URL = {https://doi.org/10.1007/s10455-011-9306-9},
}




\bib{LNS17}{article}{
   author={Leistner, Thomas},
   author={Nurowski, Pawe\l},
   author={Sagerschnig, Katja},
   title={New relations between $\rm G_2$ geometries in dimensions 5 and 7},
   journal={Internat. J. Math.},
   volume={28},
   date={2017},
   number={13},
   pages={1750094, 46},
   issn={0129-167X},
   review={\MR{3737072}},
   doi={10.1142/S0129167X1750094X},
}


\bib{Ma+}{article}{
   author={Ma, Tianyu},
   author={Flood, Keegan J.},
   author={Matveev, Vladimir S.},
   author={\v Z\'adn\'ik, Vojt\v ech},
   title={Canonical curves and Kropina metrics in Lagrangian contact
   geometry},
   journal={Nonlinearity},
   volume={37},
   date={2024},
   number={1},
   pages={Paper No. 015007, 36},
   issn={0951-7715},
   review={\MR{4679002}},
   doi={10.1088/1361-6544/ad0c2b},
}








\bib{Met13}{article}{
    AUTHOR = {Mettler, Thomas},
     TITLE = {Reduction of {$\beta$}-integrable 2-{S}egre structures},
   JOURNAL = {Comm. Anal. Geom.},
    VOLUME = {21},
      YEAR = {2013},
    NUMBER = {2},
     PAGES = {331--353},
      ISSN = {1019-8385,1944-9992},
       DOI = {10.4310/CAG.2013.v21.n2.a3},
       URL = {https://doi.org/10.4310/CAG.2013.v21.n2.a3},
}



		
	\end{biblist}
\end{bibdiv}

\end{document}
